% ============================================================
%  SY02-18《shtuka 及其模空间导论》中文译注本·主文件
% ============================================================
\documentclass[UTF8, a4paper, 12pt, fontset=fandol, scheme=chinese]{ctexart}
% ============================================================
%  SY02-18《shtuka 及其模空间导论》中文译注本·导言区
%  底本：Zhiwei Yun, Introduction to Shtukas and their moduli
%        arXiv:2411.10248（IHES 2022 暑期学校讲义，38 页）
% ============================================================


% ---------- 页面 ----------
\usepackage[a4paper, left=2.6cm, right=2.6cm, top=2.8cm, bottom=2.6cm]{geometry}
\linespread{1.22}
\setlength{\parskip}{0.25em}
\setlength{\emergencystretch}{1.5em}
\sloppy

% ---------- 数学 ----------
\usepackage{amsmath, amsthm, amssymb, amsfonts}
\usepackage{mathtools}
\usepackage{mathrsfs}
\usepackage{stmaryrd}
\usepackage{bm}
\usepackage{amscd}

% ---------- 图 ----------
\usepackage{tikz}
\usetikzlibrary{cd, arrows.meta, positioning, calc}

% ---------- 其他 ----------
\usepackage[most]{tcolorbox}
\usepackage{enumitem}
\setlist[enumerate,1]{label=(\arabic*)}
\usepackage[hyperfootnotes=false, colorlinks=true, linkcolor=blue!45!black,
            citecolor=green!40!black, urlcolor=blue!55!black]{hyperref}
\usepackage{bookmark}

% ---------- 中英混排 ----------
\xeCJKsetup{CJKecglue={\hskip 0.24em plus 0.05em minus 0.05em}}
\ctexset{
  section    = {name={,}, aftername={.\enspace}, format=\Large\bfseries\centering},
  subsection = {aftername={.\enspace}, format=\large\bfseries},
  subsubsection = {aftername={.\enspace}, format=\normalsize\bfseries}
}

% ---------- 定理环境 ----------
% 原文 amsart：\numberwithin{equation}{section}，定理类挂 [subsubsection] 计数器，
% 故印刷为 Theorem 2.4.4（节.小节.序号）与公式 (2.1)（节.序号）。
\newtheoremstyle{zhthm}{3pt}{3pt}{\itshape}{}{\bfseries}{.}{.5em}{\thmnumber{#2}.\enspace\thmname{#1}\thmnote{\ \ (#3)}}
\newtheoremstyle{zhdefn}{3pt}{3pt}{}{}{\bfseries}{.}{.5em}{\thmnumber{#2}.\enspace\thmname{#1}\thmnote{\ \ (#3)}}
\newtheoremstyle{zhrmk}{3pt}{3pt}{}{}{\bfseries}{.}{.5em}{\thmnumber{#2}.\enspace\thmname{#1}\thmnote{\ \ (#3)}}
\theoremstyle{zhthm}
\newtheorem{thm}[subsubsection]{定理}
\newtheorem{prop}[subsubsection]{命题}
\newtheorem{lem}[subsubsection]{引理}
\newtheorem{cor}[subsubsection]{推论}
\newtheorem{conj}[subsubsection]{猜想}
\newtheorem{fact}[subsubsection]{事实}
\newtheorem{ansatz}[subsubsection]{Ansatz}
\theoremstyle{zhdefn}
\newtheorem{dfn}[subsubsection]{定义}
\newtheorem{cons}[subsubsection]{构造}
\newtheorem{exam}[subsubsection]{例}
\newtheorem{ques}[subsubsection]{问题}
\theoremstyle{zhrmk}
\newtheorem{rmk}[subsubsection]{注}
\numberwithin{equation}{section}

\renewcommand{\proofname}{证明}
\renewcommand{\refname}{参考文献}
\renewcommand{\contentsname}{目\quad 录}
\renewcommand{\abstractname}{摘\quad 要}

% ---------- 译注（独立脚注序列） ----------
\usepackage{manyfoot}
\DeclareNewFootnote[plain]{trans}
\renewcommand{\thefootnotetrans}{译\arabic{footnotetrans}}
\newcommand{\zy}[1]{\footnotetrans{#1}}

% ---------- 原文宏 ----------
\newcommand{\ssszh}[1]{\par\medskip\noindent\refstepcounter{subsubsection}\if\relax\detokenize{#1}\relax\else\textbf{#1}\par\fi}
% 由 工程文件/scripts/port_macros2.py 从原文导言区自动移植，勿手改。
\renewcommand{\AA}{\mathbb{A}}
\newcommand{\BB}{\mathbb{B}}
\newcommand{\CC}{\mathbb{C}}
\newcommand{\DD}{\mathbb{D}}
\newcommand{\EE}{\mathbb{E}}
\newcommand{\FF}{\mathbb{F}}
\newcommand{\GG}{\mathbb{G}}
\newcommand{\HH}{\mathbb{H}}
\newcommand{\II}{\mathbb{I}}
\newcommand{\JJ}{\mathbb{J}}
\newcommand{\KK}{\mathbb{K}}
\newcommand{\LL}{\mathbb{L}}
\newcommand{\MM}{\mathbb{M}}
\newcommand{\NN}{\mathbb{N}}
\newcommand{\OO}{\mathbb{O}}
\newcommand{\PP}{\mathbb{P}}
\newcommand{\QQ}{\mathbb{Q}}
\newcommand{\RR}{\mathbb{R}}
\renewcommand{\SS}{\mathbb{S}}
\newcommand{\TT}{\mathbb{T}}
\newcommand{\UU}{\mathbb{U}}
\newcommand{\VV}{\mathbb{V}}
\newcommand{\WW}{\mathbb{W}}
\newcommand{\XX}{\mathbb{X}}
\newcommand{\YY}{\mathbb{Y}}
\newcommand{\ZZ}{\mathbb{Z}}
\newcommand{\calA}{\mathcal{A}}
\newcommand{\calB}{\mathcal{B}}
\newcommand{\calC}{\mathcal{C}}
\newcommand{\calD}{\mathcal{D}}
\newcommand{\calE}{\mathcal{E}}
\newcommand{\calF}{\mathcal{F}}
\newcommand{\calG}{\mathcal{G}}
\newcommand{\calH}{\mathcal{H}}
\newcommand{\calI}{\mathcal{I}}
\newcommand{\calJ}{\mathcal{J}}
\newcommand{\calK}{\mathcal{K}}
\newcommand{\calL}{\mathcal{L}}
\newcommand{\calM}{\mathcal{M}}
\newcommand{\calN}{\mathcal{N}}
\newcommand{\calO}{\mathcal{O}}
\newcommand{\calP}{\mathcal{P}}
\newcommand{\calQ}{\mathcal{Q}}
\newcommand{\calR}{\mathcal{R}}
\newcommand{\calS}{\mathcal{S}}
\newcommand{\calT}{\mathcal{T}}
\newcommand{\calU}{\mathcal{U}}
\newcommand{\calV}{\mathcal{V}}
\newcommand{\calW}{\mathcal{W}}
\newcommand{\calX}{\mathcal{X}}
\newcommand{\calY}{\mathcal{Y}}
\newcommand{\calZ}{\mathcal{Z}}
\newcommand{\cA}{\mathcal{A}}
\newcommand{\cB}{\mathcal{B}}
\newcommand{\cC}{\mathcal{C}}
\newcommand{\cD}{\mathcal{D}}
\newcommand{\cE}{\mathcal{E}}
\newcommand{\cF}{\mathcal{F}}
\newcommand{\cG}{\mathcal{G}}
\newcommand{\cH}{\mathcal{H}}
\newcommand{\cI}{\mathcal{I}}
\newcommand{\cJ}{\mathcal{J}}
\newcommand{\cK}{\mathcal{K}}
\newcommand{\cL}{\mathcal{L}}
\newcommand{\cM}{\mathcal{M}}
\newcommand{\cN}{\mathcal{N}}
\newcommand{\cO}{\mathcal{O}}
\newcommand{\cP}{\mathcal{P}}
\newcommand{\cQ}{\mathcal{Q}}
\newcommand{\cR}{\mathcal{R}}
\newcommand{\cS}{\mathcal{S}}
\newcommand{\cT}{\mathcal{T}}
\newcommand{\cU}{\mathcal{U}}
\newcommand{\cV}{\mathcal{V}}
\newcommand{\cW}{\mathcal{W}}
\newcommand{\cX}{\mathcal{X}}
\newcommand{\cY}{\mathcal{Y}}
\newcommand{\cZ}{\mathcal{Z}}
\newcommand{\bA}{\mathbf{A}}
\newcommand{\bB}{\mathbf{B}}
\newcommand{\bC}{\mathbf{C}}
\newcommand{\bD}{\mathbf{D}}
\newcommand{\bE}{\mathbf{E}}
\newcommand{\bF}{\mathbf{F}}
\newcommand{\bG}{\mathbf{G}}
\newcommand{\bH}{\mathbf{H}}
\newcommand{\bI}{\mathbf{I}}
\newcommand{\bJ}{\mathbf{J}}
\newcommand{\bK}{\mathbf{K}}
\newcommand{\bL}{\mathbf{L}}
\newcommand{\bM}{\mathbf{M}}
\newcommand{\bN}{\mathbf{N}}
\newcommand{\bO}{\mathbf{O}}
\newcommand{\bP}{\mathbf{P}}
\newcommand{\bQ}{\mathbf{Q}}
\newcommand{\bR}{\mathbf{R}}
\newcommand{\bS}{\mathbf{S}}
\newcommand{\bT}{\mathbf{T}}
\newcommand{\bU}{\mathbf{U}}
\newcommand{\bV}{\mathbf{V}}
\newcommand{\bW}{\mathbf{W}}
\newcommand{\bX}{\mathbf{X}}
\newcommand{\bY}{\mathbf{Y}}
\newcommand{\bZ}{\mathbf{Z}}
\newcommand{\frA}{\mathfrak{A}}
\newcommand{\frB}{\mathfrak{B}}
\newcommand{\frC}{\mathfrak{C}}
\newcommand{\frD}{\mathfrak{D}}
\newcommand{\frE}{\mathfrak{E}}
\newcommand{\frF}{\mathfrak{F}}
\newcommand{\frG}{\mathfrak{G}}
\newcommand{\frH}{\mathfrak{H}}
\newcommand{\frI}{\mathfrak{I}}
\newcommand{\frJ}{\mathfrak{J}}
\newcommand{\frK}{\mathfrak{K}}
\newcommand{\frL}{\mathfrak{L}}
\newcommand{\frM}{\mathfrak{M}}
\newcommand{\frN}{\mathfrak{N}}
\newcommand{\frO}{\mathfrak{O}}
\newcommand{\frP}{\mathfrak{P}}
\newcommand{\frQ}{\mathfrak{Q}}
\newcommand{\frR}{\mathfrak{R}}
\newcommand{\frS}{\mathfrak{S}}
\newcommand{\frT}{\mathfrak{T}}
\newcommand{\frU}{\mathfrak{U}}
\newcommand{\frV}{\mathfrak{V}}
\newcommand{\frW}{\mathfrak{W}}
\newcommand{\frX}{\mathfrak{X}}
\newcommand{\frY}{\mathfrak{Y}}
\newcommand{\frZ}{\mathfrak{Z}}
\newcommand{\fra}{\mathfrak{a}}
\newcommand{\frb}{\mathfrak{b}}
\newcommand{\frc}{\mathfrak{c}}
\newcommand{\frd}{\mathfrak{d}}
\newcommand{\fre}{\mathfrak{e}}
\newcommand{\frf}{\mathfrak{f}}
\newcommand{\frg}{\mathfrak{g}}
\newcommand{\frh}{\mathfrak{h}}
\newcommand{\fri}{\mathfrak{i}}
\newcommand{\frj}{\mathfrak{j}}
\newcommand{\frk}{\mathfrak{k}}
\newcommand{\frl}{\mathfrak{l}}
\newcommand{\fm}{\mathfrak{m}}
\newcommand{\frn}{\mathfrak{n}}
\newcommand{\frp}{\mathfrak{p}}
\newcommand{\frq}{\mathfrak{q}}
\newcommand{\frr}{\mathfrak{r}}
\newcommand{\frs}{\mathfrak{s}}
\newcommand{\frt}{\mathfrak{t}}
\newcommand{\fru}{\mathfrak{u}}
\newcommand{\frv}{\mathfrak{v}}
\newcommand{\frw}{\mathfrak{w}}
\newcommand{\frx}{\mathfrak{x}}
\newcommand{\fry}{\mathfrak{y}}
\newcommand{\frz}{\mathfrak{z}}
\newcommand{\tilA}{\widetilde{A}}
\newcommand{\tilB}{\widetilde{B}}
\newcommand{\tilC}{\widetilde{C}}
\newcommand{\tilD}{\widetilde{D}}
\newcommand{\tilE}{\widetilde{E}}
\newcommand{\tilF}{\widetilde{F}}
\newcommand{\tilG}{\widetilde{G}}
\newcommand{\tilH}{\widetilde{H}}
\newcommand{\tilI}{\widetilde{I}}
\newcommand{\tilJ}{\widetilde{J}}
\newcommand{\tilK}{\widetilde{K}}
\newcommand{\tilL}{\widetilde{L}}
\newcommand{\tilM}{\widetilde{M}}
\newcommand{\tilN}{\widetilde{N}}
\newcommand{\tilO}{\widetilde{O}}
\newcommand{\tilP}{\widetilde{P}}
\newcommand{\tilQ}{\widetilde{Q}}
\newcommand{\tilR}{\widetilde{R}}
\newcommand{\tilS}{\widetilde{S}}
\newcommand{\tilT}{\widetilde{T}}
\newcommand{\tilU}{\widetilde{U}}
\newcommand{\tilV}{\widetilde{V}}
\newcommand{\tilW}{\widetilde{W}}
\newcommand{\tilX}{\widetilde{X}}
\newcommand{\tilY}{\widetilde{Y}}
\newcommand{\tilZ}{\widetilde{Z}}
\newcommand{\tila}{\widetilde{a}}
\newcommand{\tilb}{\widetilde{b}}
\newcommand{\tilc}{\widetilde{c}}
\newcommand{\tild}{\widetilde{d}}
\newcommand{\tile}{\widetilde{e}}
\newcommand{\tilf}{\widetilde{f}}
\newcommand{\tilg}{\widetilde{g}}
\newcommand{\tilh}{\widetilde{h}}
\newcommand{\tili}{\widetilde{i}}
\newcommand{\tilj}{\widetilde{j}}
\newcommand{\tilk}{\widetilde{k}}
\newcommand{\tilm}{\widetilde{m}}
\newcommand{\tiln}{\widetilde{n}}
\newcommand{\tilo}{\widetilde{o}}
\newcommand{\tilp}{\widetilde{p}}
\newcommand{\tilq}{\widetilde{q}}
\newcommand{\tilr}{\widetilde{r}}
\newcommand{\tils}{\widetilde{s}}
\newcommand{\tilt}{\widetilde{t}}
\newcommand{\tilu}{\widetilde{u}}
\newcommand{\tilv}{\widetilde{v}}
\newcommand{\tilw}{\widetilde{w}}
\newcommand{\tilx}{\widetilde{x}}
\newcommand{\tily}{\widetilde{y}}
\newcommand{\tilz}{\widetilde{z}}
\newcommand{\dG}{\widehat{G}}
\newcommand{\dB}{\widehat{B}}
\newcommand{\dH}{\widehat{H}}
\newcommand{\dT}{\widehat{T}}
\newcommand{\aff}{\textup{aff}}
\newcommand{\AJ}{\textup{AJ}}
\newcommand{\alg}{\textup{alg}}
\newcommand{\Alg}{\textup{Alg}}
\newcommand{\AS}{\textup{AS}}
\newcommand{\Av}{\textup{Av}}
\newcommand{\av}{\textup{av}}
\newcommand{\BM}{\textup{BM}}
\newcommand{\Bun}{\textup{Bun}}
\newcommand{\can}{\textup{can}}
\newcommand{\ch}{\textup{char}}
\newcommand{\Ch}{\textup{Ch}}
\newcommand{\codim}{\textup{codim}}
\newcommand{\Coh}{\textup{Coh}}
\newcommand{\coker}{\textup{coker}}
\newcommand{\Coor}{\textup{Coor}}
\newcommand{\Corr}{\textup{Corr}}
\newcommand{\Div}{\textup{Div}}
\newcommand{\Dyn}{\textup{Dyn}}
\newcommand{\ev}{\textup{ev}}
\newcommand{\even}{\textup{even}}
\newcommand{\Fil}{\textup{Fil}}
\newcommand{\Fl}{\textup{Fl}}
\newcommand{\fl}{f\ell}
\newcommand{\Forg}{\textup{Forg}}
\newcommand{\Fr}{\textup{Fr}}
\newcommand{\Frac}{\textup{Frac}}
\newcommand{\Frob}{\textup{Frob}}
\newcommand{\Gal}{\textup{Gal}}
\newcommand{\geom}{\textup{geom}}
\newcommand{\glob}{\textup{glob}}
\newcommand{\Gr}{\textup{Gr}}
\newcommand{\gr}{\textup{gr}}
\newcommand{\Herm}{\textup{Herm}}
\newcommand{\Hk}{\textup{Hk}}
\newcommand{\Hilb}{\textup{Hilb}}
\newcommand{\IC}{\textup{IC}}
\newcommand{\id}{\textup{id}}
\renewcommand{\Im}{\textup{Im}}
\newcommand{\Ind}{\textup{Ind}}
\newcommand{\ind}{\textup{ind}}
\newcommand{\inv}{\textup{inv}}
\newcommand{\Irr}{\textup{Irr}}
\newcommand{\Jac}{\textup{Jac}}
\newcommand{\Lie}{\textup{Lie}\ }
\newcommand{\Loc}{\textup{Loc}}
\newcommand{\loc}{\textup{loc}}
\newcommand{\Mod}{\textup{Mod}}
\newcommand{\Mot}{\textup{Mot}}
\newcommand{\mult}{\textup{mult}}
\newcommand{\nil}{\textup{nil}}
\newcommand{\Nm}{\textup{Nm}}
\newcommand{\odd}{\textup{odd}}
\newcommand{\opp}{\textup{opp}}
\newcommand{\Out}{\textup{Out}}
\newcommand{\Perf}{\textup{Perf}}
\newcommand{\Perv}{\textup{Perv}}
\newcommand{\Pic}{\textup{Pic}}
\newcommand{\pr}{\textup{pr}}
\newcommand{\pro}{\textup{pro}}
\newcommand{\Proj}{\textup{Proj}}
\newcommand{\pt}{\textup{pt}}
\newcommand{\Quot}{\textup{Quot}}
\newcommand{\rank}{\textup{rank}\ }
\newcommand{\red}{\textup{red}}
\newcommand{\reg}{\textup{reg}}
\newcommand{\Rep}{\textup{Rep}}
\newcommand{\Res}{\textup{Res}}
\newcommand{\res}{\textup{res}}
\newcommand{\rk}{\textup{rk}}
\newcommand{\rs}{\textup{rs}}
\newcommand{\Sht}{\textup{Sht}}
\newcommand{\Span}{\textup{Span}}
\newcommand{\Spec}{\textup{Spec}\ }
\newcommand{\St}{\textup{St}}
\newcommand{\st}{\textup{st}}
\newcommand{\Stab}{\textup{Stab}}
\newcommand{\Supp}{\textup{Supp}}
\newcommand{\Sym}{\textup{Sym}}
\newcommand{\tors}{\textup{tors}}
\newcommand{\Tot}{\textup{Tot}}
\newcommand{\Tr}{\textup{Tr}}
\newcommand{\tr}{\textup{tr}}
\newcommand{\triv}{\textup{triv}}
\newcommand{\unip}{\textup{unip}}
\newcommand{\univ}{\textup{univ}}
\newcommand{\ur}{\textup{ur}}
\newcommand{\val}{\textup{val}}
\newcommand{\Vect}{\textup{Vect}}
\newcommand{\vol}{\textup{vol}}
\newcommand{\Aut}{\textup{Aut}}
\newcommand{\Hom}{\textup{Hom}}
\newcommand{\End}{\textup{End}}
\newcommand{\Map}{\textup{Map}}
\newcommand{\Mor}{\textup{Mor}}
\newcommand{\uAut}{\underline{\Aut}}
\newcommand{\uHom}{\underline{\Hom}}
\newcommand{\uEnd}{\underline{\End}}
\newcommand{\uMor}{\underline{\Mor}}
\newcommand{\RHom}{\bR\Hom}
\newcommand{\RuHom}{\bR\uHom}
\newcommand{\Ext}{\textup{Ext}}
\newcommand{\uExt}{\underline{\Ext}}
\newcommand{\GL}{\textup{GL}}
\newcommand{\gl}{\mathfrak{gl}}
\newcommand{\PGL}{\textup{PGL}}
\newcommand{\Ug}{\textup{U}}
\newcommand{\ug}{\mathfrak{u}}
\newcommand{\SL}{\textup{SL}}
\renewcommand{\sl}{\mathfrak{sl}}
\newcommand{\SU}{\textup{SU}}
\newcommand{\su}{\mathfrak{su}}
\newcommand{\GU}{\textup{GU}}
\newcommand{\SO}{\textup{SO}}
\newcommand{\GO}{\textup{GO}}
\newcommand{\Og}{\textup{O}}
\newcommand{\so}{\mathfrak{so}}
\newcommand{\Sp}{\textup{Sp}}
\renewcommand{\sp}{\mathfrak{sp}}
\newcommand{\GSp}{\textup{GSp}}
\newcommand{\Gm}{\GG_m}
\newcommand{\Ga}{\GG_a}
\newcommand{\ad}{\textup{ad}}
\newcommand{\Ad}{\textup{Ad}}
\renewcommand{\sc}{\textup{sc}}
\newcommand{\der}{\textup{der}}
\newcommand{\xch}{\mathbb{X}^*}
\newcommand{\xcoch}{\mathbb{X}_*}
\newcommand{\incl}{\hookrightarrow}
\newcommand{\isom}{\stackrel{\sim}{\to}}
\newcommand{\bij}{\leftrightarrow}
\newcommand{\surj}{\twoheadrightarrow}
\newcommand{\leftexp}[2]{{\vphantom{#2}}^{#1}{#2}}
\newcommand{\pH}{\leftexp{p}{\textup{H}}}
\newcommand{\trpose}[1]{\leftexp{t}{#1}}
\newcommand{\Ql}{\QQ_{\ell}}
\newcommand{\Qlbar}{\overline{\QQ}_\ell}
\newcommand{\const}[1]{\overline{\QQ}_{\ell,#1}}
\newcommand{\twtimes}[1]{\stackrel{#1}{\times}}
\newcommand{\hotimes}{\widehat{\otimes}}
\newcommand{\htimes}{\widehat{\times}}
\renewcommand{\j}[1]{\langle{#1}\rangle}
\newcommand{\wt}[1]{\widetilde{#1}}
\newcommand{\wh}[1]{\widehat{#1}}
\newcommand{\mat}[4]{\left(\begin{array}{cc} #1 & #2 \\ #3 & #4 \end{array}\right)}
\newcommand{\un}{\underline}
\newcommand{\bu}{\bullet}
\newcommand{\ov}{\overline}
\newcommand{\bs}{\backslash}
\newcommand{\tl}[1]{[\![#1]\!]}
\newcommand{\lr}[1]{(\!(#1)\!)}
\newcommand{\sss}{\subsubsection}
\newcommand{\xr}{\xrightarrow}
\newcommand{\op}{\oplus}
\newcommand{\ot}{\otimes}
\newcommand{\one}{\mathbf{1}}
\renewcommand{\sslash}{\mathbin{/\mkern-6mu/}}
\newcommand{\vn}{\varnothing}
\newcommand{\homo}[2]{\mathbf{H}_{#1}({#2})}
\newcommand{\homog}[2]{\textup{H}_{#1}({#2})}
\newcommand{\coho}[2]{\mathbf{H}^{#1}({#2})}
\newcommand{\cohog}[2]{\textup{H}^{#1}({#2})}
\newcommand{\cohoc}[2]{\textup{H}_{c}^{#1}({#2})}
\newcommand{\hBM}[2]{\textup{H}^{\textup{BM}}_{#1}({#2})}
\newcommand{\upH}{\textup{H}}
\newcommand{\oll}[1]{\overleftarrow{#1}}
\newcommand{\orr}[1]{\overrightarrow{#1}}
\newcommand{\olr}[1]{\overleftrightarrow{#1}}
\newcommand{\e}{\epsilon}
\newcommand{\io}{\iota}
\newcommand{\ka}{\kappa}
\renewcommand{\th}{\theta}
\newcommand{\Th}{\Theta}
\newcommand{\ph}{\varphi}
\newcommand{\s}{\sigma}
\newcommand{\Sig}{\Sigma}
\renewcommand{\t}{\tau}
\newcommand{\x}{\chi}
\newcommand{\y}{\eta}
\newcommand{\z}{\zeta}
\newcommand{\ep}{\epsilon}
\newcommand{\vp}{\varpi}
\renewcommand{\l}{\lambda}
\renewcommand{\L}{\Lambda}
\newcommand{\om}{\omega}
\newcommand{\Om}{\Omega}
\newcommand{\dm}{\diamondsuit}
\newcommand{\hs}{\heartsuit}
\newcommand{\na}{\natural}
\newcommand{\sh}{\sharp}
\newcommand{\da}{\dagger}
\newcommand{\kbar}{\overline{k}}
\newcommand{\Qbar}{\overline{\QQ}}
\newcommand{\Gk}{\Gal(\kbar/k)}
\newcommand{\GQ}{\Gal(\overline{\QQ}/\QQ)}
\newcommand{\GQp}{\Gal(\overline{\QQ}_{p}/\QQ_{p})}
\newcommand{\Aff}{\mathbf{Aff}}
\newcommand{\Sch}{\mathbf{Sch}}
\newcommand{\Gpd}{\mathbf{Gpd}}
\newcommand{\dr}{\dashrightarrow}
\newcommand{\dxr}{\dashrightarrow}
\newcommand{\br}{\breve}
\newcommand{\pl}{\partial}
\newcommand{\QGr}{\mathrm{QGr}}
\newcommand{\disj}{\mathrm{disj}}
\newcommand{\pos}{\mathrm{pos}}
\newcommand{\Sat}{\mathrm{Sat}}
\newcommand{\HN}{\mathrm{HN}}
\newcommand{\ab}{\mathrm{ab}}
\newcommand{\colim}{\mathrm{colim}}
\newcommand{\Fun}{\mathrm{Fun}}
\newcommand{\IsoSht}{\mathrm{IsoSht}}
\newcommand{\IsoCrys}{\mathrm{IsoCrys}}
\newcommand{\sep}{\mathrm{sep}}
\renewcommand{\div}{\mathrm{div}}
\newcommand{\ord}{\mathrm{ord}}
\newcommand{\Dr}{\mathrm{Dr}}
\newcommand{\Rees}{\mathrm{Rees}}
\newcommand{\sL}{\mathscr{L}}
\newcommand{\HD}{\textup{HD}}
\newcommand{\GGP}{\textup{GGP}}
\newcommand{\KR}{\textup{KR}}
\newcommand{\on}{\operatorname}


% 原文用 \renewcommand\a\alpha 这类「单个 token 作宏体」的写法，
% 移植脚本抓不到；这里补齐（逐个核对 460–491 行）。
\renewcommand{\a}{\alpha}
\renewcommand{\b}{\beta}
\newcommand{\g}{\gamma}
\newcommand{\G}{\Gamma}
\renewcommand{\d}{\delta}
\newcommand{\D}{\Delta}
\renewcommand{\r}{\rho}
\renewcommand{\c}{\circ}

\newcommand{\Fp}{\mathbb{F}_p}

% 原文用 xy-pic 画图。必须在 tikz 之后加载，否则 \xymatrix 与 tikz-cd 互相冲突。
\usepackage[all]{xy}


\begin{document}
\pagestyle{plain}

\begin{titlepage}
\centering
\vspace*{3cm}
{\Huge shtuka 及其模空间导论\par}
\vspace{1.2cm}
{\LARGE Introduction to Shtukas and their moduli\par}
\vspace{1.2cm}
{\large 中文译注本\par}
\vspace{2cm}
{\large 恽之玮（Zhiwei Yun）\par}
\vspace{0.6cm}
{\large IHES 2022 暑期学校（Langlands 纲领专题）讲义\par}
\vspace{1.6cm}
{\large 原文：arXiv:2411.10248，38 页，2024 年\par}
\vspace{0.4cm}
{\large 译者：数译（DSH agent）\par}
\vspace{0.4cm}
{\large 版本：v1.0　　日期：2026 年 11 月\par}
\vspace{1.6cm}
\rule{0.6\textwidth}{0.4pt}\par
\vspace{0.6cm}
{\small 本译本仅供学术研读；原文版权归作者所有。\par}
\vfill
\end{titlepage}

\section*{译者说明}
\addcontentsline{toc}{section}{译者说明}

\textbf{原文性质。}这是恽之玮（Zhiwei Yun）在 2022 年 IHES 暑期学校（Langlands 纲领专题）上关于「shtuka」的讲稿，属\textbf{综述性质}：定义与主要定理给全，证明多半只给出思路或引用文献。原文共五节：引言、基本定义、shtuka 的「亲戚」（Drinfeld 模、(仿射) Deligne–Lusztig 簇、有限域上的动机）、模叠的几何、上同调上的结构、进阶课题。

\textbf{编号约定。}原文是 amsart 体例：公式按节编号，形如 (2.1)；定理类环境与 \verb|subsubsection| 共用计数器，印刷为「Theorem 2.4.4」（节.小节.序号）。译文完全照此：定理类编号形如「定理 2.4.4」，公式形如 (2.1)，与原文逐条对应。

\textbf{译注标记。}文中「译$n$」是译者所加，用独立脚注序列，不占用作者脚注；作者脚注保留原编号。译注只写三类内容：易错预警、含糊之处的点明、背景直觉。

\textbf{术语约定。}核心术语在全书首次出现时保留英文原词，其后只用中文：shtuka 保留原文不译（俄语借词，原意为「小玩意」）；$G$-torsor 译「$G$-主丛」；Hecke stack 译「Hecke 叠」；excursion operator 译「游走算子」；partial Frobenius 译「部分 Frobenius」；奇点译法见 \verb|资料卡/04-术语文献勘误.md|。

\newpage
\tableofcontents
\newpage

% ============ 正文 ============

% ch01  §1 引言 / §2.1–§2.3（原文 585–962 行）

\section{引言}

shtuka 这一概念由 Drinfeld \cite{Dr-FBun} 以「$F$-丛」之名引入，其中 $F$ 指 Frobenius。事实上，shtuka 就是带某种 Frobenius 结构的丛。「shtuka」一词源自俄语，意为「小玩意」或「东西」。shtuka 推广了 Drinfeld 更早引入的「椭圆模」概念，后者是椭圆曲线的类比物。一般的 shtuka 可以看成动机的一种版本：它不定义在特征 $p$ 之中，但其系数域的特征为 $p$。

shtuka 的模空间或模叠是 Shimura 簇在函数域上的类比物。因此，它们在建立函数域上 Langlands 对应的各种情形时一直起着关键作用。

Drinfeld \cite{Dr-Pet, Dr-comp} 用秩 $2$ shtuka 的模空间证明了整体函数域上 $\GL_{2}$ 的 Ramanujan–Petersson 猜想与 Langlands 猜想。

Laumon、Rapoport 与 Shtuhler \cite{LRS} 用除代数上的 shtuka 证明了函数域上除代数的整体 Langlands 猜想的一部分，并证明了局部函数域上 $\GL_{n}$（任意 $n$）的局部 Langlands 猜想。

L. Lafforgue \cite{LLaff} 用秩 $n$ shtuka 的模叠及其紧化，对任意 $n\in\NN$ 证明了整体函数域上 $\GL_{n}$ 的 Ramanujan–Petersson 猜想与 Langlands 猜想。

Varshavsky \cite{Var} 把 shtuka 的概念从向量丛的情形推广到主 $G$-丛的情形，其中 $G$ 是基域 $\FF_{q}$ 上任一分裂约化群。我们称这样的对象为 $G$-shtuka。他证明了关于 $G$-shtuka 模叠的几何的基本结果。

更晚近地，V. Lafforgue \cite{VLaff} 利用 Varshavsky 的 $G$-shtuka 模空间，结合若干巧妙的想法（其中最重要的是游走算子），证明了函数域上任一连通约化群 $G$ 的整体 Langlands 猜想的「自守到 Galois」方向。

局部版的 shtuka 虽不在本文讨论范围之内，却在证明局部 Langlands 猜想时起重要作用。在 \cite{GL} 中，Genestier 与 V. Lafforgue 用局部 shtuka 证明了局部函数域上任一连通约化群 $G$ 的局部 Langlands 猜想的「表示到 Galois」方向。另一项令人兴奋的进展是 Scholze 关于混合特征局部域的局部 shtuka 理论，其模空间推广了 Rapoport–Zink 空间，参见 \cite{S-Berkeley}。

本文是一篇综述。我们将在 \S\ref{s:def} 中从头介绍 $\FF_{q}$ 上分裂约化群的 shtuka 概念；接着在 \S\ref{s:sib} 中描述它与 Drinfeld 模的关系，以及它与 Deligne–Lusztig 簇、有限域上动机的类比；然后在 \S\ref{s:geom} 与 \S\ref{s:coho} 中研究 shtuka 模叠上的一些重要几何构造，例如 Hecke 对应与部分 Frobenius，并讨论它们对模叠上同调的影响。这些内容为 C. Xue \cite{Xue} 的讲座提供背景。最后一节 \S\ref{s:adv} 提及若干值得进一步研究的有趣课题，例如模叠的紧化，以及把 shtuka 模叠上的特殊圈与 $L$-函数、Eisenstein 级数的高阶导数联系起来的新近工作。

关于 shtuka 已有出色的综述文章，例如 \cite{Laumon} 与 \cite{NDT-Sht}，作者正是从中习得这套理论。研究论文 \cite{Var} 与 \cite{ARH} 可作为证明细节的详尽参考。最后，本文很大一部分遵循突破性论文 \cite{VLaff} 的处理方式，读者可从中找到更多细节及其在 Langlands 纲领中的应用。

%%% thanks %%%

\subsection*{致谢}

感谢 2022 年 IHES 暑期学校（Langlands 纲领专题）的组织者邀请我就此主题作报告。感谢 Jianqiao Xia 与审稿人仔细阅读草稿并提出有益的建议。



\section{基本定义}\label{s:def}

\subsection{曲线上的 $G$-主丛}
设 $k$ 为有限域 $\FF_{q}$。对两个 $k$-叠 $\cX$ 与 $\cY$，用 $\cX\times\cY$ 表示它们在 $\Spec k$ 上的纤维积。

设 $X$ 是 $k$ 上几何连通的光滑射影曲线，$F=k(X)$ 为其函数域，$|X|$ 为 $X$ 的闭点集。对 $x\in |X|$，记 $\cO_{x}$、$\fm_{x}$、$F_{x}$、$k_{x}$ 分别为 $X$ 在 $x$ 处的完备化局部环及其极大理想、分式域与剩余域。

\ssszh{$G$-主丛}
设 $G$ 是 $k$ 上分裂连通约化群。$k$-概形 $Y$ 上的一个 $G$-主丛，是指一个概形 $\cE$，带 $G$ 的右作用以及 $G$-不变的映射 $\pi: \cE\to Y$，使得在 $Y$ 上 étale 局部地，$\cE$ 与 $G\times Y$（带 $G$ 的右平移作用）$G$-等变同构。

令 $\Bun_{G}: \Sch_{k}\to \Gpd$（其中 $\Sch_{k}$ 是 $k$-概形范畴，$\Gpd$ 是群胚范畴）为如下函子：把 $k$-概形 $S$ 对应到 $X\times S:=X\times_{k}S$ 上 $G$-主丛的群胚，其态射为 $G$-主丛的同构。则 $\Bun_{G}$ 可由 $k$ 上局部有限型的光滑 Artin 叠表示。

\begin{exam} 设 $G=\GL_{n}$。概形 $Y$ 上 $\GL_{n}$-主丛的群胚与 $Y$ 上秩 $n$ 向量丛的群胚等价。因此 $\Bun_{G}$ 同构于 $X$ 上秩 $n$ 向量丛的模叠 $\Bun_{n}$。

若 $G=\SL_{n}$，则 $\Bun_{G}(S)$ 参数化偶 $(\cV, \l)$，其中 $\cV$ 是 $X\times S$ 上秩 $n$ 的向量丛，$\l$ 是同构 $\l: \cO_{X\times S}\isom \det(\cV):=\wedge^{n}(\cV)$。
\end{exam}



\begin{exam} 设 $G=\Sp_{2n}$。概形 $Y$ 上 $\Sp_{2n}$-主丛的群胚与 $Y$ 上辛向量丛 $(\cV, \om)$ 的群胚等价。这里 $\cV$ 是 $Y$ 上秩 $2n$ 的向量丛，$\om: \cV\ot \cV\to \cO_{Y}$ 是完美交错配对。因此 $\Bun_{G}$ 同构于 $X$ 上秩 $2n$ 辛向量丛的模叠。
\end{exam}



\ssszh{水平结构}\label{sss:level}

设 $D\subset X$ 是有限子概形，$S$ 是 $k$-概形。$X\times S$ 上 $G$-主丛 $\cE$ 的一个 $D$-水平结构，是指限制所得 $D\times S$ 上 $G$-主丛 $\cE|_{D\times S}$ 的一个平凡化。令 $\Bun_{G,D}: \Sch_{k}\to \Gpd$ 为如下函子：它在 $S$ 处的取值是偶 $(\cE,\io_{D})$ 的群胚，其中 $\cE$ 是 $X\times S$ 上的 $G$-主丛，$\io_{D}$ 是 $\cE$ 的 $D$-水平结构。则 $\Bun_{G,D}$ 可由 Artin 叠表示，且遗忘映射 $\Forg_{D}: \Bun_{G,D}\to \Bun_{G}$ 是一个 $G_{D}:=\Res_{D/k}(G\times D)$-主丛。


\begin{exam} 设 $G=\Sp_{2n}$，并对某个 $x\in |X|$ 取 $D=x$。对任意 $k$-概形 $S$，$\Bun_{G,x}(S)$ 是多元组 $(\cV, \om, \{e_{1},\cdots, e_{n}, f_{1},\cdots, f_{n}\})$ 的群胚，其中 $(\cV,\om)$ 是 $X\times S$ 上秩 $2n$ 的辛向量丛，$e_{i}$ 与 $f_{i}$ 是 $\cV|_{D\times S}$ 的截面，且对所有 $1\le i,j\le n$ 有 $\om(e_{i},e_{j})=0$、$\om(f_{i},f_{j})=0$、$\om(e_{i},f_{j})=\d_{ij}$（这些都是 $\cO_{D\times S}$ 的截面之间的等式）。
\end{exam}


考虑其它的水平结构也有用处。设 $x\in |X|$。记 $L^{+}_{x}G$ 为 $G$ 在 $x$ 处的形式弧群：它是 $k_{x}$ 上的群概形，其 $R$-点（$R$ 为 $k_{x}$-代数）为 $G(R\hotimes_{k_{x}} \cO_{x})$。记 $L_{x}G$ 为 $G$ 在 $x$ 处的形式环路群：它是 $k_{x}$ 上的群 ind-概形，其 $R$-点为 $G(R\hotimes_{k_{x}} F_{x})$。参见 \cite[\S1.a]{PR}。

对 $n\in \NN$，令 $D_{nx}=\Spec \cO_{x}/\fm^{n}_{x}$。我们有 $L^{+}_{x}G$ 的截断版本 $G_{nx}:=\Res_{D_{nx}/k_{x}}(G\times D_{nx})$。令 $r_{nx}: L^{+}_{x}G\to G_{nx}$ 为约化映射。称
\begin{equation*}
L^{(n)}_{x}G=\ker(r_{nx})
\end{equation*}
为 $L^{+}_{x}G$ 的第 $n$ 个主同余子群。

设 $\bK_{x}\subset L^{+}_{x}G$ 是 $k_{x}$ 上的子群概形，且具有形状 $r_{nx}^{-1}(\ov\bK_{x,n})$，其中 $n\in\NN$，$\ov\bK_{x,n}\subset G_{nx}$ 是 $k_{x}$-子群。定义 $\Bun_{G,\bK_{x}}$ 为商叠
\begin{equation*}
\Bun_{G,\bK_{x}}:=\Bun_{G, nx}/(\Res_{k_{x}/k}\ov\bK_{x,n}).
\end{equation*}
注意 $\Res_{k_{x}/k}G_{nx}=\Res_{D_{nx}/k}(G\times D_{nx})$ 作用在 $\Bun_{G, nx}$ 上，因此其子群 $\Res_{k_{x}/k}\ov\bK_{x,n}$ 也作用。容易看出，$\Bun_{G,\bK_{x}}$ 在相差典范同构的意义下与 $n$ 的选取无关。



\begin{exam} 设 $\bI_{x}$ 是 $L^{+}_{x}G$ 的 Iwahori 子群，即它是某个 Borel 子群 $B_{x}\subset G\ot k_{x}$ 在 $r_{x}: L^{+}_{x}G\to G\ot k_{x}$ 下的原像。则 $\Bun_{G,\bI_{x}}$ 的 $S$-点与偶 $(\cE, \cF_{x})$ 是同一份数据，其中 $\cE$ 是 $X\times S$ 上的 $G$-主丛，$\cF_{x}$ 是限制 $\cE|_{\{x\}\times S}$ 的一个 Borel 约化，也就是丛 $\cE|_{\{x\}\times S}\times^{G}\fl_{G}\to \{x\}\times S$ 的一个截面（$\fl_{G}$ 是 $G$ 的旗簇）。由后一描述可见，$\Bun_{G,\bI_{x}}$ 典范地不依赖于 Borel 子群 $B_{x}$ 的选取。
\end{exam}

更一般地，当 $\bK_{x}\subset L_{x}G$ 不必包含在 $L^{+}_{x}G$ 中，而是以 $L^{+}_{x}G$ 的某个同余子群为余维有限的正规子群时，仍可定义模叠 $\Bun_{G,\bK_{x}}$。这样的 $\bK_{x}$ 称为 $x$ 处的 {\em 水平群}（level group）。为构造 $\Bun_{G,\bK_{x}}$，令 $\Bun_{G, \infty x}$ 为塔 $(\Bun_{G,nx})_{n\in\NN}$ 的逆向极限。可以证明，$\Bun_{G,\infty x}$ 上的 $L^{+}_{x}G$-作用典范地延拓为 $L_{x}G$ 的作用。于是定义 $\Bun_{G,\bK_{x}}$ 为商 $\Bun_{G,\infty x}/(\Res_{k_{x}/k}\bK_{x})$。为看出它是 Artin 叠，取一个在 $\bK_{x}$ 中正规的主同余子群 $L^{(n)}_{x}G$，令 $\ov \bK_{x,n}=\bK_{x}/L^{(n)}_{x}G$，并把 $\Bun_{G,\infty x}/(\Res_{k_{x}/k}\bK_{x})$ 等同于 $\Bun_{G,nx}/(\Res_{k_{x}/k}\ov\bK_{x,n})$。特别地，上述构造可施于抛孔子群 $\bP_{x}\subset L_{x}G$。

\begin{rmk} 对一般的水平结构 $\bK_{x}$，习惯上把 $\Bun_{G,\bK_{x}}$ 的 $S$-点说成「$X\times S$ 上带 $x$ 处 $\bK_{x}$-水平结构的 $G$-主丛」。然而这一术语可能引起误解：当 $\bK_{x}$ 不必包含在 $L^{+}_{x}G$ 中时，$\Bun_{G,\bK_{x}}$ 的 $S$-点并不诱导整个 $X\times S$ 上的 $G$-主丛，而只诱导 $(X-\{x\})\times S$ 上的 $G$-主丛。下面的例说明了这一点。
\end{rmk}



\begin{exam}\label{ex:SL parahoric} 设 $G=\SL_{n}$，并固定 $x\in |X|$。在共轭意义下，$L_{x}G$ 有 $n$ 个极大抛孔子群 $\bP^{(i)}_{x}$，由 $i\in \ZZ/n\ZZ$ 参数化。若令 $\wt G=\GL_{n}$，则可取 $\bP^{(i)}_{x}=g_{i}(L^{+}_{x}\wt G)g_{i}^{-1}\cap L_{x}G$，其中 $g_{i}$ 是 $\wt G(F_{x})$ 中任一满足 $\val_{F_{x}}(\det(g_{i}))\equiv i\pmod{n}$ 的元素。

对 $\Bun_{G, \bP^{(i)}_{x}}$ 有如下解释。考虑把 $\cV$ 送到 $\wedge^{n}(\cV)$ 的映射 $\det: \Bun_{\wt G}\to \Pic_{X}$。则存在叠的典范同构
\begin{equation*}
\Bun_{G, \bP^{(i)}_{x}}\cong {\det}^{-1}(\cO(-ix)), \quad i\in \ZZ.
\end{equation*}
当然，与 $\cO(x)$ 作张量积会把 $\det^{-1}(\cO(-ix))$ 与 $\det^{-1}(\cO(-(i-n)x))$ 等同起来，故上式右端的同构类只依赖于 $i$ 在 $\ZZ/n\ZZ$ 中的像。

换言之，对 $i\in \ZZ$，$\Bun_{G, \bP^{(i)}_{x}}$ 的 $S$-点分类偶 $(\cV,\l)$，其中 $\cV$ 是 $X\times S$ 上秩 $n$ 的向量丛，$\l$ 是同构 $\cO(-ix)\isom \wedge^{n}(\cV)$。
\end{exam}


上述定义显然可推广到在不止一个、而是有限多个点处给水平结构的情形。若 $\Sig\subset |X|$ 是有限子集，$\bK:=(\bK_{x})_{x\in \Sigma}$ 是每个 $x\in\Sig$ 处水平群的族，则把所得的模叠记作 $\Bun_{G,\bK}$。



\subsection{$G$-主丛的 Hecke 叠}
以下固定有限个位点组成的集合 $\Sigma\subset |X|$，并对每个 $x\in \Sigma$ 固定一个水平群 $\bK_{x}$。令 $U=X-\Sigma$，用 $\bK$ 记水平群族 $(\bK_{x})_{x\in \Sigma}$。于是模叠 $\Bun_{G,\bK}$ 有定义。


\ssszh{Hecke 叠}\label{sss:Hk}

设 $I$ 是有限集。令 $\Hk^{I}_{G,\bK}: \Sch_{k}\to \Gpd$ 为如下函子：它在 $k$-概形 $S$ 处的取值是多元组 $\xi=((x_{i})_{i\in I}, \cE_{0},\cE_{1} ,\a)$ 的群胚，其中
\begin{enumerate}
\item 对每个 $i\in I$，$x_{i}: S\to U$ 是态射，其图像为 $\G(x_{i})\subset X\times S$。这些 $x_{i}$ 称为 $\xi$ 的 {\em 腿}（leg）。
\item $\cE_{0},\cE_{1}\in \Bun_{G,\bK}(S)$，即 $\cE_{0}$ 与 $\cE_{1}$ 是 $X\times S$ 上的 $G$-主丛，并沿 $\{x\}\times S$ 带有 $\bK_{x}$-水平结构。
\item $\a$ 是带 $\bK$-水平结构的 $G$-主丛之间的同构 $\cE_{0}|_{X\times S-\cup_{i\in I}\G(x_{i})}\isom \cE_{1}|_{X\times S-\cup_{i\in I}\G(x_{i})}$。
\end{enumerate}

按定义，$\Hk_{G,\bK}^{I}$ 带有映射
\begin{equation*}
\Hk^{I}_{G,\bK}\to U^{I}.
\end{equation*}

当 $I=\vn$ 时，$\Hk^{I}_{G,\bK}=\Bun_{G,\bK}$。当 $\Sig=\vn$（即 $\bK$ 无数据）时，把 $\Hk^{I}_{G,\bK}$ 简写作 $\Hk^{I}_{G}$。


\ssszh{Hecke 叠的迭代版本}\label{sss:it Hk}
设 $I=I_{1}\coprod \cdots\coprod I_{r}$ 是有限集 $I$ 的分解。把同构 $\a$ 细化为 $r$ 步，便得到 Hecke 叠 $\Hk^{I}_{G,\bK}$ 的变体 $\Hk^{(I_{1},\cdots, I_{r})}_{G,\bK}$。$\Hk^{(I_{1},\cdots, I_{r})}_{G,\bK}$ 的 $S$-点由多元组 $((x_{i})_{i\in I}, (\cE_{j})_{0\le j\le r} ,(\a_{j})_{1\le j\le r})$ 组成，其中
\begin{enumerate}
\item 对每个 $i\in I$，$x_{i}: S\to U$ 是态射，其图像为 $\G(x_{i})\subset X\times S$。
\item $\cE_{0},\cdots, \cE_{r}\in \Bun_{G,\bK}(S)$。
\item 对 $1\le j\le r$，$\a_{j}$ 是带 $\bK$-水平结构的 $G$-主丛之间的同构
$$\a_{j}: \cE_{j-1}|_{X\times S-\cup_{i\in I_{j}}\G(x_{i})}\isom \cE_{j}|_{X\times S-\cup_{i\in I_{j}}\G(x_{i})}.$$
\end{enumerate}

记录丛 $\cE_{j}$ 给出态射
\begin{equation*}
p_{j}: \Hk^{(I_{1},\cdots, I_{r})}_{G,\bK}\to \Bun_{G,\bK}, \quad j=0,1,\cdots, r.
\end{equation*}

忘掉中间丛 $\cE_{1},\cdots, \cE_{r-1}$，并取复合 $\a=\a_{r}\c\cdots \c\a_{1}$，便得到映射
\begin{equation*}
\pi^{(I_{1},\cdots, I_{r})}_{I}: \Hk^{(I_{1},\cdots, I_{r})}_{G,\bK}\to \Hk^{I}_{G,\bK}.
\end{equation*}
该映射在开子概形 $U^{(I_{1},\cdots, I_{r})}_{\disj}\subset U^{I}$ 上是同构；此开子概形由满足下述条件的 $(x_{i})_{i\in I}$ 组成：对属于不同部分 $I_{j}$ 与 $I_{j'}$ 的 $i$ 与 $i'$，$x_{i}$ 与 $x_{i'}$ 的图像不相交。

更一般地，若每个 $I_{j}$ 都有分解 $I_{j}=I_{j,1}\coprod \cdots\coprod I_{j,s_{j}}$，则 $I$ 有更细的分解 $I=I_{1,1}\coprod \cdots\coprod I_{1,s_{1}}\coprod I_{2,1}\coprod\cdots\coprod I_{2,s_{2}}\coprod \cdots\coprod I_{r,s_{r}}$。把这个加细记作 $I=I'_{1}\coprod \cdots\coprod I'_{s}$（故 $s=s_{1}+\cdots+s_{r}$），其中各部分的顺序与上面相同。我们得到映射
\begin{equation*}
\pi^{(I'_{1},\cdots, I'_{s})}_{(I_{1},\cdots, I_{r})}: \Hk^{(I'_{1},\cdots, I'_{s})}_{G,\bK}\to \Hk^{(I_{1},\cdots, I_{r})}_{G,\bK}
\end{equation*}
它只记录丛 $\cE_{0},\cE_{s_{1}},\cE_{s_{1}+s_{2}},\cdots, \cE_{s}$ 以及复合 $\a_{j}:=\a_{j,s_{j}}\c\cdots\c\a_{j,1}: \cE_{j-1}\dr\cE_{j}$，$1\le j\le r$。该映射在 $\prod_{j=1}^{r}U^{(I_{j,1},\cdots, I_{j,s_{j}})}_{\disj}\subset \prod_{j=1}^{r}U^{I_{j}}=U^{I}$ 上是同构。



\ssszh{仿射 Grassmann 簇}\label{sss:Gr}
关于这一主题的更多细节，读者可参阅 X. Zhu 的综述文章 \cite{Zhu-Satake}。$G$ 的仿射 Grassmann 簇 $\Gr_{G}$ 按定义是 $k$ 上的 ind-概形，它表示 $k$-代数上函子 $R\mapsto (LG)(R)/(L^{+}G)(R)=G(R\lr{t})/G(R\tl{t})$ 的 fpqc 层化。或者，若令 $D_{R}=\Spec R\tl{t}$、$D^{\times}_{R}=\Spec R\lr{t}$，则 $\Gr_{G}(R)$ 分类偶 $(\cE, \io)$，其中 $\cE$ 是 $D_{R}$ 上的 $G$-主丛，$\io$ 是 $\cE|_{D^{\times}_{R}}$ 的平凡化，即同构 $\io: \cE|_{D^{\times}_{R}}\isom \cE_{\triv}|_{D^{\times}_{R}}$，这里 $\cE_{\triv}$ 是 $D_{R}$ 上的平凡 $G$-主丛。


设 $T\subset G$ 是分裂极大环面，Weyl 群为 $W$。回忆：仿射 Grassmann 簇 $\Gr_{G}$ 上的 $L^{+}G$-轨道由余权在 Weyl 群作用下的轨道 $\xcoch(T)/W$ 标号；对 $\l\in \xcoch(T)$，我们指定包含 $t^{\l}\in LT$ 的那条 $L^{+}G$-轨道，它只依赖于 $\l$ 的 $W$-轨道。

选取包含 $T$ 的 Borel 子群，它定义了支配余权 $\xcoch(T)^{+}$。于是有典范双射 $\xcoch(T)^{+}\isom \xcoch(T)/W$，这使我们可以用 $\xcoch(T)^{+}$ 中的元素唯一地标号 $\Gr_{G}$ 上的任一条 $L^{+}G$-轨道。对 $\l\in \xcoch(T)^{+}$，令 $\Gr_{G, \l}$ 为 $t^{\l}$ 所在的 $L^{+}G$-轨道，视作局部闭的约化子概形。我们有 $\dim\Gr_{G,\l}=\j{2\r,\l}$，其中 $2\r$ 是正根之和。

在 $\xcoch(T)^{+}$ 上定义偏序 $\le$ 如下：$\l'\le \l$ 当且仅当 $\l-\l'$ 是关于 $B$ 的正余根之和。已知 $\Gr_{G,\l}$ 的闭包为 $\Gr_{G,\le \l}=\cup_{\l'\in \xcoch(T)^{+}, \l'\le \l}\Gr_{G,\l'}$。


对闭点 $x\in |X|$，类似地定义 $\Gr_{G,x}$ 为表示 $R\mapsto \Gr_{G,x}(R)=(L_{x}G)(R)/(L^{+}_{x}G)(R)$ 的 fpqc 层化的 ind-概形。选定 $x$ 处的单值化子 $\vp_{x}$ 后，得到同构 $\Gr_{G,x}\cong \Gr_{G}\ot_{k}k_{x}$。当 $x$ 变动时，$\Gr_{G,x}$ 构成 $X$ 上的一个族，这是 Beilinson–Drinfeld Grassmann 簇的一个特例，下面就来回忆它。

\begin{exam} 当 $G=\GL_{n}$ 时，$\Gr_{G}(k)$ 是 $k\lr{t}^{\op n}$ 中秩 $n$ 的自由 $k\tl{t}$-子模的集合。设整数 $d_{1}\ge d_{2}\ge\cdots \ge d_{n}$，令 $\l=(d_{1},d_{2},\cdots, d_{n})$ 并视作对角环面 $T\subset G$ 的支配余权，则点 $t^{\l}$ 对应 $k\tl{t}$-子模
\begin{equation*}
t^{d_{1}}k\tl{t}\op t^{d_{2}}k\tl{t}\op\cdots\op t^{d_{n}}k\tl{t}\subset k\lr{t}^{\op n}.
\end{equation*}

当 $d_{n}\ge0$ 时，$\Gr_{G,\l}(k)$ 是 $k\tl{t}$-子模 $\L\subset k\tl{t}^{\op n}$ 的集合，使得作为 $k\tl{t}$-模有 $k\tl{t}^{\op n}/\L\cong \op_{i=1}^{n}k\tl{t}/(t^{d_{i}})$。

当 $\l=(d,0,\cdots, 0)$ 且 $d\ge0$ 时，$\Gr_{G,\le\l}(k)$ 是使 $\dim_{k}k\tl{t}^{\op n}/\L=d$（而非 $\le d$）的 $k\tl{t}$-子模 $\L\subset k\tl{t}^{\op n}$ 的集合。
\end{exam}



\ssszh{Beilinson–Drinfeld Grassmann 簇}\label{sss:BD}
令 $\cE_{\triv}\in \Bun_{G}(k)$ 为 $X$ 上的平凡 $G$-主丛。用笛卡儿方图
\begin{equation*}
\begin{tikzcd}
\Gr^{(I_{1},\cdots, I_{r})}_{G} \arrow[r] \arrow[d] & \Hk^{(I_{1},\cdots, I_{r})}_{G} \arrow[d, "p_{r}"] \\
\Spec k \arrow[r, "\cE_{\triv}"] & \Bun_{G}
\end{tikzcd}
\end{equation*}
定义 {\em Beilinson–Drinfeld Grassmann 簇} $\Gr^{(I_{1},\cdots, I_{r})}_{G}$。于是有腿映射 $\Gr^{(I_{1},\cdots, I_{r})}_{G}\to X^{I}$。当 $I$ 是单元集 $\{1\}$ 时，$\Gr^{\{1\}}_{G}\to X$ 在闭点 $x\in |X|$ 处的纤维典范同构于 $\Gr_{G,x}$。

\ssszh{相对位置映射}
可以构造商叠 $\QGr_{G}:=L^{+}G\bs \Gr_{G}$，它是 $\Hk^{\{1\}}_{G}$ 的局部版本：对 $k$-代数 $R$，$\QGr_{G}$ 的 $R$-点是三元组 $(\cE_{0}, \cE_{1}, \a)$，其中 $\cE_{i}$ 是 $D_{R}=\Spec R\tl{t}$ 上的 $G$-主丛，$\a: \cE_{0}|_{D^{\times}_{R}}\isom \cE_{1}|_{D^{\times}_{R}}$。于是 $\QGr_{G}$ 的点与 $\xcoch(T)^{+}$ 一一对应。对 $\l\in \xcoch(T)^{+}$，令 $\QGr_{G, \le \l}$ 为 $\Gr_{G,\le \l}$ 的像，它是 $\QGr_{G}$ 的闭约化子叠。

令 $\Aut(D)$ 为 $k\tl{t}$ 的自同构构成的 $k$ 上 pro-代数群。则 $\Aut(D)$ 通过坐标变换右作用在 $\QGr_{G}$ 上：$\th: D\to D$ 把 $(\cE_{0},\cE_{1}, \a)$ 送到 $(\th^{*}\cE_{0}, \th^{*}\cE_{1}, \th^{*}\a)$。显然 $\QGr_{G, \le\l}$ 在 $\Aut(D)$ 的作用下不变。


对 Hecke 叠 $\Hk^{\{1\}}_{G,\bK}$，如下定义求值映射
\begin{equation*}
\ev^{\{1\}}_{\Hk}: \Hk^{\{1\}}_{G,\bK}\to \QGr_{G}/\Aut(D)
\end{equation*}
只需构造一个 $\Aut(D)$-主丛 $\wt\Hk^{\{1\}}_{G,\bK}\to \Hk^{\{1\}}_{G,\bK}$，并配以 $\Aut(D)$-等变映射 $\wt\ev^{\{1\}}_{\Hk}: \wt\Hk^{\{1\}}_{G,\bK}\to \QGr_{G}$。定义 $\wt\Hk^{\{1\}}_{G,\bK}$ 为如下模叠：其 $R$-点为 $(x, t_{x}, \cE_{0}, \cE_{1}, \a)$，其中 $(x,\cE_{0},\cE_{1}, \a)\in\Hk^{\{1\}}_{G,\bK}(R)$，而 $t_{x}$ 是 $X\ot R$ 沿 $\G(x)$ 的形式完备化上的正则函数，且对 $s\in \Spec R$ 的每个几何点，$t_{x}$ 都是 $\wh\cO_{X\times s, x_{s}}$ 的单值化子。于是 $t_{x}$ 给出连续同构 $\wh \cO_{X\ot R, \G(x)}\cong R\tl{t}$，从而给出映射 $t^{\sh}_{x}: \Spec R\tl{t}\to X\ot R$。映射 $\wt\ev^{\{1\}}_{\Hk}$ 便把 $(x, t_{x}, \cE_{0}, \cE_{1}, \a)$ 送到 $(t^{\sh, *}_{x}\cE_{0}, t^{\sh, *}_{x}\cE_{1}, t^{\sh, *}_{x}\a)$。

\ssszh{有界的 Hecke 叠}\label{sss:bd Hk}
对 $\l\in \xcoch(T)^{+}$，定义闭子叠
\begin{equation*}
\Hk^{\{1\}, \le \l}_{G,\bK}\subset \Hk^{\{1\}}_{G,\bK}
\end{equation*}
为 $\QGr_{G,\le \l}/\Aut(D)\subset \QGr_{G}/\Aut(D)$ 在上述映射 $\ev^{\{1\}}_{\Hk}$ 下的原像。粗略地说，$\Hk^{\{1\}}_{G,\bK}$ 的 $S$-点 $(x,\cE_{0},\cE_{1}, \a)$ 落在 $\Hk^{\{1\}, \le \l}_{G,\bK}$ 中，当且仅当 $\a$ 的相对位置 $\le \l$。

现在设有限集 $I$ 有分解 $I=I_{1}\coprod I_{2}\coprod \cdots\coprod I_{r}$。令 $\un\l=(\l_{i})_{i\in I}\in (\xcoch(T)^{+})^{I}$。我们想定义 $\Hk^{(I_{1},\cdots, I_{r})}_{G,\bK}$ 的闭子叠 $\Hk^{(I_{1},\cdots, I_{r}), \le \un\l}_{G,\bK}$。构造分几步完成：
\begin{enumerate}
\item 令 $U^{I_{j}}_{\disj}\subset U^{I_{j}}$ 为使 $\G(x_{i})$ 两两不相交的 $(x_{i})_{i\in I_{j}}$ 构成的开子集。则有映射
\begin{equation}\label{ev disj}
\ev^{(I_{1},\cdots, I_{r})}_{\Hk,\disj}: \Hk^{(I_{1},\cdots, I_{r})}_{G,\bK}|_{U^{I_{1}}_{\disj}\times\cdots\times U^{I_{r}}_{\disj}}\to (\QGr_{G}/\Aut(D))^{I}
\end{equation}
其第 $i$ 个分量记录 $\a_{j}$（其中 $i\in I_{j}$）沿 $\G(x_{i})$ 的相对位置。定义
\begin{equation*}
\Hk^{(I_{1},\cdots, I_{r}), \le \un\l}_{G,\bK, \disj}\subset \Hk^{(I_{1},\cdots, I_{r})}_{G,\bK}|_{U^{I_{1}}_{\disj}\times\cdots\times U^{I_{r}}_{\disj}}
\end{equation*}
为 $\prod_{i\in I}\QGr_{G,\le \l_{i}}/\Aut(D)$ 在 $\ev^{(I_{1},\cdots, I_{r})}_{\Hk,\disj}$ 下的原像。


\item 接着定义 $\Hk^{(I_{1},\cdots, I_{r}), \le \un\l}_{G,\bK}$ 为 $\Hk^{(I_{1},\cdots, I_{r}), \le \un\l}_{G,\bK, \disj}$ 在 $\Hk^{(I_{1},\cdots, I_{r})}_{G,\bK}$ 中的闭包。
\end{enumerate}

定义 Beilinson–Drinfeld Grassmann 簇的有界版本 $\Gr^{(I_{1},\cdots, I_{r}), \le\un\l}_{G}$ 为 $\Hk^{(I_{1},\cdots, I_{r}), \le\un\l}_{G}$ 在典范映射 $\Gr^{(I_{1},\cdots, I_{r})}_{G}\to \Hk^{(I_{1},\cdots, I_{r})}_{G}$ 下的原像。

\begin{exam} 设 $G=\GL_{n}$，并对所有 $i\in I$ 取 $\l_{i}=(1,0,\cdots, 0)$。
\begin{enumerate}
\item 迭代版本。设 $I=\{1,\cdots, r\}$ 分解为单点集 $I_{i}=\{i\}$。则 $\Hk^{(\{1\},\cdots,\{ r\}),\le\un\l}_{G}$ 的 $S$-点是
$$((x_{i})_{1\le i\le r}, (\cV_{j})_{0\le j\le r}, (\a_{j})_{1\le j\le r}),$$
其中 $\cV_{j}$ 是 $X\times S$ 上秩 $n$ 的向量丛，$\a_{j}: \cV_{j-1}\incl \cV_{j}$ 是凝聚层的单射，其余核在图像 $\G(x_{j})$ 上局部自由且秩为 $1$，$1\le j\le r$。

\item 非迭代版本。$\Hk^{I,\le\un\l}_{G}$ 的 $S$-点是 $((x_{i})_{i\in I}, \cV_{0}, \cV_{1}, \a)$，其中 $\cV_{0}$ 与 $\cV_{1}$ 是 $X\times S$ 上秩 $n$ 的向量丛，$\a: \cV_{0}\incl \cV_{1}$ 是凝聚层的单射，使得 $\det(\a): \det(\cV_{0})\to \det(\cV_{1})$ 的除子作为 $X\times S$ 上的除子是 $\sum_{i\in I}\G(x_{i})$。
\end{enumerate}

\end{exam}


\begin{rmk} 我们的记号与 \cite[\S1]{VLaff} 不同。Hecke 叠 $\Hk^{(I_{1},\cdots, I_{r})}_{G,\bK}$ 在 {\em 同上} 中记作 $\textup{Hecke}_{\bK,I}^{(I_{1},\cdots, I_{r})}$。在 \cite[定义 1.2]{VLaff} 中引入了迭代 Hecke 叠的两个有界版本，分别记作 $\textup{Hecke}^{(I_{1},\cdots, I_{r})}_{\bK, I, \un\l}$ 与 $\textup{Hecke}^{(I_{1},\cdots, I_{r})}_{\bK, I, \lesssim\un\l}$。可以证明，我们的 Hecke 叠 $\Hk^{(I_{1},\cdots, I_{r}),\le\un\l}_{G, \bK}$ 是 $\textup{Hecke}^{(I_{1},\cdots, I_{r})}_{\bK, I, \un\l}$ 的约化结构，而后者又包含于 $\textup{Hecke}^{(I_{1},\cdots, I_{r})}_{\bK, I, \lesssim\un\l}$。
\end{rmk}

%%%%%%%%%%%%%%
\subsection{$G$-shtuka}\label{ss:Sht}
对 $k=\FF_{q}$ 上任意概形 $Y$，用 $\Fr_{Y}$ 或简写 $\Fr: Y\to Y$ 表示把正则函数取 $q$ 次幂的 Frobenius 态射。

仍固定有限个位点组成的集合 $\Sigma\subset |X|$ 以及水平群 $\bK=(\bK_{x})_{x\in \Sig}$。令 $U=X-\Sigma$。

\ssszh{$\Sht_{G}$ 的定义}
设有限集 $I$ 有分解 $I=I_{1}\coprod \cdots\coprod I_{r}$。用笛卡儿方图
\begin{equation}\label{defn Sht}
\begin{tikzcd}
\Sht_{G,\bK}^{(I_{1},\cdots, I_{r})} \arrow[r] \arrow[d] & \Hk_{G,\bK}^{(I_{1},\cdots, I_{r})} \arrow[d, "{({p_0},{p_r})}"'] \\
\Bun_{G,\bK} \arrow[r, "{(\id,\Fr)}"] & \Bun_{G,\bK}\times \Bun_{G,\bK}
\end{tikzcd}
\end{equation}
定义 Artin 叠 $\Sht_{G,\bK}^{(I_{1},\cdots, I_{r})}$。

由构造，有记录腿的映射
\begin{equation*}
\pi:=\pi^{(I_{1},\cdots, I_{r})}_{G,\bK}: \Sht_{G,\bK}^{(I_{1},\cdots, I_{r})}\to U^{I}.
\end{equation*}


\ssszh{点函子}\label{sss:Sht points}
按定义，对 $k$ 上任意测试概形 $S$，$\Sht_{G,\bK}^{I}$ 的 $S$-点是多元组
\begin{equation*}
\xi=((x_{i})_{i\in I}, \cE, \a: \cE\dr {}^{\t}\cE)
\end{equation*}
的群胚，其中
\begin{itemize}
\item 对 $i\in I$，$x_{i}\in U(S)$，其图像为 $\G(x_{i})\subset U\times S$。这些 $x_{i}$ 称为 shtuka $\xi$ 的 {\em 腿}。
\item $\cE$ 是 $X\times S$ 上的 $G$-主丛，并沿 $\{x\}\times S$ 带有 $\bK_{x}$-水平结构。
\item ${}^{\t}\cE=(\id_{X}\times \Fr_{S})^{*}\cE$。
\item $\a$ 是保持 $\bK$-水平结构的 $G$-主丛同构 $\cE|_{X\times S-\cup_{i\in I}\G(x_{i})}\isom ({}^{\t}\cE)|_{X\times S-\cup_{i\in I}\G(x_{i})}$。
\end{itemize}


更一般地，对分解 $I=I_{1}\coprod \cdots\coprod I_{r}$，$\Sht_{G,\bK}^{(I_{1},\cdots, I_{r})}$ 的 $S$-点是多元组
\begin{equation*}
((x_{i})_{i\in I},(\cE_{j})_{0\le j\le r}, (\a_{j}: \cE_{j-1}\dr \cE_{j})_{1\le j\le r}, \io)
\end{equation*}
的群胚，其中
\begin{itemize}
\item 对 $i\in I$，$x_{i}\in U(S)$，其图像为 $\G(x_{i})\subset X\times S$。
\item 对 $0\le j\le r$，$\cE_{j}$ 是 $X\times S$ 上的 $G$-主丛，并沿 $\{x\}\times S$ 带有 $\bK_{x}$-水平结构。
\item 对 $1\le j\le r$，$\a_{j}$ 是保持 $\bK$-水平结构的 $G$-主丛同构
\begin{equation*}
\cE_{j-1}|_{X\times S-\cup_{i\in I_{j}}\G(x_{i})}\isom \cE_{j}|_{X\times S-\cup_{i\in I_{j}}\G(x_{i})}.
\end{equation*}
\item $\io$ 是保持 $\bK$-水平结构的 $G$-主丛同构
\begin{equation*}
\io:\cE_{r}\isom {}^{\t}\cE_{0}=(\id_{X}\times \Fr_{S})^{*}\cE_{0}.
\end{equation*}
\end{itemize}

\begin{exam}\label{ex:Sht no leg}
当 $I=\vn$ 时，按定义 $\Sht^{\vn}_{G,\bK}$ 是对角态射 $\D: \Bun_{G,\bK}\to \Bun_{G,\bK}^{2}$ 与 Frobenius 图像 $(\id,\Fr): \Bun_{G,\bK}\to \Bun_{G,\bK}^{2}$ 的纤维积。一般事实是：对 $k$ 上任意 Artin 叠 $\frX$，$\D: \frX\to \frX^{2}$ 与 Frobenius 图像 $(\id,\Fr):\frX\to \frX^{2}$ 的纤维积同构于群胚 $\frX(k)$，视作 $k$ 上的离散叠。在我们的情形有
\begin{equation*}
\Sht^{\vn}_{G,\bK}=\coprod_{\cE\in \Bun_{G,\bK}(k)}[(\Spec k)/\Aut(\cE)(k)].
\end{equation*}
\end{exam}



\ssszh{有界版本}
令 $\un\l=(\l_{i})_{i\in I}\in (\xcoch(T)^{+})^{I}$。用笛卡儿方图
\begin{equation*}
\begin{tikzcd}
\Sht_{G,\bK}^{(I_{1},\cdots, I_{r}),\le\un\l}\arrow[r]\arrow[d] & \Hk_{G,\bK}^{(I_{1},\cdots, I_{r}),\le\un\l}\arrow[d, "{(p_{0},p_{r})}"] \\
\Bun_{G,\bK}\arrow[r, "{(\id,\Fr)}"] & \Bun_{G,\bK}\times \Bun_{G,\bK}
\end{tikzcd}
\end{equation*}
定义 $\Sht^{(I_{1},\cdots, I_{r})}_{G,\bK}$ 的有界版本 $\Sht^{(I_{1},\cdots, I_{r}), \le \un\l}_{G,\bK}$。于是 $\Sht^{(I_{1},\cdots, I_{r}), \le \un\l}_{G,\bK}$ 是 $\Sht_{G,\bK}^{(I_{1},\cdots, I_{r})}$ 的闭子叠。



\begin{rmk} 我们的记号与 \cite[定义 2.1]{VLaff} 不同。叠 $\Sht^{(I_{1},\cdots, I_{r})}_{G,\bK}$ 在 {\em 同上} 中记作 $\textup{Cht}_{\bK,I}^{(I_{1},\cdots, I_{r})}$。至于有界版本，叠 $\textup{Cht}_{\bK,I, \un\l}^{(I_{1},\cdots, I_{r})}$ 是在 {\em 同上} 中引入的，并且存在闭嵌入 $\Sht^{(I_{1},\cdots, I_{r}), \le \un\l}_{G,\bK}\incl \textup{Cht}_{\bK,I, \un\l}^{(I_{1},\cdots, I_{r})}$，它在取约化结构后诱导同构。
\end{rmk}


\begin{exam}\label{ex:Gm} 设 $G=\GL_{1}=\Gm$，$I$ 为有限集。此时一个界由整数列 $\un\l:=(\l_{i})_{i\in I}\in \ZZ^{I}$ 给出。注意 $\ZZ^{I}$ 上的偏序是平凡的：$\l'\le \l$ 当且仅当 $\l'=\l$。叠 $\Sht^{I, \le \un\l}_{\Gm}$ 另有如下描述：它适合笛卡儿方图
\begin{equation*}
\begin{tikzcd}
\Sht^{I, \le \un\l}_{\Gm} \arrow[r]\arrow[d] & \Pic_{X}\arrow[d, "\bL_{\Pic_{X}}"] \\
X^{I}\arrow[r, "\AJ_{\un\l}"] & \Pic_{X}
\end{tikzcd}
\end{equation*}
这里 $\bL_{\Pic_{X}}$ 是 $\Pic_{X}$ 的 Lang 同源，它把线丛 $\cL$ 送到 ${}^{\t}\cL\ot \cL^{-1}$。Abel–Jacobi 映射 $\AJ_{\un\l}$ 把 $(x_{i})_{i\in I}$ 送到线丛 $\cO_{X}(\sum_{i\in I}\l_{i}x_{i})$。

由于 $\bL_{\Pic_{X}}$ 的像落在 $\Pic^{0}_{X}$ 中，而 $\AJ_{\un\l}$ 落在 $\Pic^{|\un\l|}_{X}$ 中（其中 $|\un\l|=\sum_{i\in I}\l_{i}$），可见当 $|\un\l|\ne0$ 时 $\Sht^{I, \le \un\l}_{G} =\vn$。
\end{exam}

\begin{exam}\label{ex:Dr} 设 $G=\GL_{n}$，并设 $D=\sum_{x\in \Sig}m_{x}x$ 是 $X$ 上的有效除子。令 $\bK_{D}$ 为 $x\in \Sig$ 处主同余子群 $\bK_{x,m_{x}}$ 构成的族。令 $I=\{1,2\}$，并取分解 $I_{1}=\{1\}$、$I_{2}=\{2\}$。令 $\l_{1}=(0,\cdots, 0,-1)$ 与 $\l_{2}=(1,0,\cdots,0)$ 为两个极小余权，$\un\l=(\l_{1},\l_{2})$。则 $\Sht^{(I_{1},I_{2}), \le\un\l}_{G,\bK_{D}}(S)$ 分类如下数据
\begin{equation*}
\begin{tikzcd}
\cV_{0} & \cV_{1}\arrow[l, "\a_{1}"'] \arrow[r, "\a_{2}"] & \cV_{2}\arrow[r, "\io"{above}, "{\sim}"{below}] & {}^{\t}\cV_{0}
\end{tikzcd}
\end{equation*}
其中 $\cV_{i}$ 是 $X\times S$ 上秩 $n$ 的向量丛，并沿 $D\times S$ 带有平凡化；$\a_{1}$ 与 $\a_{2}$ 是与 $D\times S$ 上的平凡化相容的凝聚层单射，其余核分别在映射 $x_{1}: S\to U=X-\Sig$ 与 $x_{2}: S\to U$ 的图像上局部自由。这正是 Drinfeld 在 \cite{Dr-FBun} 中最初考虑的那类模叠。
\end{exam}

\ssszh{容许性}
令 $\Om=\xcoch(T)/\ZZ\Phi^{\vee}$，其中 $\ZZ\Phi^{\vee}$ 是余根格。则 $\Om$ 是 $\Gr_{G}$（以及 $\QGr_{G}$）的连通分支群。由 \cite[引理 2.2]{Var}，连通分支集 $\pi_{0}(\Bun_{G})$ 也可典范地等同于 $\Om$。更一般地，若每个 $\bK_{x}$ 都包含在 $L_{x}G$ 的某个抛孔子群中，则 $\pi_{0}(\Bun_{G,\bK})$ 仍可典范地等同于 $\Om$，其中平凡丛所在的分支对应 $0\in \Om$。


令 $\un\l=(\l_{i})_{i\in I}\in (\xcoch(T)^{+})^{I}$，$\om$ 为 $\sum_{i\in I}\l_{i}\in \xcoch(T)$ 在 $\Om$ 中的像。若 $\om=0\in \Om$，则称 $\un\l$ 是 {\em $G$-容许的}（见 \cite[定义 2.14]{Var}）。

\begin{lem}[Varshavsky {\cite[命题 2.16(d)]{Var}}] 设对每个 $x\in \Sig$，$\bK_{x}$ 都包含在某个抛孔子群中。设 $I=I_{1}\coprod\cdots\coprod I_{r}$ 是有限集 $I$ 的分解，$\un\l=(\l_{i})_{i\in I}\in (\xcoch(T)^{+})^{I}$。

则 $\Sht^{(I_{1},\cdots, I_{r}), \le\un\l}_{G,\bK}$ 非空当且仅当 $\un\l$ 是 $G$-容许的。
\end{lem}
\begin{proof}[证明概要]

对 $\g\in \Om\cong \pi_{0}(\Bun_{G,\bK})$，以 $\Bun_{G, \bK}^{\g}$ 记 $\Bun_{G,\bK}$ 相应的连通分支。把映射 $(p_{0},p_{r}): \Hk^{(I_{1},\cdots, I_{r})}_{G,\bK}\to \Bun_{G,\bK}\times\Bun_{G,\bK}$ 限制到 $\Hk^{(I_{1},\cdots, I_{r}), \le\un\l}_{G,\bK}$ 上，得
\begin{equation*}
(p^{\le\un\l}_{0},p^{\le\un\l}_{r}):  \Hk^{(I_{1},\cdots, I_{r}), \le\un\l}_{G,\bK}\to \coprod_{\g\in \Om} \Bun_{G,\bK}^{\g}\times \Bun^{\g+\om}_{G,\bK}
\end{equation*}
其中 $\om$ 是 $\sum\l_{i}$ 在 $\Om$ 中的像。

若 $\om\ne0$，则 $(p^{\le\un\l}_{0},p^{\le\un\l}_{r})$ 的像与 $\Bun_{G,\bK}^{\g}\times \Bun^{\g}_{G,\bK}$（对任意 $\g\in \Om$）都不相交。然而 $\Bun_{G,\bK}$ 上 Frobenius 的图像落在 $\coprod_{\g\in \Om} \Bun_{G,\bK}^{\g}\times \Bun^{\g}_{G,\bK}$ 之上。由笛卡儿方图 \eqref{defn Sht} 即知 $\Sht^{(I_{1},\cdots, I_{r}), \le\un\l}_{G,\bK}=\vn$。

反之，设 $\un\l$ 是 $G$-容许的。由于 $\Hk^{(I_{1},\cdots, I_{r}),\le\un\l}_{G, \bK}\to \Hk^{I,\le\un\l}_{G,\bK}$ 满射，映射 $\Sht^{(I_{1},\cdots, I_{r}),\le\un\l}_{G, \bK}\to \Sht^{I,\le\un\l}_{G,\bK}$ 也满射。因此只需证明 $\Sht^{I,\le\un\l}_{G,\bK}\ne\vn$。

令 $\l=\sum_{i}\l_{i}$。因 $\un\l$ 容许，支配余权 $\l$ 落在 $G$ 的余根格中，故 $0\le \l$。把 $\Sht^{I,\le\un\l}_{G,\bK}$ 限制到对角线 $\D(U)\subset U^{I}$，得 $\Sht^{I,\le\un\l}_{G,\bK}|_{\D(U)}\cong \Sht^{\{1\},\le\l}_{G,\bK}$。这一点将在 \S\ref{ss:geom fact} 中讨论。于是只需证明 $\Sht^{\{1\},\le\l}_{G,\bK}\ne\vn$。因 $0\le \l$，有 $\Sht^{\{1\},\le0}_{G,\bK}\subset \Sht^{\{1\},\le\l}_{G,\bK}$。显然 $\Sht^{\{1\},\le0}_{G,\bK}=\Sht^{\vn}_{G,\bK}\times U$，由例 \ref{ex:Sht no leg} 知它非空。
\end{proof}

% ch02  §2.4 / §3.1–§3.3 / §4.1–§4.2（原文 963–1295 行）

\subsection{Iso-shtuka}

回忆：$F$ 是 $X$ 的函数域。

\ssszh{iso-shtuka 的张量范畴}
设 $\IsoSht(F)$ 为如下二元组 $(V, \phi)$ 构成的范畴：
\begin{itemize}
\item $V$ 是有限维 $\br F:=F\ot_{k} \ov k$-向量空间。
\item $\phi: V\isom V$ 是 $F$-线性且 $(\ov k, \Fr)$-线性的同构。也就是说，$\phi$ 是双射，且对 $v\in V$、$c\in \ov k$ 与 $a\in F$，有
$\phi(cv)=c^{q}\phi(v)$ 与 $\phi(av)=a\phi(v)$。
\end{itemize}
我们把 $\IsoSht(F)$ 中的对象称为函数域 $F$ 上的\emph{iso-shtuka}。$(V,\phi)\in \IsoSht(F)$ 的秩定义为 $\dim_{\br F}V$。于是 $\IsoSht(F)$ 在运算 $(V_{1},\phi_{1})\ot (V_{2}, \phi_{2})=(V_{1}\ot_{\br F} V_{2}, \phi_{1}\ot\phi_{2})$ 下成为刚性 $F$-线性张量范畴。$(V,\phi)$ 的对偶对象是 $(V^{\vee}=\Hom_{\br F}(V,\br F), \phi^{\vee})$，其中 $\phi^{\vee}$ 由 $\j{\phi^{\vee}(\xi),\phi(v)}=\j{\xi, v}$（$\xi\in V^{\vee}$，$v\in V$）定义。

我们应把 $F$ 上的 iso-shtuka 看作 $X$ 的通用点上的 shtuka。更确切地说，每个 $\ov k$-点 $((x_{i})_{i\in I}, \cE, \a)\in \Sht^{I}_{\GL_{n},\bK}(\ov k)$ 都对应一个如下定义的 iso-shtuka。设 $V$ 是 $\cE$ 在 $X_{\ov k}$ 的通用点 $\y=\Spec \br F$ 处的纤维。则 $\a$ 限制到 $\y$ 上给出 $\br F$-线性同构 $\a: V\isom V\ot_{\ov k,\Fr}\ov k$。对 $v\in V$ 定义 $\phi(v)=\a^{-1}(v\ot 1)$。则 $(V,\phi)$ 是 $\IsoSht(F)$ 中秩为 $n$ 的对象。

对 $F$ 上的约化群 $G$，我们把\emph{$G$-iso-shtuka} 定义为 $F$-线性张量函子
\begin{equation}\label{GisoSh}
\Rep_{F}(G)\to \IsoSht(F).
\end{equation}
若 $G$ 是分裂的，则 $\ov k$-点 $((x_{i})_{i\in I}, \cE, \a)\in \Sht^{I}_{G,\bK}(\ov k)$ 给出一个 $G$-iso-shtuka，其定义如下：它把 $G$ 的有限维 $F$-表示 $W$ 送到 iso-shtuka $(W_{\cE,\y}, \phi_{W})$，其中 $W_{\cE}=(\cE|_{U_{\ov k}})\times^{G}W$ 是由 $\cE$ 与 $W$ 相伴的 $U_{\ov k}$ 上的向量丛，$W_{\cE,\y}$ 是它的通用纤维；而 $\phi_{W}$ 是 $\a_{W}: W_{\cE}\dr {}^{\t}(W_{\cE})\cong W_{{}^{\t}\cE}$ 在 $X_{\ov k}$ 的通用点 $\y$ 处限制之逆。

\ssszh{等晶（isocrystal）}
显然，iso-shtuka 是等晶的整体类比物。设 $x\in |X|$，令 $\br F_{x}:=\wh{F_{x}^{\ur}}$ 为 $F_{x}$ 的最大非分歧扩域之完备化。设 $\Fr_{x}\in\Aut(\br F_{x})$ 是 $k_{x}$ 的自同构 $a\mapsto a^{\# k_{x}}$ 的提升。回忆：关于 $\Fr_{x}$ 的、$\br F_{x}$ 上等晶的范畴 $\IsoCrys(F_{x})$ 由如下二元组 $(V, \phi)$ 组成：
\begin{itemize}
\item $V$ 是有限维 $\br F_{x}$-向量空间。
\item $\phi: V\isom V$ 是 $(\br F_{x}, \Fr_{x})$-线性同构。也就是说，$\phi$ 是双射，且对 $v\in V$、$a\in \br F_{x}$，有
$\phi(av)=\Fr_{x}(a)\phi(v)$。
\end{itemize}
于是 $\IsoCrys(F_{x})$ 是 $F_{x}$-线性刚性张量范畴。

给定 iso-shtuka $(V,\phi)$ 与位置 $x\in|X|$，可以通过基变换作出 $\br F_{x}$ 上的等晶
\begin{equation*}
(V,\phi)\mapsto (V\ot_{\br F}\br F_{x}, \phi^{d_{x}}\ot \Fr_{x}).
\end{equation*}
这里 $d_{x}=[k_{x}:k]$。上述构造给出 $F$-线性张量函子
\begin{equation}\label{Sht fib functor}
\om_{x}: \IsoSht(F)\to \IsoCrys(F_{x}).
\end{equation}

\ssszh{单对象的分类}
回忆等晶的 Dieudonné–Manin 分类：$\IsoCrys(F_{x})$ 是半单 $F_{x}$-线性张量范畴，其单对象由斜率 $\l\in \QQ$ 分类。斜率为 $\l=m/n$（$n\ge1$，$m\in \ZZ$ 且与 $n$ 互素）的单对象 $M(\l)\in \IsoCrys(F_{x})$ 在 $\br F_{x}$ 上的维数为 $n$。对 $\l_{1}, \l_{2}\in \QQ$，$M(\l_{1})\ot M(\l_{2})$ 同构于若干个 $M(\l_{1}+\l_{2})$ 的直和。$M(\l)$ 的对偶同构于 $M(-\l)$。$M(\l)$ 的自同态代数是 $F_{x}$ 上的中心可除代数，其不变量为 $-\l\mod\ZZ\in \QQ/\ZZ$。

Drinfeld \cite[\S2]{Dr-Pet} 描述了 $\IsoSht(F)$ 作为 Abel 范畴的结构。为陈述它，设 $F^{\sep}$ 是 $F$ 的可分闭包，考虑带 $\G_{F}:=\Gal(F^{\sep}/F)$ 作用的 $\QQ$-向量空间 $F^{\sep,\times}\ot_{\ZZ}\QQ$（通常简写为 $F^{\sep,\times}\ot\QQ$）。对每个元素 $a\in F^{\sep,\times}\ot\QQ$，存在唯一的极小有限扩张 $L_{a}/F$（在 $F^{\sep}$ 内），使得 $a$ 落在 $L^{\times}_{a}\ot\QQ\to F^{\sep,\times}\ot\QQ$ 的像中。

\begin{thm}[Drinfeld \cite{Dr-Pet}]\label{th:isoSht}
\begin{enumerate}
\item 范畴 $\IsoSht(F)$ 半单。
\item $\IsoSht(F)$ 中单对象的同构类与 $F^{\sep,\times}\ot\QQ$ 上的 $\G_{F}$-轨道自然一一对应。
\item 设 $(V,\phi)$ 是 $\IsoSht(F)$ 中的单对象，其同构类对应于 $a\in F^{\sep,\times}\ot\QQ$ 的 $\G_{F}$-轨道。按上文定义 $L_{a}$，并记 $\div(a)=\sum_{y\in |L_{a}|}\ord_{y}(a)\cdot y$（对 $L_{a}$ 的位置求和）。则
\begin{enumerate}
\item $\End(V,\phi)$ 是 $L_{a}$ 上的中心可除代数，其在 $y\in |L_{a}|$ 处的不变量为 $-\ord_{y}(a)[k_{y}:k]\mod \ZZ\in \QQ/\ZZ$。
\item $\dim_{\br F}V=[L_{a}:F]d_{a}$，其中 $d_{a}\in \NN$ 是有理数集合 $\{\ord_{y}(a)[k_{y}:k]\}_{y\in |L_{a}|}$ 的最小公分母。
\item 由于 $L_{a}\subset \End(V,\phi)$，可以把 $(V,\phi)$ 看作 $\IsoSht(L_{a})$ 中的对象，于是把 $(V,\phi)$ 在位置 $y\in |L_{a}|$ 处局部化、得到 $\om_{y}(V,\phi)\in \IsoCrys(L_{a,y})$ 是有意义的。此时 $\om_{y}(V,\phi)$ 是若干斜率为 $\ord_{y}(a)$ 的单等晶的直和。
\end{enumerate}

\end{enumerate}
\end{thm}

我们给出定理第 (2) 部分一个方向的构造。由单对象 $(V,\phi)\in \IsoSht(F)$，按如下方式得到 $a\in (F^{\sep,\times}\ot\QQ)/\G_{F}$。存在有限扩张 $k_{n}/k$ 以及 $V$ 在 $F\ot_{k}k_{n}$ 上的下降 $V_{n}$，使得 $\phi$ 限制为 $\id_{F}\ot \Fr_{k_{n}/k}$-线性映射 $\phi_{n}: V_{n}\to V_{n}$。把 $a$ 定义为 $\l^{1/n}$ 的 Galois 轨道，其中 $\l\in F^{\sep,\times}$ 是 $F\ot_{k}k_{n}$-线性映射 $(\phi_{n})^{n}$ 的任一特征值。改变 $n$、下降 $V_{n}$ 与 $(\phi_{n})^{n}$ 的特征值 $\l$ 的选择，元素 $\l^{1/n}\in F^{\sep, \times}$ 只会乘上一个被 $\G_{F}$ 作用的单位根。由此得到 $F^{\sep,\times}\ot\QQ$ 中一个良定义的 $\G_{F}$-轨道。

\ssszh{张量结构}
为描述 $\IsoSht(F)$ 上的张量结构，考虑如下函数
\begin{equation*}
\chi: K_{0}(\IsoSht(F))\to \ZZ[F^{\sep, \times}\ot\QQ]^{\G_{F}}
\end{equation*}
它把对应于 $a\in F^{\sep,\times}\ot\QQ$ 的 Galois 轨道的单对象 $(V,\phi)$ 送到函数 $\chi(V,\phi): F^{\sep, \times}\ot\QQ\to \ZZ$：该函数在 $a$ 的 $\G_{F}$-轨道上支撑、在其上取常值，且总和等于 $\dim_{\br F} V$。

下面的命题描述 $\IsoSht(F)$ 中单对象如何彼此作张量积。
\begin{prop}\label{p:IsoSht tensor} 映射 $\chi$ 是单射环同态。
\end{prop}
这可以通过计算两个单对象 $V_{1},V_{2}\in \IsoSht(F)$ 的半单 $F$-代数 $\End(V_{1}\ot V_{2})\cong\End(V_{1})\ot_{ F}\End(V_{2})$ 的局部不变量来证明。

遗忘函子给出纤维函子 $\om: \IsoSht(F)\to \Vect(\br F)$，因此 $\IsoSht(F)$ 是 Tannaka 范畴。它在 $\br F$ 上的 Tannaka 群是可对角化群 $\DD_{\br F}$，特征群为 $\xch(\DD_{\br F})=F^{\sep,\times}\ot\QQ$，并带有 $\Gal(F^{\sep}/\br F)$ 的作用。

%%%%%%%%%%%%%%
%%%%%%%%%%%%%%


\section{亲戚与表亲}\label{s:sib}
本节介绍若干与 shtuka 密切相关的构造。我们可以把 Drinfeld 模看作 shtuka 的同卵双胞胎，把 Deligne–Lusztig 簇看作年幼的弟妹，而把有限域上的动机看作表亲。由于作者学识所限，我们略去了 Scholze \cite{S-Berkeley} 在过去十年发现的一个重要的混合特征表亲。

\subsection{Drinfeld 模}

我们回忆 Drinfeld 模的概念，以及 Drinfeld 模与某类特殊 shtuka 之间的词典。

\ssszh{Drinfeld 模}
设 $\infty\in |X|$，$A=\G(X-\{\infty\}, \cO_{X})$，这是 $k$ 上的 Dedekind 整环。

我们回忆秩 $n$ 的 Drinfeld $A$-模（按 \cite{Dr-Ell} 的称法，也叫椭圆模）在试验 $k$-概形 $S$ 上的概念。它是一个二元组 $(\cG, \io)$，其中 $\cG$ 是 $S$ 上的一维加性群概形（在 $S$ 上 Zariski 局部同构于 $\GG_{a,S}$），而 $\io: A\to \End_{S}(\cG)$ 是 $k$-代数同态，且满足：对每个几何点 $s=\Spec K \to S$ 以及任意 $K$-群概形同构 $\cG_{s}\cong\GG_{a,K}$，$\io$ 在 $s$-纤维上的限制
\begin{equation*}
\io_{s}: A\to \End_{K}(\cG_{s})\cong K\j{\t}
\end{equation*}
（其中 $\t\in \End_{K}(\cG_{s})\cong\End_{K}(\GG_{a,K})$ 是 Frobenius 态射 $x\mapsto x^{q}$）形如
\begin{equation*}
\io_{s}(a)=\sum_{i=0}^{-n[k_{\infty}:k]\ord_{\infty}(a)}c_{i}(a)\t^{i}, \quad \forall a\in A-\{0\}
\end{equation*}
且对 $a\ne0$，$c_{-n[k_{\infty}:k]\ord_{\infty}(a)}(a)\ne0$。

$A$ 在 $\cG$ 的李代数上的作用给出环同态 $A\to \G(S,\cO_{S})=\End_{S}(\Lie\cG)$，等价地给出映射 $\xi: S\to \Spec A=X-\{\infty\}$。这称为 Drinfeld $A$-模 $(\cG,\io)$ 的\emph{特征}。

设 $\Dr_{n}(S)$ 是 $S$ 上秩 $n$ 的 Drinfeld $A$-模构成的群胚。

\ssszh{Drinfeld 模簇}\label{sss:one leg}
为把 Drinfeld 模与 shtuka 联系起来，我们考虑例 \ref{ex:Dr} 中模问题的一个变体。

为简单起见，假设 $\infty$ 的剩余域为 $k$。设 $\Sht_{n}^{\Dr}$ 是其 $S$-点为 $(x, \cE, \{\cE(\frac{i}{n})\}_{i\in \ZZ}, \phi)$ 的模叠，其中
\begin{enumerate}
\item $x: S\to U:=X-\{\infty\}$。
\item $\cE$ 是 $X\times S$ 上秩 $n$ 的向量丛，且对任意几何点 $s\in S$，凝聚层 Euler 示性数 $\chi(X_{s}, \cE_{s})=0$（即 $\dim \upH^{0}(X_{s},\cE_{s})=\dim\upH^{1}(X_{s},\cE_{s})$）。
\item $\cE(-1)=\cE\ot\cO_{X}(-\infty)\subset \cE(-\frac{n-1}{n})\subset\cdots\subset \cE(-\frac{1}{n})\subset \cE(0)=\cE$ 是 $\cE|_{\{\infty\}\times S}$ 的一个完全旗，即对 $i=0,1,\cdots, n-1$，$\cE(-\frac{i}{n})/\cE(-\frac{i+1}{n})$ 在 $\{\infty\}\times S$ 上局部自由、秩为一。我们把 $\cE(\frac{i}{n})$ 的定义延拓到所有 $i\in \ZZ$，要求 $\cE(1+\frac{i}{n})=\cE(\frac{i}{n})\ot\cO_{X}(\infty)$。
\item $\phi: {}^{\t}\cE\to \cE(\frac{1}{n})$，其余核在 $\G(x)$ 上局部自由、秩为一；且对所有 $i=1,\cdots,n-1$，$\phi$ 把 ${}^{\t}(\cE(-\frac{i}{n}))$ 送到 $\cE(-\frac{i-1}{n})$。
\end{enumerate}

设 $I=\{1,2\}$，分解为单点集 $I_{1}=\{1\}$ 与 $I_{2}=\{2\}$。设 $\l_{1}=(0,\cdots,0,-1)$ 与 $\l_{2}=(1,0,\cdots ,0)$ 是 $\GL_{n}$ 的两个极小余权。在例 \ref{ex:Dr} 中我们定义了模叠 $\pi: \Sht^{(I_{1},I_{2}), (\l_{1},\l_{2})}_{\GL_{n}} \to X^{2}$。于是有映射
\begin{equation}\label{Dr to Sht}
\Sht_{n}^{\Dr}\to \Sht^{(I_{1},I_{2}), (\l_{1},\l_{2})}_{\GL_{n}}|_{U\times\{\infty\}}
\end{equation}
它把 $(x, \cE, \{\cE(\frac{i}{n})\}_{i\in \ZZ}, \phi)$ 送到两条leg $(x, \infty)$ 与向量丛
\begin{equation*}
\xymatrix{\cE_{0}=\cE & \cE_{1}:={}^{\t}(\cE(\frac{-1}{n}))\ar@{_{(}->}[l]_-{\phi}\ar@{^{(}->}[r] & \cE_{2}:={}^{\t}\cE}
\end{equation*}
可以证明 \eqref{Dr to Sht} 是闭嵌入；其像中的 $\ov k$-点由这样的 shtuka 组成：它在 $\infty$ 处的等晶是斜率为 $1/n$ 的超奇异等晶。

\ssszh{词典}\label{sss:dict}
据 \cite{Mum}，存在群胚的典范等价
\begin{equation*}
\Dr_{n}(S)\cong \Sht^{\Dr}_{n}(S).
\end{equation*}

我们描述两个方向的函子，更多细节见 \cite[Chapter 6]{Goss}。以下假设 $S$ 局部 Noether 环。

从 Drinfeld 模到 shtuka。设 $(\cG,\io)$ 是 $S$ 上秩 $n$ 的 Drinfeld 模。则 $M=\un\Hom_{S}(\cG,\GG_{a,S})$ 是 $S$ 上的拟凝聚层，且作为 $\cO_{S}\j{\t}=\un\End_{S}(\GG_{a,S})$-模局部自由、秩为一。此外，$M$ 通过 $\io$ 带有与 $\cO_{S}\j{\t}$-作用交换的 $A$-作用。我们将如下定义 $M$ 上一族由 $\cO_{S}$-子模构成的典范递增滤过 $M_{i}$。在 $S$ 上局部，取 $m_{0}\in M$ 作为自由 $\cO_{S}\j{\t}$-模 $M$ 的一组基。令 $M'(\frac{i}{n})$ 是 $M$ 中由 $m_{0}, \t m_{0},\cdots, \t^{i-1}m_{0}$ 生成的 $\cO_{S}$-子模。然后对滤过 $\{M'(\frac{i}{n})\}$ 作重整：定义 $M(\frac{i}{n})$ 为那些 $m\in M$，使得当 $\ord_{\infty}(a)$ 充分大（依赖于 $m$）时，$am\in M'(\frac{i}{n}+\ord_{\infty}(a))$。$m_{0}$ 的两种局部选择对 $i\gg0$ 给出相同的 $M'(\frac{i}{n})$（因为 $M$ 作为 $\cO_{S}\j{\t}$-模的任一其他基 $m_{1}$ 都形如 $m_{1}=(c_{0}+c_{1}\t+\cdots +c_{N}\t^{N})m_{0}$，其中 $c_{1}\in \cO_{S}^{\times}$，且 $c_{i}$ 对 $i\ge1$ 幂零\zy{原文此处作 $c_{1}\in\cO_{S}^{\times}$；按 $\cO_{S}\j{\t}$ 的单位群应为常数项 $c_{0}\in\cO_{S}^{\times}$，否则 $m_{1}$ 不可能是基。[依据：$\cO_{S}\j{\t}$ 中元素可逆当且仅当其常数项可逆]}）。这说明滤过 $\{M(\frac{i}{n})\}_{i\ge1}$ 与局部基 $m_{0}$ 的选择无关，从而定义了 $M$ 的一个由凝聚 $\cO_{S}$-子模构成的典范滤过。由构造可见 $\t M(\frac{i}{n})\subset M(\frac{i+1}{n})$。注意 $A$ 按在 $\infty$ 处的极点阶滤过：$A=\cup_{j\ge0}A_{j}$。令 $\Rees(A)=\op_{j\ge0}A_{j}$，则 $X=\Proj (\Rees(A))$。由构造，$A_{j}$ 把 $M(\frac{i}{n})$ 送到 $M(\frac{i}{n}+j)$。令 $N(\frac{i}{n})$ 为分次 $\Rees(A)\ot\cO_{S}$-模，其 $j$ 次分量为 $M(\frac{i}{n}+j)$。令 $\cE(\frac{i}{n})$ 为 $X\times S=\Proj_{S}(\Rees(A)\ot\cO_{S})$ 上与 $N(\frac{i}{n})$ 相伴的凝聚层。可以验证 $\cE(\frac{i}{n})$ 是 $X\times S$ 上秩 $n$ 的向量丛，且分次嵌入 $N(\frac{i}{n})\incl N(\frac{i+1}{n})$ 诱导单射 $\cE(\frac{i}{n})\incl \cE(\frac{i+1}{n})$，其余核在 $\{\infty\}\times S$ 上秩为一。最后，映射 $\phi: {}^{\t}\cE={}^{\t}\cE(0)\to \cE(\frac{1}{n})$ 来自 $\Fr_{S}$-线性映射 $\phi_{N}: N(0)\to N(\frac{1}{n})$，后者在 $j$ 次分量上是 $\t: M(j)\to M(\frac{1}{n}+j)$。

从 shtuka 到 Drinfeld 模。设 $(x, \cE, \{\cE(\frac{i}{n})\}_{i\in \ZZ}, \phi)\in \Sht^{\Dr}(S)$。令 $\pr_{S}: X\times S\to S$ 为投影。令 $M$ 为拟凝聚 $\cO_{S}$-模 $M=\varinjlim_{i}\pr_{S*}\cE(\frac{i}{n})$，它带有 $A$-作用。映射 $\phi$ 给出 $\Fr_{S}$-线性映射 $\phi_{M}: M\to M$。可以证明 $M$ 是 $\cO_{S}$ 上斜多项式环 $\cO_{S}\j{\t}$ 上的局部自由秩一模，其中 $\t$ 通过 $\phi_{M}$ 作用。由于 $\un\End_{S}(\GG_{a,S})=\cO_{S}\j{\t}$，秩一局部自由 $\cO_{S}\j{\t}$-模等同于 $\un\Aut(\GG_{a,S})\cong \cO_{S}\j{\t}^{\times}$ 的 Zariski 局部平凡主丛，也等同于 $\GG_{a}$ 在 $S$ 上的 Zariski 局部平凡形式 $\cG$。由于 $M$ 带有与 $\cO_{S}\j{\t}$-作用交换的 $A$-作用，$\cG$ 也带有 $A$-作用，这给出 Drinfeld 模的数据。

\ssszh{$t$-动机}
G. Anderson 的 $t$-动机理论 \cite{And} 是 $A=k[t]$ 情形的 Drinfeld 模的推广。一方面，若把 Drinfeld 模看作椭圆曲线的函数域类比，则 Anderson 的 $t$-动机是更一般动机的类比，包括高维 Abel 簇的类比。另一方面，\S\ref{sss:dict} 中的词典可推广到 $t$-动机：$t$-动机可以转成 $X=\PP^{1}$ 上的 shtuka，其中一条leg固定在 $\infty$，另一条leg移动。细节参见 Anderson 的原始论文 \cite{And} 与 D. Goss 的专著 \cite{Goss}。

\subsection{（仿射）Deligne–Lusztig 簇}
设 $\cB$ 是 $G$ 的旗簇，即 $G$ 的 Borel 子群的模空间。$\cB\times\cB$ 上（对角作用下的）$G$-轨道集合构成 Coxeter 群 $(W,S)$，它同构于 $G$ 的极大分裂环面所附的 Weyl 群。对 $(B,B')\in \cB^{2}$，若 $(B,B')$ 落在由 $w\in W$ 标号的 $G$-轨道中，记 $\pos(B,B')=w$；若 $(B,B')$ 落在该轨道的闭包中，记 $\pos(B,B')\le w$。

把 $G$-主丛换成 Borel 子群，可以定义 $\Hk^{(I_{1},\cdots, I_{r})}_{G}$ 的类比（至少在每个 $I_{j}$ 为单点集时）。设 $\un w=(w_{j})_{1\le j\le r}\in W^{r}$。定义 $Y^{\un w}$ 为 $\cB^{r+1}$ 中由满足 $\pos(B_{j-1},B_{j})=w_{j}$（$1\le j\le r$）的 Borel 子群 $(B_{0},\cdots, B_{r})$ 构成的子概形。类似地，定义 $Y^{\le \un w}\subset \cB^{r+1}$ 为满足 $\pos(B_{j-1},B_{j})\le w_{j}$（$1\le j\le r$）的闭子概形。

于是可以定义 $\Sht^{(I_{1},\cdots, I_{r})}_{G}$ 的类比。用如下 Cartesian 图定义 $X^{\un w}$ 与 $X^{\le \un w}$：
\begin{equation*}
\xymatrix{ X^{\un w}\ar[r]\ar[d] & Y^{\un w}\ar[d]^{(p_{0}, p_{r})} & X^{\le\un w}\ar[r]\ar[d]  & Y^{\le\un w}\ar[d]^{(p_{0}, p_{r})} \\
\cB\ar[r]^-{(\id,\Fr)} & \cB^{2} & \cB\ar[r]^-{(\id,\Fr)} & \cB^{2}
}
\end{equation*}
这里 $p_{j}$ 是第 $j$ 个投影 $\cB^{r+1}\to \cB$。注意当 $r=1$ 时，$X^{w}$ 就是 Deligne–Lusztig 簇 \cite[Definition 1.4]{DL}，它分类满足 $\pos(B,\Fr(B))=w$ 的 Borel 子群 $B$。Deligne–Lusztig 簇在 Lie 型有限群的不可约表示的构造与分类中起关键作用。

类似地，若把 $G$ 换成环路群 $LG$，$B$ 换成 Iwahori 子群 $\bI\subset LG$，$\cB$ 换成仿射旗簇 $\Fl_{G}=LG/\bI$，$W$ 换成扩张仿射 Weyl 群 $\tilW$，则可类似地定义 $\cY^{\un w}\subset \cY^{\le\un w}\subset (\Fl_{G})^{r+1}$，其中 $\un w\in (\tilW)^{r}$。再把它们与 $\Fl_{G}$ 上 Frobenius 的图相交，便得到\emph{仿射 Deligne–Lusztig 簇} $\cX^{\un w}$ 与 $\cX^{\le \un  w}$ 的特殊情形。

更一般的仿射 Deligne–Lusztig 簇由 Rapoport \cite[\S4]{Ra} 定义。为定义它们，固定 $b\in LG$，并把上一段中的 Frobenius 映射 $\Fl_{G}\to \Fl_{G}$ 换成 $b\c\Fr$，其中 $b$ 通过左平移作用。所得 $\Fl_{G}$ 的子 ind-概形记作 $\cX^{\un w}(b)$ 与 $\cX^{\le\un w}(b)$，其同构类只依赖 $b$ 的 $\Fr$-共轭类。

\subsection{有限域上的动机}
现在我们指出 shtuka 与有限域上动机之间的相似之处。主要参考文献是 Milne \cite{Milne} 与 Langlands–Rapoport \cite{LR}。

固定素数 $p$。对 $\FF_{p}$ 的每个有限扩张 $\FF_{q}$，存在 $\FF_{q}$ 上（纯）动机的 $\QQ$-线性张量范畴 $\Mot(\FF_{q})$，其 hom 空间用乘积上的圈模数值等价来定义。类似地有 $\ov\FF_{p}$ 上（纯）动机的 $\QQ$-线性张量范畴 $\Mot(\ov\FF_{p})$，它是 $\Mot(\FF_{q})$ 当 $q$ 遍历 $\FF_{p}$ 的有限扩张时的余极限。已知这些范畴半单 \cite{Jan}。

$\Mot(\ov\FF_{p})$ 作为 Abel 范畴有如下描述，它推广了有限域上 Abel 簇的 Honda–Tate 理论。回忆：权为 $n\in \ZZ$ 的 Weil $p$-数是代数数 $\pi$，它对所有素数 $\ell\ne p$ 都是 $\ell$-进单位，且 $\pi$ 的所有 阿基米德 范数都等于 $p^{n/2}$。设 $W_{p}\subset \ov \QQ$ 为各种整数权的 Weil $p$-数之集，$\mu_{\infty}\subset \ov\QQ^{\times}$ 为单位根群；后者被 Galois 群 $\G_{\QQ}=\Gal(\ov\QQ/\QQ)$ 置换。给定 $\pi\in W_{p}/\mu_{\infty}$，定义数域 $L_{\pi}=\cap_{n\in\NN}\QQ[\pi^{n}]$。

下面的定理应与定理 \ref{th:isoSht} 对照。
\begin{thm}[参见 {\cite[Theorem 2.18, 2.19]{Milne}}]\label{th:mot}
\begin{enumerate}
\item $\Mot(\ov\FF_{p})$ 中单对象的同构类与 $W_{p}/\mu_{\infty}$ 上的 $\G_{\QQ}$-轨道自然一一对应。
\item 若单动机 $X\in \Mot(\ov\FF_{p})$ 对应于权为 $w$ 的 $\pi\in W_{p}/\mu_{\infty}$ 的 Galois 轨道，则
\begin{enumerate}
\item $\End(X)$ 是数域 $L_{\pi}$ 上的中心可除代数，其在位置 $v\in |L_{\pi}|$ 处的不变量为
\begin{equation*}
\inv_{v}(\End(X))=\begin{cases}  -\ord_{v}(\pi)[k_{v}:\FF_{p}], & v\mid p \\ \frac{1}{2}, &  v\mbox{ 实，$w$ 为奇数} ; \\ 0, & \text{其他情形. } \end{cases}
\end{equation*}

\item $X$ 的秩为 $[L_{\pi}:\QQ]d_{\pi}$，其中 $d_{\pi}$ 是有理数集合 $\{\ord_{v}(\pi)[k_{v}:\FF_{p}]\}_{v|p}$ 的最小公分母。
\end{enumerate}
\end{enumerate}
\end{thm}

为描述 $\Mot(\ov\FF_{p})$ 上的张量结构，考虑如下线性函数
\begin{equation*}
\chi: K_{0}(\Mot(\ov\FF_{p}))\to \ZZ[W_{p}/\mu_{\infty}]^{\G_{\QQ}}
\end{equation*}
它把对应于 $\pi\in W_{p}/\mu_{\infty}$ 的 Galois 轨道的单动机 $X$ 送到函数 $\chi(X): W_{p}/\mu_{\infty}\to \ZZ$：该函数在 $\pi$ 的 $\G_{\QQ}$-轨道上支撑，并取常值 $d_{\pi}$。下述命题与命题 \ref{p:IsoSht tensor} 类似：

\begin{prop}[参见 {\cite[Corollary 2.20]{Milne}}] 映射 $\chi$ 是单射环同态。
\end{prop}

\ssszh{纤维函子}
对素数 $\ell\ne p$，取动机的 $\ell$-进 étale 上同调给出纤维函子
\begin{equation*}
\om_{\ell}: \Mot(\ov\FF_{p})\ot\QQ_{\ell}\to \Vect(\Ql).
\end{equation*}
对 $\ell=p$，取 $\ov\FF_{p}$ 上动机的晶体上同调给出张量函子
\begin{equation*}
\om_{p}: \Mot(\ov\FF_{p})\ot\QQ_{p}\to \IsoCrys(\QQ_{p}).
\end{equation*}
假定 Tate 猜想，$\Mot(\ov\FF_{p})\ot\RR$ 也容许到 $\RR$-向量空间的纤维函子。这些纤维函子是 \eqref{Sht fib functor} 中定义的 $\IsoSht(F)$ 上函子 $\om_{x}$ 的类比。

\ssszh{$\QQ$ 上的 iso-shtuka？}
对照定理 \ref{th:mot} 与定理 \ref{th:isoSht}，这提示应把 $\ov\FF_{p}$ 上的动机看作数域 $\QQ$ 上的 iso-shtuka，其leg在 $p$ 与 $\infty$。这一类比在 $\ov\FF_{p}$ 上椭圆曲线的动机与秩 $2$ 的 Drinfeld 模对应的 shtuka 之间尤其贴切。

在 \cite{Kot} 中，Kottwitz 把半单 $\QQ$-线性 Abel 范畴定义为 $\QQ$ 上以投射环面为带的 gerbe 的表示。猜想它对所有素数 $p$ 都把 $\Mot(\ov\FF_{p})$ 作为全子范畴包含在内。在 \cite{SchICM} 中，Scholze 建议把 Kottwitz 的范畴看作 $\IsoSht(\QQ)$（带有任意有限条leg）。

\begin{ques} 若 $\ov\FF_{p}$ 上的动机给出 $\QQ$ 上leg在 $p$ 与 $\infty$ 的 iso-shtuka，那么什么样的几何对象给出 $\QQ$ 上leg在若干素数的 iso-shtuka？
\end{ques}

%%%%%%%%%%%%%%%%%%%%%
\section{模叠的几何}\label{s:geom}
本节总结 $G$-shtuka 模叠的基本几何性质与结构。正是它们造就了这些模叠上同调上的对称性，下一节将研究这些对称性。

\subsection{可表示性}
关于 shtuka 模叠的几何，最基本的结果如下。
\begin{thm}[Varshavsky {\cite[Proposition 2.16(a)]{Var}}] 设 $I=I_{1}\coprod\cdots\coprod I_{r}$ 是有限集的分解，且 $\un\l=(\l_{i})_{i\in I}\in (\xcoch(T)^{+})^{I}$。则叠 $\Sht^{(I_{1},\cdots, I_{r}), \le \un\l}_{G,\bK}$ 是 $k$ 上局部有限型的 Deligne–Mumford 叠。
\end{thm}

众所周知，$\Bun_{G,\bK}$ 是 $k$ 上局部有限型的光滑 Artin 叠。此外，映射 $p_{i}: \Hk_{G,\bK}^{(I_{1},\cdots, I_{r}), \le \un\l}\to \Bun_{G,\bK}$ 有限型。由定义 $\Sht^{(I_{1},\cdots, I_{r}), \le\un\l}_{G,\bK}$ 的图 \eqref{defn Sht} 可见，它是 $k$ 上局部有限型的 Artin 叠。

\ssszh{Harder–Narasimhan 截断}\label{sss:HN}
可以如下用 $k$ 上有限型的开子叠覆盖 $\Sht^{(I_{1},\cdots, I_{r}), \le \un\l}_{G,\bK}$。

先考虑无水平结构的情形。设 $\xcoch(T)^{+}_{\QQ}$ 是 $\xcoch(T)_{\QQ}$ 中的支配锥。设 $\HN\subset \xcoch(T)^{+}_{\QQ}$ 为满足如下整性条件的 $\mu\in \xcoch(T)_{\QQ}^{+}$ 之集：若 $L_{\mu}$ 是以 $\{\a|\j{\a,\mu}=0\}$ 为根的标准 Levi 子群，则 $\mu\in \xcoch(L_{\mu,\ab})_{\QQ}=\xcoch(ZL_{\mu})_{\QQ}\subset\xcoch(T)_{\QQ}$（这里 $L_{\mu,\ab}$ 是 $L_{\mu}$ 的交换化）。每个几何点 $\cE\in \Bun_{G}$ 都有关于标准抛物子群 $P\subset G$ 的典范 $P$-约化 $\cE_{P}$（推广向量丛的 Harder–Narasimhan 滤过），使得：设 $L$ 是 $P$ 的 Levi 商，则诱导的 $L_{\ab}$-主丛 $\cE_{L_{\ab}}=\cE_{P}/(U_{P}[L,L])$ 的次数 $\deg(\cE_{L_{\ab}})\in \xcoch(L_{\ab})$ 满足：对 $U_{P}$ 中所有根 $\a$，$\j{\a, \deg(\cE_{L_{\ab}})}>0$。因此每个几何点 $\cE\in \Bun_{G}$ 都有良定义的 Harder–Narasimhan 点 $\deg(\cE_{L_{\ab}})\in \HN$。

考虑 $\HN$ 上的偏序：$\mu\le \mu'$ 当且仅当 $\mu'- \mu $ 是单余根的 $\QQ_{\ge0}$-线性组合。存在开子叠 ${}^{\le\mu}\Bun_{G}\subset \Bun_{G}$，其几何点的 HN 点 $\le \mu$。已知 ${}^{\le\mu}\Bun_{G}$ 在 $k$ 上有限型。

容易看出，记录 $\cE_{0}$ 的遗忘映射 $p_{0}: \Sht^{(I_{1},\cdots, I_{r}), \le \un\l}_{G}\to \Bun_{G}$ 有限型。现对 $\mu\in \HN$，令 ${}^{\le \mu}\Sht^{(I_{1},\cdots, I_{r}), \le \un\l}_{G}=p_{0}^{-1}({}^{\le \mu}\Bun_{G})$，则 ${}^{\le \mu}\Sht^{(I_{1},\cdots, I_{r}), \le \un\l}_{G}$ 在 ${}^{\le \mu}\Bun_{G}$ 上有限型，从而在 $k$ 上有限型。当 $\mu\in \HN$ 变动时，$\{{}^{\le \mu}\Sht^{(I_{1},\cdots, I_{r}),\le\un\l}_{G}\}_{\mu\in\HN}$ 覆盖 $\Sht^{(I_{1},\cdots, I_{r}),\le\un\l}_{G}$。

对一般的水平结构 $\bK=(\bK_{x})_{x\in \Sig}$，可取 $n_{x}\in\NN$，使得 $\bK_{x}$ 包含主同余子群 $L^{(n_{x})}_{x}G$。令 $D=\sum_{x\in \Sig}n_{x}x$ 与 $\bK_{D}=(L^{(n_{x})}_{x}G)_{x\in \Sig}$，则 $\Sht^{(I_{1},\cdots, I_{r})}_{G,\bK_{D}}\to \Sht^{(I_{1},\cdots, I_{r})}_{G,\bK}$ 是有限 étale 映射（见 \S\ref{ss:change level}），因此只需用 $k$ 上有限型的开子叠覆盖 $\Sht^{(I_{1},\cdots, I_{r})}_{G,\bK_{D}}$。有遗忘映射 $\Sht^{(I_{1},\cdots, I_{r}),\le\un\l}_{G,\bK_{D}}\xr{p_{0}}\Bun_{G,D}\to \Bun_{G}$。对 $\mu\in \HN$，取 ${}^{\le \mu}\Sht^{(I_{1},\cdots, I_{r}),\le\un\l}_{G,\bK_{D}}$ 为 ${}^{\le\mu}\Bun_{G}$ 的原像，则 ${}^{\le \mu}\Sht^{(I_{1},\cdots, I_{r}),\le\un\l}_{G,\bK_{D}}$ 在 $\Sht^{(I_{1},\cdots, I_{r}),\le\un\l}_{G,\bK_{D}}$ 中是开的，且在 $k$ 上有限型。当 $\mu\in \HN$ 变动时，${}^{\le \mu}\Sht^{(I_{1},\cdots, I_{r}),\le\un\l}_{G,\bK_{D}}$ 覆盖 $\Sht^{(I_{1},\cdots, I_{r}),\le\un\l}_{G,\bK_{D}}$。

\medskip\noindent Varshavsky \cite[Prop.2.16(1)]{Var} 实际上证明了更精确的结果：每个有限型开子叠 ${}^{\le\mu}\Sht^{(I_{1},\cdots, I_{r}), \le\un\l}_{G,\bK}$ 都是一个拟射影概形被一个有限（常）群作用后的商。论证思路是：通过过渡到更深的水平结构 $\bK'_{x}$（仍在 $x\in \Sig$ 处），其在 $\bK_{x}$ 中正规，可使 ${}^{\le\mu}\Bun_{G,\bK'}$ 成为拟射影概形。于是由图 \eqref{defn Sht} 清楚可见，${}^{\le\mu}\Sht^{(I_{1},\cdots, I_{r}), \le\un\l}_{G,\bK'}$ 是拟射影概形。由 \S\ref{ss:change level} 的讨论，映射 ${}^{\le\mu}\Sht^{(I_{1},\cdots, I_{r}), \le\un\l}_{G,\bK'}\to {}^{\le\mu}\Sht^{(I_{1},\cdots, I_{r}), \le\un\l}_{G,\bK}$ 是有限群 $\prod_{s\in \Sig}K_{x}/K'_{x}$ 的主齐性空间，其中 $K_{x}=\bK_{x}(k_{x})$，$K'_{x}=\bK'_{x}(k_{x})$。

\subsection{局部模型}
人们希望了解 shtuka 模叠更精细的几何信息，例如其光滑性与维数。这些可以借助仿射 Grassmann 簇中的仿射 Schubert 簇给出的局部模型来理解。

设 $I=I_{1}\coprod\cdots\coprod I_{r}$ 是分成\emph{单点集}的分解。回忆相对位置映射 \eqref{ev disj}：
\begin{equation}\label{ev Hk singleton}
\ev^{(I_{1},\cdots, I_{r})}_{\Hk}: \Hk^{(I_{1},\cdots, I_{r})}_{G,\bK}\to (\QGr/\Aut(D))^{I}.
\end{equation}
注意由于每个 $I_{j}$ 都是单点集，该映射在整个 $U^{I}$ 上有定义。考虑复合
\begin{equation}\label{ev Sht singleton}
\ev^{(I_{1},\cdots, I_{r})}_{\Sht}: \Sht^{(I_{1},\cdots, I_{r})}_{G,\bK}\xr{} \Hk^{(I_{1},\cdots, I_{r})}_{G,\bK}\xr{\ev^{(I_{1},\cdots, I_{r})}_{\Hk}} (\QGr/\Aut(D))^{I}.
\end{equation}

\begin{prop}\label{p:fm sm} 设 $I=I_{1}\coprod\cdots\coprod I_{r}$ 是分成单点集的分解。则映射 \eqref{ev Hk singleton} 与 \eqref{ev Sht singleton} 形式光滑。
\end{prop}
\begin{proof}[证明梗概]

先证 \eqref{ev Sht singleton} 形式光滑。

回忆 $D=\Spec k\tl{t}$。设 $\wh U\to U$ 是典范 $\Aut(D)$-主丛 $(x,t)$，其中 $t$ 是 $x\in U$ 处的局部参数。则我们有映射
\begin{equation*}
\wh\ev: \wh\Sht^{(I_{1},\cdots, I_{r})}_{G,\bK}:=\wh U^{I}\times_{U^{I}}\Sht^{(I_{1},\cdots, I_{r})}_{G,\bK}\to  \wh U^{I}\times \QGr^{I}
\end{equation*}
使得 $\ev^{(I_{1},\cdots, I_{r})}$ 由投影到 $\QGr^{I}$ 并在两边对 $\Aut(D)^{I}$ 取商得到。因此只需证明 $\wh\ev$ 形式光滑。

为记号简洁，只处理无水平结构的情形，一般情形类似。可把 $I$ 与 $\{1,\cdots, r\}$ 等同，使 $I_{j}=\{j\}$。设 $R'$ 是 $k$-代数，$J\subset R'$ 是平方为零的理想，$R=R'/J$。假设给定 $R'$-点 $(x'_{j}, t'_{j})\in \wh X(R')$（$1\le j\le r$）、$R'$-点 $\z'_{j}=(\cF'_{j}\dxr{} \cG'_{j})\in \QGr(R')$（$1\le j\le r$，其中 $(\cF'_{j}, \cG'_{j})$ 是 $D_{R'}$ 上的一对 $G$-主丛，并在 $D^{\times}_{R'}$ 上带同构），以及 $\wh \Sht^{(I_{1},\cdots, I_{r})}_{G}$ 的 $R$-点 $\xi=((x_{j},t_{j})_{1\le j\le r}, (\cE_{j})_{0\le j\le r}, (\a_{j})_{1\le j\le r}, \io)$，使得 $\wh\ev(\xi)=((x'_{j}, t'_{j},\z'_{j})|_{\Spec R})_{1\le j\le r}$。特别地，$(x_{j},t_{j})=(x'_{j}, t'_{j})|_{\Spec R}$。我们希望把 $\xi$ 提升为 $\wh \Sht^{(I_{1},\cdots, I_{r})}_{G}$ 的 $R'$-点 $\xi'$，使得 $\wh\ev(\xi')=(x'_{j}, t'_{j}, \z'_{j})_{1\le j\le r}$。这意味着需要把 $\cE_{j}$ 提升为 $X_{R'}$ 上的 $G$-主丛 $\cE'_{j}$（$0\le j\le r$），把同构 $\a_{j}$ 提升为 $\a'_{j}: \cE'_{j-1}|_{X_{R'}-\G(x'_{j})}\isom \cE'_{j}|_{X_{R'}-\G(x'_{j})}$（$1\le j\le r$），并把同构 $\io: \cE_{r}\cong(\id_{X}\times \Fr_{R})^{*}\cE_{0}$ 提升为同构 $\io': \cE'_{r}\cong(\id_{X}\times \Fr_{R'})^{*}\cE'_{0}$。

首先，$\cE'_{r}$ 已由 $\cE_{0}$ 如下确定。由于 $J^{2}=0$ 是平方零加厚，$\Fr_{R'}$ 分解为 $R'\surj R\xr{\ov\Fr_{R'}} R'$。设 $\ph: \Spec R'\to \Spec R$ 是由 $\ov\Fr_{R'}$ 给出的映射。若 $\cE'_{0}$ 是 $\cE_{0}$ 到 $X_{R'}$ 的任一扩张，则 $(\id_{X}\times \Fr_{R'})^{*}\cE'_{0}=(\id_{X}\times \ph)^{*}\cE_{0}$。因此，为得到所需同构 $\io': \cE'_{r}\cong {}^{\t}\cE'_{0}$，必须把 $\cE'_{r}$ 定义为 $(\id_{X}\times \ph)^{*}\cE_{0}$。

接着把沿 $\G(x_{j})$ 的修正 $\a_{j}: \cE_{j-1}\dxr\cE_{j}$ 延拓为沿 $\G(x'_{j})$ 的修正 $\a'_{j}: \cE'_{j-1}\dxr\cE'_{j}$。我们从 $j=r$ 开始归纳进行。每一步给定 $X_{R'}$ 上的 $\cE'_{j}$，试图构造 $\cE'_{j-1}$ 与 $\a'_{j}$。用 $t'_{j}$ 把 $D_{x'_{j}}$ 与 $D_{R'}$ 等同，从而也把 $D_{x_{j}}$ 与 $D_{R}$ 等同。同构 $\wh\ev(\xi)=(x_{j}, t_{j},\z'_{j}|_{\Spec R})_{1\le j\le r}$ 给出图之间的同构
\begin{equation*}
\xymatrix{\cE_{j-1}|_{D_{x_{j}}}\ar[d]^{\e_{j-1}}_{\wr}\ar@{-->}[r]^{\a_{j}} & \cE_{j}|_{D_{x_{j}}}\ar[d]^{\e_{j}}_{\wr}\\
\cF'_{j}|_{D_{R}}\ar@{-->}[r]^{\z_{j}} & \cG'_{j}|_{D_{R}}
}
\end{equation*}
先把同构 $\e_{j}$ 提升为 $\e'_{j}: \cE'_{j}|_{D_{x'_{j}}}\cong \cG'_{j}$（之所以可行，是因为 $L^{+}G$ 形式光滑）。然后把 $\cE'_{j-1}$ 定义为 $\cE'_{j}|_{X_{R'}-\G(x'_{j})}$ 与 $\cF'_{j}$ 沿同构 $\cE'_{j}|_{D^{\times}_{x'_{j}}}\xr{\e'_{j}} \cG'_{j}|_{D^{\times}_{R'}}\xr{\z_{j}^{-1}}\cF'_{j}|_{D^{\times}_{R'}}$ 的粘合。于是有典范同构 $\cE'_{j-1}|_{X_{R}}\cong \cE_{j-1}$，因为后者由 $\cE_{j}$ 与 $\cF'_{j}|_{D_{R}}$ 沿 $\cE_{j}|_{D^{\times}_{x_{j}}}\xr{\e_{j}}\cG'_{j}|_{D^{\times}_{R}}\xr{\z_{j}^{-1}}\cF'_{j}|_{D^{\times}_{R}}$ 粘合而成。这就完成了 $\cE'_{j}$ 的归纳构造。因此 \eqref{ev Sht singleton} 形式光滑。

为证 \eqref{ev Hk singleton} 形式光滑，论证类似：只需从 $\cE_{r}$ 到 $X_{R'}$ 的任一提升 $\cE'_{r}$ 出发，其存在性来自 $\Bun_{G}$ 的光滑性。
\end{proof}

过渡到有界版本，命题 \ref{p:fm sm} 中证明的形式光滑性蕴含有界版本相对位置映射的光滑性。更强地，仿射 Grassmann 簇在如下意义下是模叠的 étale 局部模型。

\begin{prop}[{\cite[Theorem 2.20]{Var}}，{\cite[Theorem 2.11]{VLaff}}]\label{p:local model} 设 $I=I_{1}\coprod\cdots\coprod I_{r}$ 是有限集 $I$ 的分解，$\un\l\in (\xcoch(T)^{+})^{I}$。

\begin{enumerate}
\item 叠 $\Sht^{(I_{1},\cdots, I_{r}), \le \un\l}_{G,\bK}$ 在 $U^{I}$ 上 étale 局部同构于 $\Gr^{(I_{1},\cdots, I_{r}), \le\un\l}_{G}$，后者又 étale 局部同构于乘积 $\prod_{j=1}^{r}\Gr^{I_{j},\le\un\l_{j}}_{G}$（其中 $\un\l_{j}=\{\l_{i}\}_{i\in I_{j}}$）。
\item 特别地，若 $\un\l$ 是 $G$-容许的，则
\begin{equation*}
\dim \Sht^{(I_{1},\cdots, I_{r}),\le\un\l}_{G,\bK}=\sum_{i\in I}(\j{2\r, \l_{i}}+1).
\end{equation*}
\item 设对每个 $j$，$I_{j}=I_{j,1}\coprod\cdots\coprod I_{j,s_{j}}$ 是分解，并令 $I=I'_{1}\coprod\cdots\coprod I'_{r'}$ 是 $I=I_{1}\coprod\cdots\coprod I_{r}$ 由此得到的加细。则态射 $\pi^{(I'_{1},\cdots, I'_{r'}),\le\un\l}_{(I_{1},\cdots, I_{r})}: \Sht^{(I'_{1},\cdots, I'_{r'}), \le \un\l}_{G,\bK}\to \Sht^{(I_{1},\cdots, I_{r}), \le \un\l}_{G,\bK}$ 在目标上 étale 局部同构于自然映射之积 $\prod_{j=1}^{r}\Gr^{(I_{j,1},\cdots, I_{j,s_{j}}),\le\un\l_{j}}_{G}\to \prod_{j=1}^{r}\Gr^{I_{j},\le\un\l_{j}}_{G}$。
\end{enumerate}
\end{prop}
证明最终归结为一个切向计算，并用到了 Frobenius 映射的微分为零。

若 $\l_{i}$ 都是极小的，则 $\Gr^{\le\l_{i}}_{G}$ 光滑，从而 $\Gr^{(I_{1},\cdots, I_{r}), \le\un\l}_{G}$ 在 $U^{I}$ 上光滑。此时命题 \ref{p:local model} 给出：
\begin{cor}\label{c:Sht sm} 设 $I=I_{1}\coprod\cdots\coprod I_{r}$ 是分成单点集的分解，且对所有 $i\in I$，$\l_{i}\in \xcoch(T)^{+}$ 是极小支配余权。则 $\Sht^{(I_{1},\cdots, I_{r}),\le\un\l}_{G,\bK}$ 在 $U^{I}$ 上光滑，纯相对维数为 $\sum_{i\in I}\j{2\r, \l_{i}}$。
\end{cor}

% ch03  §4.3–§4.7 / §5.1（原文 1296–1647 行）

\subsection{分解结构}\label{ss:geom fact}

对一簇 $U^{I}$ 上的几何对象 $A_{I}$（叠或层），每个有限集 $I$ 对应其中一个。所谓 $\{A_{I}\}$ 上的分解结构，通常包含两个方面。其一是 $A_{I}$ 限制到 $U^{I}_{\disj}$ 时的行为：例如当 $A_{I}$ 是 $U^{I}$ 上的叠时，可要求 $A_{I}|_{U^{I}_{\disj}}$ 配有与 $(\prod_{i\in I}A_{\{i\}})|_{U^{I}_{\disj}}$ 的同构。其二是 $A_{I}$ 在部分对角上的行为：例如可要求 $A_{I}|_{\D(U)}\cong A_{\{1\}}$。这类分解结构的典型例子由 $X^{I}$ 上的 Beilinson--Drinfeld Grassmannian $\Gr^{I}$ 承载。

在 shtuka 的情形，上述第一个方面只在弱的意义下成立，其内容本质上就是命题 \ref{p:fm sm} 与命题 \ref{p:local model}。下面把第二个方面写清楚。

对任意叠 $\frX$ 与有限集之间的任意映射 $\th: I\to J$，有诱导映射
\begin{equation*}
\D_{\th, \frX}: \frX^{J}\to \frX^{I}
\end{equation*}
它把 $(x_{j})_{j\in J}$ 送到 $(x_{\th(i)})_{i\in I}$。当 $\th$ 是满射时，$\D_{\th,\frX}$ 是 $\frX^{I}$ 中的部分对角。

Hecke 叠的如下性质由定义直接看出。若 $\th: I\surj J$ 是有限集之间的满射，则存在典范同构
\begin{equation}\label{Hk diagonal}
\Hk^{I}_{G,\bK}|_{\D_{\th,U}(U^{J})}\cong \Hk^{J}_{G,\bK}.
\end{equation}
此外，该同构与满射的复合 $I\surj J\surj K$ 相容。

特别地，取 $J$ 为单点集，便得到
\begin{equation*}
\Hk^{I}_{G,\bK}|_{\D(U)}\cong \Hk^{\{1\}}_{G,\bK}.
\end{equation*}

同样的性质对 $\Sht^{I}_{G,\bK}$ 也成立。若 $\th: I\surj J$ 是有限集之间的满射，则存在典范同构
\begin{equation*}
\Sht^{I}_{G,\bK}|_{\D_{\th,U}(U^{J})}\cong \Sht^{J}_{G,\bK}.
\end{equation*}
此外，该同构与满射的复合 $I\surj J\surj K$ 相容。

特别地，取 $J$ 为单点集，便得到
\begin{equation*}
\Sht^{I}_{G,\bK}|_{\D(U)}\cong \Sht^{\{1\}}_{G,\bK}.
\end{equation*}

\subsection{水平结构的改变}\label{ss:change level}

\ssszh{水平提升}

对每个 $x\in \Sig$，设 $\bK^{\#}_{x}\lhd \bK_{x}$ 是正规子群，其商 $\bL_{x}=\bK_{x}/\bK^{\#}_{x}$ 有限。令 $\bK^{\#}=(\bK^{\#}_{x})_{x\in \Sig}$。则自然映射
\begin{equation*}
\Sht^{(I_{1},\cdots, I_{r})}_{G,\bK^{\#}} \to \Sht^{(I_{1},\cdots, I_{r})}_{G,\bK}
\end{equation*}
是 $L$-主丛，其中 $L:=\prod_{x\in \Sig}\bL_{x}(k_{x})$ 为有限群。

对这样的正规子群 $\bK^{\#}$ 取投影极限，便得到 $U^{I}$ 上的叠
\begin{equation*}
\Sht^{(I_{1},\cdots, I_{r})}_{G,\infty\Sig} :=\varprojlim_{\bK^{\#}}\Sht^{(I_{1},\cdots, I_{r})}_{G,\bK^{\#}}.
\end{equation*}
其投影
\begin{equation*}
\Sht^{(I_{1},\cdots, I_{r})}_{G,\infty\Sig}\to \Sht^{(I_{1},\cdots, I_{r})}_{G,\bK}
\end{equation*}
是 pro-\'etale 映射，且是以 $K=\prod_{x\in \Sig}K_{x}$ 为结构群的主丛，其中 $K_{x}=\bK_{x}(k_{x})$。特别地有：

\begin{lem}\label{l:change level Sig}
对每个 $x\in \Sig$，设 $\bK'_{x}\subset \bK_{x}$ 是水平群，并令 $\bK'=(\bK'_{x})_{x\in \Sig}$。则自然映射
\begin{equation*}
\Sht^{(I_{1},\cdots, I_{r})}_{G,\bK'} \to \Sht^{(I_{1},\cdots, I_{r})}_{G,\bK}
\end{equation*}
是有限 \'etale 的，其几何纤维与陪集 $\prod_{x\in \Sig}K_{x}/K'_{x}$ 一一对应（其中 $K'_{x}=\bK'_{x}(k_{x})$）。
\end{lem}

\ssszh{在 $x$ 处的满水平结构}\label{sss:full x}

若 $x\notin \Sig$，则可定义 $(U-\{x\})^{I}$ 上的叠
\begin{equation*}
\Sht^{(I_{1},\cdots, I_{r})}_{G,\bK,\infty x}=\varprojlim_{\bK^{\da}_{x}}\Sht^{(I_{1},\cdots, I_{r})}_{G,\bK^{\da}}
\end{equation*}
其中 $\bK^{\da}$ 表示沿 $\{x\}\cup \Sig$ 的水平结构，即在 $\Sig$ 上取 $\bK$、在 $x$ 处取 $\bK^{\da}_{x}$。这是 $\Sht^{(I_{1},\cdots, I_{r})}_{G, \bK}|_{(U-\{x\})^{I}}$ 上的 $G(\cO_{x})$-主丛。

我们断言，$G(\cO_{x})$ 在 $\Sht^{(I_{1},\cdots, I_{r})}_{G, \bK,\infty x}$ 上的作用可典范地延拓为 $G(F_{x})$ 的作用。事实上，设
\begin{equation*}
\xi=((x_{i})_{i\in I}, (\cE_{j})_{0\le j\le r}, (\a_{j})_{1\le j\le r}, \io, (\ka_{j})_{0\le j\le r})
\end{equation*}
是 $\Sht^{(I_{1},\cdots, I_{r})}_{G,\bK,\infty x}$ 的一个 $R$-点，其中 $\ka_{j}: G\times D_{x, R}\isom \cE_{j}|_{D_{x,R}}$ 是对应于 $x$ 处满水平结构的平凡化。这里 $D_{x,R}=\Spec(\cO_{x}\wh\ot R)$，$D_{x,R}^{\times}=\Spec(F_{x}\wh\ot R)$。设 $g\in G(F_{x})$，定义 $\xi\cdot g$ 为多元组
\begin{equation*}
((x_{i})_{i\in I}, (\cE^{g}_{j})_{0\le j\le r}, (\a'_{j})_{1\le j\le r}, \io', (\ka'_{j})_{0\le j\le r}).
\end{equation*}
其中 $\cE^{g}_{j}$ 是把 $\cE_{j}|_{(X-\{x\})_{R}}$ 与 $D_{x,R}$ 上的平凡 $G$-主丛用同构
\begin{equation*}
G\times D^{\times}_{x, R}\xr{g} G\times D^{\times}_{x, R} \xr{\ka_{j}} \cE_{j}|_{D^{\times}_{x,R}}.
\end{equation*}
粘合所得的 $G$-主丛。

上式第一个映射是 $g$ 的左乘。映射 $\a'_{j}$ 与 $\io'$ 由 $\a_{j}$ 与 $\io$ 诱导得到，因为按定义 $\G(x_{j})$ 与 $\{x\}_{R}$ 不交。最后，$\cE'_{j}|_{D_{x,R}}$ 的平凡化 $\ka'_{j}$ 就是由 $\cE'_{j}$ 的粘合构造给出的重言平凡化。

\subsection{中心作用}\label{sss:central action}

设 $Z$ 是 $G$ 的中心。对每个 $x\in \Sig$，令 $\bK_{Z,x}=\bK_{x}\cap L^{+}_{x}Z$，这是 $Z$ 在 $x$ 处的水平群。记 $\bK_{Z}$ 为诸 $\bK_{Z,x}$（$x\in \Sig$）的汇集。由于 $Z$ 交换，叠 $\Bun_{Z,\bK_{Z}}$ 是群叠。特别地，$\Bun_{Z,\bK_{Z}}(k)$ 是 Picard 群胚。

$\Bun_{Z,\bK_{Z}}$ 在 $\Bun_{G,\bK}$ 上有典范作用。为看出这一点，先考虑如下情形：对每个 $x\in \Sig$，$\bK_{x}$ 是主同余子群 $L^{(n_{x})}_{x}G$，其中 $n_{x}\in\NN$。此时 $\Bun_{G,\bK}=\Bun_{G,D}$，其中 $D=\sum n_{x}x$，且 $\Bun_{Z,\bK_{Z}}=\Bun_{Z,D}$。

考虑 $S$-点 $\cE\in\Bun_{G,D}(S)$，它对应于 $X\times S$ 上的右 $G$-主丛 $\cE$ 连同 $\cE|_{D\times S}$ 的平凡化 $\io_{D}$。再考虑 $S$-点 $\cQ\in \Bun_{Z,D}(S)$，它对应于 $X\times S$ 上的右 $Z$-主丛 $\cQ$ 连同 $\cQ|_{D\times S}$ 的平凡化 $\ka_{D}$。由于 $Z$ 在 $G$ 中中心，它在 $\cE$ 上的自然作用与 $G$ 的作用交换，故 $Z$ 通过 $G$-主丛自同构作用于 $\cE$。

作扭纤维积 ${}_{\cQ}\cE:=\cQ\times^{Z}_{X\times S}\cE$，即 $\cQ\times_{X\times S}\cE$ 对 $Z$ 的对角作用的商。$\cE$ 上的右 $G$-作用使 ${}_{\cQ}\cE$ 具有 $G$-主丛结构。利用 $\cE|_{D\times S}$ 与 $\cQ|_{D\times S}$ 的平凡化 $\io_{D}$ 与 $\ka_{D}$，得到 ${}_{\cQ}\cE|_{D\times S}$ 的平凡化 $\io'_{D}$。于是定义作用 $(\cQ,\ka_{D})\cdot (\cE,\io_{D})$ 的结果为 $({}_{\cQ}\cE, \io'_{D})$。

对更一般的水平结构，仍可如下得到所需的作用。对每个 $x$ 选取足够小的主同余子群 $L^{(n_{x})}_{x}G\subset \bK_{x}$，并如 \S\ref{sss:level} 那样把 $\Bun_{G,\bK}$（相应地 $\Bun_{Z,\bK_{Z}}$）实现为 $\Bun_{G,D}$（相应地 $\Bun_{Z,D}$）对 $G_{D}$（相应地 $Z_{D}$）的某个代数子群的商。细节从略。

现在考虑上述构造的 shtuka 版本。我们断言，$\Bun_{Z,\bK_{Z}}(k)$ 在 $\Sht^{(I_{1},\cdots, I_{r})}_{G,\bK}$ 上有典范作用。事实上，对任意 $S$-点
\begin{equation*}
\xi=((x_{i})_{i\in I}, (\cE_{j})_{0\le j\le r}, (\a_{j})_{1\le j\le r}, \io)
\end{equation*}
（属于 $\Sht^{(I_{1},\cdots, I_{r})}_{G,\bK}$）以及 $\cQ\in \Bun_{Z,\bK_{Z}}(k)$，用上述构造得到 ${}_{\cQ}\cE_{j}\in \Bun_{G,\bK}(S)$。

原有的同构 $\a_{j}: \cE_{j-1}|_{X\times S-\cup_{i\in I_{j}}\G(x_{i})}\isom  \cE_{j}|_{X\times S-\cup_{i\in I_{j}}\G(x_{i})}$ 诱导同构
\begin{equation*}
\a'_{j}: {}_{\cQ}\cE_{j-1}|_{X\times S-\cup_{i\in I_{j}}\G(x_{i})}\isom  {}_{\cQ}\cE_{j}|_{X\times S-\cup_{i\in I_{j}}\G(x_{i})},
\end{equation*}
而 $\io:\cE_{r}\isom{}^{\t}\cE_{0}$ 诱导同构
\begin{equation*}
\io': {}_{\cQ}\cE_{r}\cong {}_{{}^{\t}\cQ}({}^{\t}\cE_{0})={}^{\t}({}_{\cQ}\cE_{0}),
\end{equation*}
因为 ${}^{\t}\cQ=\cQ$。于是 $\cQ$ 在 $\Sht^{(I_{1},\cdots, I_{r})}_{G,\bK}$ 上的作用把 $\xi$ 送到
\begin{equation*}
\xi':=((x_{i})_{i\in I}, ({}_{\cQ}\cE_{j})_{0\le j\le r}, (\a'_{j})_{1\le j\le r}, \io')\text{。}
\end{equation*}

容易看出，对任意支配余权组 $\un\l\in (\xcoch(T)^{+})^{I}$，$\Bun_{Z,\bK_{Z}}(k)$ 在 $\Sht^{(I_{1},\cdots, I_{r})}_{G,\bK}$ 上的作用保持有界版本 $\Sht^{(I_{1},\cdots, I_{r}),\le\un\l}_{G,\bK}$。

\begin{exam} 考虑 $G=\GL_{n}$ 且无水平结构的情形。此时 $\Bun_{Z}=\Pic_{X}$，而线丛 $\cL\in \Pic_{X}(k)$ 在 $\Sht^{(I_{1},\cdots, I_{r})}_{\GL_{n}}$ 上的作用就是把每个向量丛 $\cV_{j}$ 与 $\cL$ 作张量积。

另一方面，考虑 $G=\SL_{n}$ 且无水平结构的情形。此时 $\Bun_{Z}\cong \Pic_{X}[n]$，即 $X$ 的 Picard 叠的 $n$-挠子叠。对 $\cL\in \Bun_{Z}(k)=\Pic_{X}[n](k)$，即带同构 $\e: \cO_{X}\cong \cL^{\ot n}$ 的线丛 $\cL$，它在 $\Sht^{(I_{1},\cdots, I_{r})}_{\SL_{n}}$ 上的作用仍由把每个向量丛 $\cV_{j}$ 与 $\cL$ 作张量积给出。注意，$\det(\cV_{j})$ 原有的平凡化连同 $\e$ 诱导了 $\det(\cV_{j}\ot \cL)\cong \det(\cV_{j})\ot \cL^{\ot n}$ 的平凡化。
\end{exam}

\subsection{Hecke 对应}

\ssszh{球面 Hecke 对应}

设 $x\notin \Sig$，$K_{x}=G(\cO_{x})$，并取 $[g_{x}]\in K_{x}\bs G(F_{x})/K_{x}$。定义
\begin{equation*}
 \Sht^{(I_{1},\cdots, I_{r})}_{G,\bK}[g_{x}]:= \Sht^{(I_{1},\cdots, I_{r})}_{G,\bK,\infty x}/(K_{x}\cap g_{x}K_{x}g_{x}^{-1}).
\end{equation*}
我们把 $\Sht^{(I_{1},\cdots, I_{r})}_{G,\bK}[g_{x}]$ 看作对应
\begin{equation*}
\begin{tikzcd}
& \Sht^{(I_{1},\cdots, I_{r})}_{G,\bK}[g_{x}] \arrow[dl, "\oll{c}"'] \arrow[dr, "\orr{c}"] & \\
\Sht^{(I_{1},\cdots, I_{r})}_{G,\bK}|_{(U-\{x\})^{I}} & & \Sht^{(I_{1},\cdots, I_{r})}_{G,\bK}|_{(U-\{x\})^{I}}
\end{tikzcd}
\end{equation*}
其中 $\oll{c}$ 是对应于包含 $K_{x}\cap g_{x}K_{x}g_{x}^{-1}$ 的自然投影；映射 $\orr{c}$ 是复合
\begin{equation*}
 \Sht^{(I_{1},\cdots, I_{r})}_{G,\bK,\infty x}/(K_{x}\cap g_{x}K_{x}g_{x}^{-1})\xr{\cdot g_{x}} \Sht^{(I_{1},\cdots, I_{r})}_{G,\bK,\infty x}/(g_{x}^{-1}K_{x}g_{x}\cap K_{x})\to \Sht^{(I_{1},\cdots, I_{r})}_{G,\bK}.
\end{equation*}
这里第一个映射用到 \S\ref{sss:full x} 中构造的右 $G(F_{x})$-作用；第二个映射对应于包含 $g^{-1}_{x}K_{x}g_{x}\cap K_{x}\subset K_{x}$。

\ssszh{分歧点处的 Hecke 对应}

对 $x\in \Sig$ 与 $[g_{x}]\in  K_{x}\bs G(F_{x})/K_{x}$，定义对应
\begin{equation}\label{Hk x Sig}
\begin{tikzcd}
& \Sht^{(I_{1},\cdots, I_{r})}_{G,\bK}[g_{x}] \arrow[dl, "\oll{c}"'] \arrow[dr, "\orr{c}"] & \\
\Sht^{(I_{1},\cdots, I_{r})}_{G,\bK} & &  \Sht^{(I_{1},\cdots, I_{r})}_{G,\bK}
\end{tikzcd}
\end{equation}
为此令
\begin{equation*}
\Sht^{(I_{1},\cdots, I_{r})}_{G,\bK}[g_{x}]:=\Sht^{(I_{1},\cdots, I_{r})}_{G,\bK^{x}, \infty x}/(K_{x}\cap g_{x}K_{x}g_{x}^{-1}).
\end{equation*}
这里 $\bK^{x}$ 是水平群组 $(\bK_{x})_{x\in \Sig-\{x\}}$。映射 $\oll{c}$ 是自然投影，而映射 $\orr{c}$ 是先作用 $g_{x}$、再取自然投影，与球面 Hecke 修正的情形相同。

对 $k$-概形 $S$，$\Sht^{(I_{1},\cdots, I_{r})}_{G,\bK}[g_{x}](S)$ 分类 $(x_{i})_{i\in I}$（其中 $x_{i}\in U(S)$）连同图表
\begin{equation*}
\begin{tikzcd}[column sep=large, row sep=large]
\cE_{0} \arrow[d, dashed, "f_{0}"] \arrow[r, dashed, "\a_{1}"] &
\cE_{1} \arrow[d, dashed, "f_{1}"] \arrow[r, dashed, "\a_{2}"] &
\cdots \arrow[r, dashed, "\a_{r}"] &
\cE_{r} \arrow[d, dashed, "f_{r}"] \arrow[r, "\io", "\sim"'] &
{}^{\t}\cE_{0} \arrow[d, dashed, "{}^{\t}f_{0}"] \\
\cE'_{0} \arrow[r, dashed, "\a'_{1}"] &
\cE'_{1} \arrow[r, dashed, "\a'_{1}"] &
\cdots \arrow[r, dashed, "\a'_{r}"] &
\cE'_{r} \arrow[r, "\io'", "\sim"'] &
{}^{\t}\cE'_{0}
\end{tikzcd}
\end{equation*}
这里上下两行都是 $\Sht^{(I_{1},\cdots, I_{r})}_{G,\bK}$ 的 $S$-点，支点为 $(x_{i})_{i\in I}$；竖直箭头是 $U\times S$ 上 $G$-主丛的同构
\begin{equation*}
f_{i}: \cE_{i}|_{U\times S}\isom \cE'_{i}|_{U\times S},
\end{equation*}
它可延拓为 $x'\in \Sigma-\{x\}$ 处 $\bK_{x'}$-水平结构的同构，并在 $x$ 处有相对位置 $[g_{x}]\in K_{x}\bs G(F_{x})/K_{x}$。

\ssszh{有界版本}

给定界 $\un\l\in (\xcoch(T)^{+})^{I}$，有带界 $\un\l$ 的 Hecke 对应的有界版本
\begin{equation*}
\Sht^{(I_{1},\cdots, I_{r}), \le \un\l}_{G,\bK}[g_{x}]
\end{equation*}
它视为 $\Sht^{(I_{1},\cdots, I_{r}), \le \un\l}_{G,\bK}|_{(U-\{x\})^{I}}$（若 $x\notin \Sig$）或 $\Sht^{(I_{1},\cdots, I_{r}), \le \un\l}_{G,\bK}$（若 $x\in \Sig$）的自对应。

\subsection{部分 Frobenius}

存在典范态射
\begin{equation}\label{pl Fr}
\Fr^{\pl}: \Sht^{(I_{1},\cdots, I_{r})}_{G,\bK}\to \Sht^{(I_{2},\cdots, I_{r}, I_{1})}_{G,\bK}
\end{equation}
它在 $S$-点层面的定义如下：把 $S$-点
\begin{equation*}
\begin{tikzcd}[column sep=small]
\xi=((x_{i})_{i\in I}, \quad \cE_{0} \arrow[r, dashed, "\a_{1}"] & \cdots \arrow[r, dashed, "\a_{r-1}"] & \cE_{r-1} \arrow[r, dashed, "\a_{r}"] & \cE_{r} \arrow[r, "\io", "\sim"'] & {}^{\t}\cE_{0}
\end{tikzcd}
\end{equation*}
送到
\begin{equation*}
\begin{tikzcd}[column sep=small]
\Fr^{\pl}(\xi)=( (x'_{i})_{i\in I}, \quad \cE_{1} \arrow[r, dashed, "\a_{2}"] & \cdots \arrow[r, dashed, "\a_{r}"] & \cE_{r} \arrow[r, dashed, "{{}^{\t}\a_{1}\c\io}"] & {}^{\t}\cE_{1} \arrow[r, equal, "{\io'=\id}"] & {}^{\t}\cE_{1}.
\end{tikzcd}
\end{equation*}
这里
\begin{equation*}
x'_{i}=\begin{cases} x_{i}\c\Fr_{S}, & i\in I_{1}\\
x_{i}, & i\notin I_{1}.\end{cases}
\end{equation*}
显然
\begin{equation*}
\underbrace{\Fr^{\pl}\c \Fr^{\pl}\c\cdots \c \Fr^{\pl}}_{r \text{ 次}}= \Fr\in \End(\Sht^{(I_{1},\cdots, I_{r})}_{G,\bK}).
\end{equation*}

有交换图
\begin{equation}\label{pl Fr UI}
\begin{tikzcd}
\Sht^{(I_{1},\cdots, I_{r})}_{G,\bK} \arrow[r, "\Fr^{\pl}"] \arrow[d, "\pi^{(I_{1},\cdots, I_{r})}"] & \Sht^{(I_{2},\cdots, I_{1})}_{G,\bK} \arrow[d, "\pi^{(I_{2},\cdots, I_{1})}"] \\
U^{I} \arrow[r, "\Fr_{I_{1}}"] & U^{I}
\end{tikzcd}
\end{equation}
这里 $\Fr_{I_{1}}:U^{I}\to U^{I}$ 是在 $I_{1}$ 的各分量上取 Frobenius、在其余分量上取恒等的映射。由于 $U^{I}$ 上的 $\Fr_{I_{1}}$ 与 $\Fr^{\pl}$ 都是万有同胚，该图是拓扑笛卡尔图。

\section{上同调上的结构}\label{s:coho}

本节总结 $G$-shtuka 模叠的上同调（相对于支点空间而言）上的基本对称性：分解性、Hecke 代数作用与部分 Frobenius 作用。在 C.Xue 的讲义 \cite{Xue} 中，这些对称性将用于进一步研究 $G$-shtuka 模叠的上同调。

\subsection{Shtuka 模叠的上同调}

\ssszh{Satake 层}

回忆 Lusztig、Beilinson--Drinfeld、Ginzburg、Mirkovi\'c--Vilonen 的几何 Satake 等价（见 \cite{MV}）。有 $\Qlbar$ 上的张量范畴等价
\begin{equation*}
\Sat: \Rep(\dG,\Qlbar)\isom \Perv_{\Aut(D)}(\QGr_{G})
\end{equation*}
其中右端的幺半结构由卷积给出，交换约束由 Beilinson 与 Drinfeld 的融合构造给出。这里 $\QGr_{G}$ 上的 $\Aut(D)$-等变反常层，指的是 $\Gr_{G}$ 上 $L^{+}G\rtimes\Aut(D)$-等变的反常层。对任意有限集 $I$，定义函子
\begin{equation*}
\Sat_{I}: \Rep(\dG^{I},\Qlbar)\to\Perv((\QGr_{G}/\Aut(D))^{I}).
\end{equation*}
对 $V=\boxtimes_{i\in I}V_{i}$（其中 $V_{i}\in \Rep(\dG)$），定义 $\Sat_{I}(V)=\boxtimes_{i\in I}\Sat(V_{i})$，再按可加性延拓。

下面在 $\Hk^{I}_{G,\bK}$ 上定义层的复形。设 $r=|I|$，并取双射 $\s: I\isom \{1,2,\cdots, r\}$。令 $I_{j}=\{\s^{-1}(j)\}\subset I$，便得到把 $I$ 分成单点的分解 $I=I_{1}\coprod\cdots\coprod I_{r}$。回忆 \eqref{ev Hk singleton} 中的相对位置映射
\begin{equation*}
\ev^{(I_{1},\cdots, I_{r})}_{\Hk}: \Hk^{(I_{1},\cdots, I_{r})}_{G,\bK}\to (\QGr_{G}/\Aut(D))^{I}.
\end{equation*}
定义拉回
\begin{equation*}
\IC^{(I_{1},\cdots, I_{r})}_{\Hk}(V):=(\ev^{(I_{1},\cdots, I_{r})}_{\Hk})^{*}\Sat_{I}(V)\in D_{c}(\Hk^{(I_{1},\cdots, I_{r})}_{G,\bK}, \Qlbar).
\end{equation*}
若 $V=\boxtimes_{i\in I}V_{i}$ 且每个 $V_{i}$ 不可约、最高权为 $\l_{i}\in\xcoch(T)^{+}$，则 $\IC^{(I_{1},\cdots, I_{r})}_{\Hk}(V)$ 支撑在闭子叠 $\Hk^{(I_{1},\cdots, I_{r}), \le\un\l}_{G,\bK}$ 上，其中 $\un\l=(\l_{i})_{i\in I}$。

又有投影
\begin{equation*}
\pi^{(I_{1},\cdots, I_{r})}_{I}: \Hk^{(I_{1},\cdots, I_{r})}_{G,\bK}\to \Hk^{I}_{G,\bK}.
\end{equation*}
于是定义
\begin{equation*}
\IC^{I,\s}_{\Hk}(V):=\pi^{(I_{1},\cdots, I_{r})}_{I, !}\IC^{(I_{1},\cdots, I_{r})}_{\Hk}(V)\in D_{c}(\Hk^{I}_{G,\bK}, \Qlbar).
\end{equation*}
同样，若 $V=\boxtimes_{i\in I}V_{i}$ 且 $V_{i}$ 不可约、最高权为 $\l_{i}\in\xcoch(T)^{+}$，则 $\IC^{I}_{\Hk}(V)$ 支撑在 $\Hk^{I, \le\un\l}_{G,\bK}$ 上，其中 $\un\l=(\l_{i})_{i\in I}$。

\begin{lem}\label{l:IC Hk}
对任意 $V\in \Rep(\dG^{I})$，复形 $\IC^{I,\s}_{\Hk}(V)$ 在相差一个典范同构的意义下与 $I$ 上顺序 $\s$ 的选取无关。我们记之为 $\IC^{I}_{\Hk}(V)$。

此外，若 $V=\boxtimes_{i\in I}V_{i}$ 且 $V_{i}$ 不可约、最高权为 $\l_{i}$，则 $\IC^{I,\s}_{\Hk}(V)$ 在相差一个平移的意义下同构于 $\Hk^{I, \le\un\l}_{G,\bK}$ 的相交复形，其中 $\un\l=(\l_{i})_{i\in I}$。
\end{lem}
\begin{proof}[证明概要] 设 $V=\boxtimes_{i\in I}V_{i}$ 且 $V_{i}$ 不可约、最高权为 $\l_{i}$。对分成单点的分解 $I=I_{1}\coprod\cdots\coprod I_{r}$，由命题 \ref{p:fm sm}，相对位置映射 $\ev^{(I_{1},\cdots, I_{r})}_{\Hk}$ 形式光滑。因此 $\IC^{(I_{1},\cdots, I_{r})}_{\Hk}(V)$ 在相差一个平移的意义下典范同构于 $\Hk^{(I_{1},\cdots, I_{r}), \le\un\l}_{G,\bK}$ 的相交复形。映射 $\pi^{(I_{1},\cdots, I_{r})}_{I}$ 在 \cite[命题 4.2 之后]{MV} 的意义下分层小，这意味着 $\IC^{I,\s}_{\Hk}(V)$ 在相差一个平移的意义下典范同构于 $\Hk^{I, \le\un\l}_{G,\bK}$ 的相交复形。故它与 $\s$ 的选取无关。
\end{proof}

更一般地，若 $I=I_{1}\coprod\cdots \coprod I_{r}$ 是 $I$ 的任意分解，则可给每个 $I_{j}$ 排序，并把这分解细化为分成单点的分解 $I=I'_{1}\coprod\cdots \coprod I'_{s}$。于是有双射 $\s: I\isom \{1,2,\cdots, s\}$，使 $I'_{j}=\{\s^{-1}(j)\}$（$1\le j\le s$），且 $\s(I_{j})$ 是 $\{1,2,\cdots, s\}$ 的连续子集。有典范投影
\begin{equation*}
\pi^{(I'_{1},\cdots, I'_{s})}_{(I_{1},\cdots, I_{r})}: \Hk^{(I'_{1},\cdots, I'_{s})}_{G,\bK}\to \Hk^{(I_{1},\cdots, I_{r})}_{G,\bK}.
\end{equation*}
对 $V\in \Rep(\dG^{I}, \Qlbar)$ 定义
\begin{equation*}
\IC^{(I_{1},\cdots, I_{r}),\s}_{\Hk}(V):=\pi^{(I'_{1},\cdots, I'_{s})}_{(I_{1},\cdots, I_{r}), !}\IC^{(I'_{1},\cdots, I'_{s})}_{\Hk}(V)\in D_{c}(\Hk^{(I_{1},\cdots, I_{r})}_{G,\bK}, \Qlbar).
\end{equation*}

与引理 \ref{l:IC Hk} 相同的论证给出以下结论。

\begin{lem}\label{l:IC Hk refine} 设 $V\in \Rep(\dG^{I})$。
\begin{enumerate}
\item $\Hk^{(I_{1},\cdots, I_{r})}_{G,\bK}$ 上的复形 $\IC^{(I_{1},\cdots, I_{r}),\s}_{\Hk}(V)$ 在相差一个典范同构的意义下，与细化分解 $I=I_{1}\coprod\cdots \coprod I_{r}$ 的 $I$ 上顺序 $\s$ 的选取无关。记此复形为 $\IC^{(I_{1},\cdots, I_{r})}_{\Hk}(V)$。

\item 设 $I=I'_{1}\coprod\cdots \coprod I'_{r'}$ 是分解 $I=I_{1}\coprod\cdots \coprod I_{r}$ 的任意细化。则存在典范同构
\begin{equation*}
\IC^{(I_{1},\cdots, I_{r})}_{\Hk}(V)\cong \pi^{(I'_{1},\cdots, I'_{r'})}_{(I_{1},\cdots, I_{r}), !}\IC^{(I'_{1},\cdots, I'_{r'})}_{\Hk}(V).
\end{equation*}
且这些同构与进一步的细化相容。

\item 若 $V=\boxtimes_{i}V_{i}$ 且每个 $V_{i}$ 不可约、最高权为 $\l_{i}$，则 $\IC^{(I_{1},\cdots, I_{r})}_{\Hk}(V)$ 在相差一个平移的意义下同构于 $\Hk^{(I_{1},\cdots, I_{r}),\le\un\l}_{G,\bK}$ 的相交复形，其中 $\un\l=(\l_{i})_{i\in I}$。
\end{enumerate}
\end{lem}

\ssszh{$\Sht^{I}_{G}$ 的上同调复形}

回忆典范映射
\begin{equation*}
\r^{(I_{1},\cdots, I_{r})}: \Sht^{(I_{1},\cdots, I_{r})}_{G,\bK}\to \Hk^{(I_{1},\cdots, I_{r})}_{G,\bK}.
\end{equation*}
对 $V\in \Rep(\dG^{I}, \Qlbar)$ 定义
\begin{equation*}
\IC^{(I_{1},\cdots, I_{r})}_{\Sht}(V):=(\r^{(I_{1},\cdots, I_{r})})^{*}\IC^{(I_{1},\cdots, I_{r})}_{\Hk}(V)\in D_{c}(\Sht^{(I_{1},\cdots, I_{r})}_{G,\bK}).
\end{equation*}

另一种做法是：当 $I=I_{1}\coprod\cdots \coprod I_{r}$ 是分成单点的分解时，可考虑 $\Sht^{(I_{1},\cdots, I_{r})}_{G,\bK}$ 的相对位置映射 \eqref{ev Sht singleton}。于是
\begin{equation*}
\IC^{(I_{1},\cdots, I_{r})}_{\Sht}(V)\cong (\ev_{\Sht}^{(I_{1},\cdots, I_{r})})^{*}\Sat_{I}(V).
\end{equation*}

\begin{lem}\label{l:IC Sht refine} 设 $V\in \Rep(\dG^{I})$。
\begin{enumerate}
\item 对分解 $I=I_{1}\coprod\cdots \coprod I_{r}$ 的任意细化 $I=I'_{1}\coprod\cdots \coprod I'_{r'}$，存在典范同构
\begin{equation*}
\IC^{(I_{1},\cdots, I_{r})}_{\Sht}(V)\cong \pi^{(I'_{1},\cdots, I'_{r'})}_{(I_{1},\cdots, I_{r}), !}\IC^{(I'_{1},\cdots, I'_{r'})}_{\Sht}(V),
\end{equation*}
且与进一步的细化相容。
\item 若 $V=\boxtimes_{i\in I}V_{i}$，其中 $V_{i}$ 不可约、最高权为 $\l_{i}$，则 $\IC^{(I_{1},\cdots, I_{r})}_{\Sht}(V)$ 在相差 $\#I$ 的平移的意义下同构于 $\Sht^{(I_{1},\cdots, I_{r}), \le\un\l}_{G,\bK}$ 的相交复形，其中 $\un\l=(\l_{i})_{i\in I}$。
\end{enumerate}
\end{lem}
\begin{proof}
(1) 由引理 \ref{l:IC Hk refine} 以及笛卡尔方图
\begin{equation*}
\begin{tikzcd}[column sep=large, row sep=large]
\Sht^{(I'_{1},\cdots, I'_{r'})}_{G,\bK} \arrow[r, "\r^{(I'_{1},\cdots, I'_{r'})}"] \arrow[d, "\pi^{(I'_{1},\cdots, I'_{r'})}_{(I_{1},\cdots, I_{r})}"] & \Hk^{(I'_{1},\cdots, I'_{r'})}_{G,\bK} \arrow[d, "\pi^{(I'_{1},\cdots, I'_{r'})}_{(I_{1},\cdots, I_{r})}"] \\
\Sht^{(I_{1},\cdots, I_{r})}_{G,\bK} \arrow[r, "\r^{(I_{1},\cdots, I_{r})}"] & \Hk^{(I_{1},\cdots, I_{r})}_{G,\bK}
\end{tikzcd}
\end{equation*}
的固有基变换得到。

(2) 先设各 $I_{j}$ 都是单点集，此时结论由命题 \ref{p:fm sm} 得到。

现考虑一般的分解 $I=I_{1}\coprod\cdots\coprod I_{r}$。把每个 $I_{j}$ 分成单点，再合起来便得到分解 $I=I'_{1}\coprod\cdots\coprod I'_{r'}$，其中每个 $I'_{i}$ 是单点，且 $I_{j}$ 是若干连续部分的并。由命题 \ref{p:local model}，映射 $\pi: \Sht^{(I'_{1},\cdots, I'_{r'}),\le\un\l}_{G,\bK}\to\Sht^{(I_{1},\cdots, I_{r}),\le\un\l}_{G,\bK}$ 在 \'etale 局部同构于映射 $\Gr^{(I'_{1},\cdots, I'_{r'}),\un\le\l}_{G}\to \Gr^{(I_{1},\cdots, I_{r}),\un\le\l}_{G}$，而后者分层小。故由 (1)，
\begin{equation*}
\begin{aligned}
&\IC^{(I_{1},\cdots, I_{r})}_{\Sht}(V)\cong \pi_{!}\IC^{(I'_{1},\cdots, I'_{r'})}_{\Sht}(V)\\
&\qquad\cong\pi_{!}\IC(\Sht^{(I'_{1},\cdots, I'_{r'}),\le\un\l}_{G,\bK})[-\#I]
\cong\IC(\Sht^{(I_{1},\cdots, I_{r}),\le\un\l}_{G,\bK})[-\#I].
\end{aligned}
\end{equation*}
\end{proof}

\begin{dfn} 设 $I$ 是带分解 $I=I_{1}\coprod\cdots \coprod I_{r}$ 的有限集。回忆映射 $\pi^{I}: \Sht^{(I_{1},\cdots, I_{r})}_{G,\bK}\to U^{I}$。定义 ind-可构造复形
\begin{equation*}
\cH^{(I_{1},\cdots, I_{r})}_{G,\bK}(V):=\pi^{I}_{!}\IC^{(I_{1},\cdots, I_{r})}_{\Sht}(V)\in D(U^{I}, \Qlbar).
\end{equation*}
当讨论中 $(G,\bK)$ 固定时，就把 $\cH^{(I_{1},\cdots, I_{r})}_{G,\bK}(V)$ 简记为 $\cH^{(I_{1},\cdots, I_{r})}(V)$。

特别地，有 ind-可构造复形 $\cH^{I}(V)\in D(U^{I}, \Qlbar)$。
\end{dfn}

\begin{rmk}
为看出 $\cH^{I}(V)$ 是 ind-可构造的，我们使用 \S\ref{sss:HN} 中讨论的 Harder--Narasimhan 截断。一般水平结构可归结为主同余水平结构的情形，故不妨设 $\bK_{x}\subset L^{+}_{x}G$ 对一切 $x\in \Sig$ 成立。对 $\mu\in \HN$，令 ${}^{\le \mu}\pi^{I}: {}^{\le\mu}\Sht^{I}_{G,\bK}\to U^{I}$ 为 $\pi^{I}$ 的限制。令 ${}^{\le \mu}\IC_{\Sht}^{I}(V)$ 为 $\IC^{I}_{\Sht}(V)$ 到 ${}^{\le\mu}\Sht^{I}_{G,\bK}$ 的限制。则 ${}^{\le\mu}\cH^{I}(V):={}^{\le \mu}\pi^{I}_{!}({}^{\le \mu}\IC_{\Sht}^{I}(V))$ 可构造，因为 ${}^{\le \mu}\IC_{\Sht}^{I}(V)$ 的支撑在 $U^{I}$ 上有限型。于是有
\begin{equation}\label{HV colim}
\cH^{I}(V)=\colim_{\mu\in \HN}{}^{\le\mu}\cH^{I}(V)
\end{equation}
作为 $U^{I}$ 上的 ind-可构造复形。这里对 $\mu\le \mu'$，过渡映射 ${}^{\le\mu}\cH^{I}(V)\to {}^{\le\mu'}\cH^{I}(V)$ 由开浸入 ${}^{\le\mu}\Sht^{I}_{G,\bK}\incl {}^{\le\mu'}\Sht^{I}_{G,\bK}$ 诱导。
\end{rmk}

\begin{cor}[引理 \ref{l:IC Sht refine} 的推论]\label{c:same HV}
对分解 $I=I_{1}\coprod\cdots \coprod I_{r}$ 的任意细化 $I=I'_{1}\coprod\cdots \coprod I'_{r'}$，存在典范同构
\begin{equation*}
\cH^{(I_{1},\cdots, I_{r})}(V)\cong \cH^{(I'_{1},\cdots, I'_{r'})}(V),
\end{equation*}
且与进一步的细化相容。特别地，对任意分解 $I=I_{1}\coprod\cdots \coprod I_{r}$，存在典范同构
\begin{equation*}
\cH^{(I_{1},\cdots ,I_{r})}(V)\cong\cH^{I}(V)\in D(U^{I}, \Qlbar).
\end{equation*}
\end{cor}
\begin{proof}
对引理 \ref{l:IC Sht refine}(1) 中的同构作用 $\pi^{(I_{1},\cdots, I_{r})}_{!}$。
\end{proof}

% ch04  §5.2–§5.4 / §6.1–§6.3（原文 1648–2027 行）

\subsection{Hecke 作用}

设 $V\in \Rep(\dG^{I})$。

\ssszh{}
对任意 $x\in \Sig$，$G(F_{x})$ 上具紧支集、双 $K_{x}$-不变的函数构成的向量空间 $\bH_{x}:=\Qlbar[K_{x}\bs G(F_{x})/K_{x}]$ 在卷积下是 $\Qlbar$-代数，单位元为 $\one_{K_{x}}$。有环同态
\begin{equation*}
\bH_{x}\to \End(\cH^{I}(V))
\end{equation*}
定义如下。

由推论 \ref{c:same HV}，只需在 $\cH^{(I_{1},\cdots, I_{r})}(V)$ 上构造 $\cH_{x}$ 的作用\zy{原文此处作 $\cH_{x}$；据同段前文，作用在 $\cH^{I}(V)$ 上的应是 Hecke 代数 $\bH_{x}$，$\cH$ 与 $\bH$ 疑为笔误。}，其中各 $I_{j}$ 都是单点集。又可将情形化归为 $V$ 不可约，此时 $V=\boxtimes_{i\in I}V_{i}$，其中 $V_{i}$ 是 $\dG$ 的不可约表示，最高权为 $\l_{i}$。令 $\un\l=(\l_{i})_{i\in I}$。对 $[g_{x}]\in K_{x}\bs G(F_{x})/K_{x}$，回忆 Hecke 对应 \eqref{Hk x Sig}，其中 $\oll{c}$ 与 $\orr{c}$ 由引理 \ref{l:change level Sig} 都是有限 étale 的。

类似地可以构造相对位置映射
\begin{equation*}
\ev_{\Sht[g_{x}]}^{(I_{1},\cdots, I_{r})}: \Sht^{(I_{1},\cdots, I_{r})}_{G,\bK}[g_{x}]\to (\QGr_{G}/\Aut(D))^{I}
\end{equation*}
它与映射 $\oll{c}$、$\orr{c}$ 以及 \eqref{ev Sht singleton} 中 $\Sht^{(I_{1},\cdots, I_{r})}_{G,\bK}$ 的相对位置映射相容。于是得到典范同构
\begin{equation}\label{corr IC}
\oll{c}^{*}\IC_{\Sht}^{I}(V)\isom \orr{c}^{*}\IC_{\Sht}^{I}(V)\cong\orr{c}^{!}\IC_{\Sht}^{I}(V).
\end{equation}
$\one_{K_{x}g_{x}K_{x}}$ 在 $\cH^{I}(V)$ 上的作用定义为如下复合：
\begin{equation*}
\begin{aligned}
\cH^{I}(V)=\pi^{I}_{!}\IC_{\Sht}^{I}(V)&\to \pi^{I}_{!}\oll{c}_{!}\oll{c}^{*}\IC_{\Sht}^{I}(V)\\
&\xr{\eqref{corr IC}} \pi^{I}_{!}\orr{c}_{!}\orr{c}^{!}\IC_{\Sht}^{I}(V)\to \pi^{I}_{!}\IC_{\Sht}^{I}(V)=\cH^{I}(V)
\end{aligned}
\end{equation*}
这里用到了伴随的单位与余单位：$\id \to \oll{c}_{*}\oll{c}^{*}\cong \oll{c}_{!}\oll{c}^{*}$（因为 $\oll{c}$ 是 proper 的），以及 $\orr{c}_{!}\orr{c}^{!}\to\id$。

\ssszh{球 Hecke 作用}
对 $x\in |U|=|X|-\Sig$，类似的构造给出环同态
\begin{equation*}
\bH_{x}:=\Qlbar[G(\cO_{x})\bs G(F_{x})/G(\cO_{x})]\to \End(\cH^{I}(V)|_{(U-\{x\})^{I}}).
\end{equation*}

\begin{rmk} 对 $x\in |U|$，$\bH_{x}$ 在 $\cH^{I}(V)|_{(U-\{x\})^{I}}$ 上的作用已被 V. Lafforgue 延拓到整个 $\cH^{I}(V)$ 上：办法是把它看成游走算子的一个特例。
\end{rmk}

设 $\Sig'\subset |U|$ 是任意有限个位点构成的集合。把上面定义的各 Hecke 作用合起来，得到环同态
\begin{equation*}
\bH_{\Sig'}=\ot_{x\in \Sig'}\bH_{x}\to \End(\cH^{I}(V)|_{(U-\Sig')^{I}}).
\end{equation*}
取 $\Sig'$ 递增时的余极限，就得到整体 Hecke 代数在 $\cH^{I}(V)$ 限制到一般点 $\y^{I}\to X^{I}$ 之上的作用
\begin{equation*}
\bH_{G,\bK}:=\ot'_{x\in |X|}\bH_{x}\to \End(\cH^{I}(V)|_{\y^{I}}).
\end{equation*}

\subsection{因子化结构}

\ssszh{预期的形态}
设 $V=\boxtimes_{i\in I}V_{i}$，其中 $V_{i}\in \Rep(\dG)$。则\emph{预期} $\cH^{I}(V)$ 满足下述因子化性质。设 $\pi$ 是 $G(\AA)$ 的尖点自守表示。令 $\pi(\bK)\subset \pi$ 为在 $\prod_{x\notin\Sig}G(\cO_{x})\times \prod_{x\in \Sig}K_{x}$ 下不变的那些向量构成的子空间，其中 $K_{x}=\bK_{x}(k_{x})$。注意 $\pi(\bK)$ 是 $\bH_{G,\bK}$-模。由 V. Lafforgue 的工作 \cite{VLaff}，可以对 $\pi$ 附上一个半单 Langlands 参数 $\r_{\pi}$，它是 $U$ 上的 $\dG$-局部系统。利用 V. Lafforgue 的游走算子，可以定义 $H^{i}\cH^{I}(V)$（即 $\cH^{I}(V)$ 的第 $i$ 个上同调层，$i\in \ZZ$）的 $\pi$-等型子层，记作 $H^{i}\cH^{I}(V)[\pi]$。作为初步近似，可以把 $H^{i}\cH^{I}(V)[\pi]$ 定义为 $\Hom_{\bH_{G,\bK}}(\pi(\bK), H^{i}\cH^{I}(V))$（它可能比真正的 $H^{i}\cH^{I}(V)[\pi]$ 更大；微妙之处在于 $\bH_{G,\bK}$-模 $\pi(\bK)$ 未必能决定 $\r_{\pi}$ 的同构类）。

我们作一个简化假设：$\Aut(\r_{\pi})=Z\dG$。例如当 $G=\GL_{n}$ 且 $\pi$ 尖点时（此时 $\r_{\pi}$ 是不可约的秩 $n$ 局部系统），已知确实如此。于是预期
\begin{equation*}
H^{0}\cH^{I}(V)[\pi]\stackrel{?}{\simeq} M_{\pi}\ot \boxtimes_{i\in I}(V_{i,\r_{\pi}}),
\end{equation*}
其中 $V_{i,\r_{\pi}}$ 是由 $\r_{\pi}$ 与 $\dG$-表示 $V_{i}$ 诱导出的 $U$ 上的局部系统，$M_{\pi}$ 是依赖于 $\pi$ 的有限维向量空间。还预期此时 $H^{i}\cH^{I}(V)[\pi]=0$ 对 $i\ne0$ 成立。更精确的陈述是数域上自守形式的 Arthur 重数公式的类比，见 Xinwen Zhu 的文章 \cite[\S4.7]{Zhu}。

上述预期特别地暗示：在 $\cH^{I}(V)$ 的几何一般茎上，有 $I$ 份基本群 $\pi_{1}(U,\ov \y)$ 的相互交换的作用。这一事实由 Cong Xue \cite{X} 证明，推广了 V. Lafforgue \cite{VLaff} 的结果。他们的证明依赖部分 Frobenius（将在 \S\ref{sss:pF} 中讨论）与 Drinfeld 引理。

\ssszh{在部分对角线上的行为}
下面把 $\cH^{I}(V)$ 限制到部分对角线上时的行为讲清楚。

有限集之间的映射 $\th: I\to J$ 诱导群同态 $\D_{\dG, \th}: \dG^{J}\to \dG^{I}$，从而诱导表示上的限制函子
\begin{equation*}
\Res_{\th}: \Rep(\dG^{I})\to \Rep(\dG^{J}).
\end{equation*}

\begin{prop}
设 $\th: I\surj J$ 是有限集之间的满射，相应地有部分对角映射 $\D_{\th,U}: U^{J}\incl U^{I}$。则有典范同构
\begin{equation*}
\D_{\th,U}^{*}\cH^{I}(V)\isom \cH^{J}(\Res_{\th}V).
\end{equation*}
此外，该同构对 $V\in \Rep(\dG^{I})$ 具有函子性，并与满射的复合 $I\surj J\surj K$ 相容。
\end{prop}
\begin{proof}[证明概要]
由 proper 基变换，只需对 $\Hk^{J}_{G,\bK}$ 上的 Satake 层验证相应的陈述：
\begin{equation*}
\IC^{I}_{\Hk}(V)|_{\Hk^{J}_{G,\bK}}\cong \IC^{J}_{\Hk}(\Res_{\th}(V)).
\end{equation*}
这里用同构 \eqref{Hk diagonal} 把 $\Hk^{J}_{G,\bK}$ 看作 $\Hk^{I}_{G,\bK}$ 的子叠（经由 $\D_{\th, U}: U^{J}\incl U^{I}$ 拉回）。取双射 $\s: J\isom \{1,\cdots, r\}$，并对 $1\le j\le r'$ 令 $I_{j}=\th^{-1}(\s'^{-1}(j))$。于是有笛卡尔方块
\begin{equation*}
\begin{tikzcd}
\Hk^{(\{1\},\cdots, \{r\})}_{G,\bK} \arrow[r,"\y"] \arrow[d] & \Hk^{(I_{1},\cdots, I_{r})}_{G,\bK} \arrow[d] \\
\Hk^{J}_{G,\bK} \arrow[r] & \Hk^{I}_{G,\bK}
\end{tikzcd}
\end{equation*}
由 proper 基变换与引理 \ref{l:IC Hk refine}，只需证明
\begin{equation}\label{IC fact b}
\y^{*}\IC^{(I_{1},\cdots, I_{r})}_{\Hk}(V)\cong \IC^{(\{1\},\cdots ,\{r\})}_{\Hk}(\Res_{\th}V).
\end{equation}

下面处理 $J$ 本身是单点集的情形。一般情形不再需要新的想法。于是设 $J=\{1\}$ 为单点集。取双射 $\s: I\isom \{1,\cdots, s\}$，并定义
\begin{equation*}
\Hk^{[s]}_{G,\bK}=\Hk^{(\{1\},\cdots,\{s\})}_{G,\bK}\times_{U^{s}}\D(U).
\end{equation*}
于是 $\Hk^{[s]}_{G,\bK}$ 分类带 $\bK$-水平结构的 $G$-主丛的修正 $\cE_{0}\dxr\cE_{1}\dxr\cdots\dxr\cE_{s}$，且这些修正都落在同一个点 $x\in U$ 上。令 $\QGr^{[s]}_{G}$ 是这样的模叠：其 $R$-点是圆盘 $D_{R}$ 上的 $G$-主丛 $(\cE_{0},\cdots, \cE_{s})$，并带有同构 $\a_{i}: \cE_{i-1}|_{D_{R}^{\times}}\isom\cE_{i}|_{D^{\times}_{R}}$（$1\le i\le s$）。于是有求值映射 $\ev^{[s]}: \Hk^{[s]}_{G,\bK}\to \QGr^{[s]}_{G}/\Aut(D)$，以及一个三个方块全为笛卡尔的图表
\begin{equation*}
\begin{tikzcd}
 & \QGr^{[s]}/\Aut(D) \arrow[rr,"\d"] \arrow[dd,"\g"' pos=0.22] & & (\QGr/\Aut(D))^{s} \\
\Hk^{[s]}_{G,\bK} \arrow[rr,"\wt\y" pos=0.62] \arrow[dd] \arrow[ur,"\ev^{[s]}"] & & \Hk^{(\{1\},\cdots,\{s\})}_{G,\bK} \arrow[ur,"\ev^{(\{1\},\cdots,\{s\})}"] \arrow[dd] & \\
 & \QGr/\Aut(D) & & \\
\Hk^{\{1\}}_{G,\bK} \arrow[ur,"\ev^{\{1\}}"] \arrow[rr,"\y"] & & \Hk^{I}_{G,\bK} &
\end{tikzcd}
\end{equation*}
映射 $\d$ 由 $\QGr^{[s]}\to (\QGr)^{s}$ 诱导，后者把 $((\cE_{i})_{0\le i\le r},(\a_{i})_{1\le i\le r})$ 送到 $(\cE_{i-1},\cE_{i}, \a_{i})_{1\le i\le s}$；映射 $\g$ 由 $\QGr^{[s]}\to \QGr$ 诱导，后者把 $((\cE_{i})_{0\le i\le r},(\a_{i})_{1\le i\le r})$ 送到 $(\cE_{0},\cE_{s}, \a_{r}\c\cdots\c\a_{1})$。

再由 proper 基变换，\eqref{IC fact b} 化归为同构
\begin{equation*}
\g_{!}\d^{*}\Sat_{\{1,\cdots , s\}}(V)\cong \Sat(\Res^{\dG^{s}}_{\dG} V).
\end{equation*}
由线性性，只需对 $V=\boxtimes_{i=1}^{s} V_{i}$ 证明；此时左端是张量范畴 $\Perv_{\Aut(D)}(\QGr)$ 中的卷积 $\Sat(V_{1})\star\cdots\star\Sat(V_{s})$，右端是 $\Sat(V_{1}\ot\cdots\ot V_{s})$。所求同构来自 Satake 等价的幺半结构。
\end{proof}

\subsection{部分 Frobenius}\label{sss:pF}
设 $I=I_{1}\coprod\cdots\coprod I_{r}$ 是分解为\emph{单点集}的分解。这给出双射 $\s: I\isom \{1,\cdots, r\}$，使 $I_{j}=\{\s^{-1}(j)\}$。回忆 \eqref{pl Fr} 中定义的部分 Frobenius 映射 $\Fr^{\pl}$。
我们有交换图表
\begin{equation}\label{pl Frob ev}
\begin{tikzcd}
\Sht^{(I_{1},\cdots, I_{r})}_{G,\bK} \arrow[r,"\ev^{(I_{1},\cdots, I_{r})}"] \arrow[d,"\Fr^{\pl}"'] & (\QGr/\Aut(D))^{I} \arrow[d,"{(\Fr, \id,\cdots, \id)}"] \\
\Sht^{(I_{2},\cdots, I_{r}, I_{1})}_{G,\bK} \arrow[r,"\ev^{(I_{2},\cdots, I_{1})}"] & (\QGr/\Aut(D))^{I}
\end{tikzcd}
\end{equation}
这里右侧的竖直映射在第一个因子 $I_{1}$ 上是 Frobenius，在其余因子上是恒等映射。

对最高权为 $\l$ 的不可约表示 $V_{\l}\in \Rep(\dG)$，在 $\Sat(V_{\l})$ 上固定一个 Weil 结构，使其在 $\Gr_{G,\l}$ 上的限制为 $\Qlbar[\j{2\r, \l}](\j{\r,\l})$（若 $\j{\r,\l}$ 是半整数，这里涉及 $q^{1/2}$ 的选取）。用加性把这些 Weil 结构延拓到所有 $V\in \Rep(\dG)$ 的 $\Sat(V)$ 上。更一般地，对任意有限集 $I$ 与 $V\in \Rep(\dG^{I})$，这些 Weil 结构延拓为 $\Sat_{I}(V)\in \Perv((\QGr/\Aut(D))^{I})$ 上的 Weil 结构，并与外张量积相容。

现在设 $V=\boxtimes_{i\in I}V_{i}\in \Rep(\dG^{I})$。对 $1\le j\le r$，把 $V_{\s^{-1}(j)}$ 也记作 $V_{j}$。由图表 \eqref{pl Frob ev} 并利用上面固定的 $\Sat(V_{1})$ 上的 Weil 结构，得到同构
\begin{equation}\label{pl Fr pullback IC}
\IC^{(I_{1},\cdots, I_{r})}_{\Sht}(V)\cong \Fr^{\pl*}\IC^{(I_{2},\cdots, I_{1})}_{\Sht}(V).
\end{equation}
另一方面，对拓扑笛卡尔图表 \eqref{pl Fr UI} 用 proper 基变换，给出同构
\begin{equation*}
\Fr_{I_{1}}^{*}\cH^{I}(V)= \Fr_{I_{1}}^{*}\pi^{(I_{2},\cdots, I_{1})}_{!}\IC^{(I_{2},\cdots, I_{1})}_{\Sht}(V)\isom  \pi^{(I_{1},\cdots, I_{r})}_{!}\Fr^{\pl*}\IC^{(I_{2},\cdots, I_{1})}_{\Sht}(V)
\end{equation*}
结合 \eqref{pl Fr pullback IC} 便得到同构
\begin{equation*}
\phi_{I_{1}}: \Fr_{I_{1}}^{*}\cH^{I}(V)\cong \pi^{(I_{1},\cdots, I_{r})}_{!}\Fr^{\pl*}\IC^{(I_{2},\cdots, I_{1})}_{\Sht}(V)\cong \pi^{(I_{1},\cdots, I_{r})}_{!}\IC^{(I_{1},\cdots, I_{r})}_{\Sht}(V)=\cH^{I}(V).
\end{equation*}
该同构与 $I-I_{1}$ 分解为单点集的方式无关。此外，这类同构按加性延拓到所有 $V\in \Rep(\dG^{I})$。

由于可以取双射 $\s$ 使 $I_{1}$ 成为任意给定的元素 $i\in I$，故对一切 $i\in I$ 与 $V\in\Rep(\dG^{I})$ 有典范同构
\begin{equation*}
\phi_{i}:  \Fr_{i}^{*}\cH^{I}(V)\isom  \cH^{I}(V).
\end{equation*}
这里 $\Fr_{i}: U^{I}\to U^{I}$ 在第 $i$ 个因子上是 Frobenius，在其余因子上是恒等映射。

下面的事实由定义容易验证。
\begin{prop}
同构族 $(\phi_{i})_{i\in I}$ 彼此在典范同构的意义下交换，它们的复合是 $\cH^{I}(V)$ 上的典范 Weil 结构，即由 $\Sat_{I}(V)$ 上的 Weil 结构诱导出的同构 $\Fr^{*}_{U^{I}}\cH^{I}(V)\isom  \cH^{I}(V)$。
\end{prop}

特别地，对任意子集 $J\subset I$ 有典范同构 $\phi_{J}: \Fr_{J}^{*}\cH^{I}(V)\isom  \cH^{I}(V)$，其中 $\Fr_{J}: U^{I}\to U^{I}$ 在 $U^{J}$ 上是 Frobenius，在其余因子上是恒等映射。它定义为各 $\phi_{i}$（$i\in J$）按任意次序的复合。

\section{进阶课题}\label{s:adv}
本节概述 shtuka 模空间的若干方面，这些方面值得进一步研究。

\subsection{紧化}
映射 $\pi^{(I_{1},\cdots, I_{r}),\le \un\l}_{G,\bK}: \Sht^{(I_{1},\cdots, I_{r}),\le \un\l}_{G,\bK}\to U^{I}$ 不是 proper 的，这给研究它的直接像复形、尤其是其通用纤维的上同调带来困难。例如，如果它是 proper 的，复形 $\cH^{I}(V)$ 就会自动是纯的，这将对函数域上自守形式的 Ramanujan–Petersson 猜想产生影响。

由于 $\pi^{(I_{1},\cdots, I_{r}),\le \un\l}_{G,\bK}$ 不是有限型的，要作紧化，必须先限制到 Harder–Narasimhan 截断 ${}^{\le \mu}\Sht^{(I_{1},\cdots, I_{r}),\le \un\l}_{G,\bK}$，如 \S\ref{sss:HN} 那样。

在 \cite{Dr-comp} 中，Drinfeld 对 $G=\GL_{2}$ 的 Drinfeld 例 \ref{ex:Dr} 构造了 $\Sht^{(I_{1},I_{2}),(\l_{1},\l_{2})}_{\GL_{2}, \bK}$（的截断）的紧化。在 \cite{Laff-comp} 中，无水平结构时，Drinfeld 的构造被推广到 $\GL_{n}$。在 \cite{LLaff} 中，带水平结构的 Drinfeld 例 \ref{ex:Dr} 的紧化被构造出来（性质稍逊），并用于证明函数域上的整体 Langlands 对应。在 \cite{NDT} 中，Ng\^o Dac Tuan 用几何不变量理论给出了紧化的另一种构造。

这些构造的基本想法是：放宽 \S\ref{sss:Sht points} 的模描述中的同构 $\io: \cE_{r}\isom {}^{\t}\cE_{0}$，把它换成一个 $G$ 所附的 Vinberg 幺半群中的点。当 $G=\GL_{n}$ 时，这样的点通常称为两个向量空间或向量丛之间的完全同态（complete homomorphism）。对一般的 $G$，Ng\^o Dac Tuan 与 Varshavsky 的一份未发表手稿沿这条路线给出了 shtuka 模空间紧化的构造。

据作者所知，描述 $\Sht^{(I_{1},\cdots, I_{r}),\le\un\l}_{G,\bK}$ 的紧化的奇点仍是一个公开问题。

\subsection{代数圈}
我们简要综述 shtuka 模叠上代数圈的一些例子。为简单起见，下面的讨论中不考虑水平结构。

\ssszh{地平圈}
\S\ref{sss:it Hk} 中给出的迭代 Hecke 叠的定义对任意代数群都说得通，不必要求约化。因此迭代 shtuka 的模叠对 $k$ 上任意代数群也有定义。我们把这一构造用于抛物子群 $P\subset G$，其 Levi 商为 $L$。设有限集 $I$ 有分解 $I=I_{1}\coprod \cdots \coprod I_{r}$，$S$ 是 $k$ 上的测试概形，则 $\Sht^{(I_{1},\cdots, I_{r})}_{P}(S)$ 分类如下的组
$$((x_{i})_{i\in I}, (\cE_{P,j})_{0\le j\le r}, (\a_{j}: \cE_{P, j-1}\dr\cE_{P, j})_{1\le j\le r}, \io)$$
这里与 \S\ref{sss:Sht points} 中 $\Sht^{(I_{1},\cdots, I_{r})}_{G}(S)$ 的描述唯一的差别是：$\cE_{P,j}$ 是 $X$ 上的 $P$-主丛。由该构造对换群的函子性，典范映射 $G\leftarrow{P}\to L$ 诱导映射
\begin{equation*}
\begin{tikzcd}
\Sht_{G}^{(I_{1},\cdots, I_{r})} & \Sht_{P}^{(I_{1},\cdots, I_{r})} \arrow[l,"\frp"'] \arrow[r,"\frq"] & \Sht_{L}^{(I_{1},\cdots, I_{r})}.
\end{tikzcd}
\end{equation*}
映射 $\frp$ 把 $\cE_{P,j}$ 变成诱导的 $G$-主丛 $\cE_{j}=\cE_{P,j}\times^{P}G$；$\frq$ 把 $\cE_{P,j}$ 变成诱导的 $L$-主丛 $\cE_{L,j}=\cE_{P,j}/U_{P}$（其中 $U_{P}$ 是 $P$ 的幂幺根）。

要定义 $\Sht^{(I_{1},\cdots, I_{r})}_{P}$ 的有界版本，只需取 $L$ 的一族容许的支配余权 $\un\mu=(\mu_{i})_{i\in I}$，并令
\begin{equation*}
\Sht_{P}^{(I_{1},\cdots, I_{r}), \le\un\mu}=\frq^{-1}(\Sht_{L}^{(I_{1},\cdots, I_{r}), \le\un\mu} ).
\end{equation*}
\cite[Proposition 5.7]{Var} 证明了：若限制到 $U^{I}$ 的\emph{一般点 $\y^{I}$} 之上，则映射
$$\frp^{\le\un\mu}_{\y^{I}}: \Sht_{P,\y^{I}}^{(I_{1},\cdots, I_{r}), \le\un\mu}\to \Sht_{G,\y^{I}}^{(I_{1},\cdots, I_{r})}$$
（即 $\frp$ 的限制）是有限且非分歧的。$\Sht_{P,\y^{I}}^{(I_{1},\cdots, I_{r}), \le\un\mu}$ 的基本圈的像定义了 $\Sht_{G,\y^{I}}^{(I_{1},\cdots, I_{r})}$ 上的\emph{地平圈}。

粗略地说，地平圈的圈类对应 $\Sht_{G}^{(I_{1},\cdots, I_{r})}$ 上同调中的 Eisenstein 部分。对 \ref{ex:Dr} 中的 $\GL_{2}$ 例子，地平圈的几何在 \cite[\S4]{Dr-FBun} 中被分析，它与 Eisenstein 级数的关系在 \cite[\S3]{Dr-Pet} 中有详细研究。

一般而言，按自守形式空间谱分解的风格来描述 $\Sht_{G}^{(I_{1},\cdots, I_{r}),\le\un\l}$ 上同调的 Eisenstein 部分，仍是一个公开问题。$G=\GL_{2}$ 情形的定性陈述见 \cite[\S7]{YZ1}，一般 $G$ 的有限性结果见 \cite{X}。

\ssszh{特殊圈}

Shimura 簇上的特殊圈由 Shimura 子簇给出，也就是说，它们来自 $G$ 的约化子群。类似地可以讨论 shtuka 模叠上的特殊圈。具体地说，若 $\ph: H\to G$ 是 $k$ 上约化群之间的同态，则由 shtuka 概念对换群的函子性，它诱导模叠之间的映射
$$\th: \Sht^{(I_{1},\cdots, I_{r})}_{H}\to \Sht^{(I_{1},\cdots, I_{r})}_{G}.$$

对有界版本，可以先取 $H$ 的一族支配余权 $\un\mu$，再为 $G$ 取足够大的 $\un\l$，使 $\th$ 把 $\Sht^{(I_{1},\cdots, I_{r}), \le\un\nu}_{H}$ 送到 $\Sht^{(I_{1},\cdots, I_{r}), \le\un\l}_{G}$。Breutmann \cite{Br} 证明了：在允许对 $\Sht^{(I_{1},\cdots, I_{r})}_{H}$ 加抛孔水平结构的一般情形下，$\th$ 是概形的、有限且非分歧的。另见 \cite{Y-cycle} 中的另一个证明。于是对适当的 $\un\l$，$\th$ 的像定义了 $\Sht^{(I_{1},\cdots, I_{r}),\le\un\l}_{G}$ 的 Chow 群中的一个元素。

我们给出几个具有算术意义的特殊圈的例。其中有些并不严格落在上述情形之内。

\begin{exam}[Heegner–Drinfeld 圈]\label{ex:HD}
设 $\nu: X'\to X$ 是非分歧二重覆叠。设 $r$ 为偶数，并把 $I=\{1,2,\cdots, r\}$ 分解为单点集 $I_{i}=\{i\}$。设 $\un\mu=(\mu_{i})_{1\le i\le r}$，其中每个 $\mu_{i}$ 或是 $(1,0,\cdots,0)$，或是 $(0,\cdots, 0,-1)$（$\GL_{n}$ 的极小余权）。令 $\Sht^{\un\mu}_{n,X'}$ 表示 $X'$ 上秩 $n$ 的迭代 shtuka、界由 $\un\mu$ 给出的模叠 $\Sht^{(I_{1},\cdots, I_{r}),\le\un\mu}_{\GL_{n},X'}$。

设 $\un\l=(\l_{i})_{1\le i\le r}$ 是 $\GL_{2n}$ 的由 $\un\mu$ 决定的极小支配余权组：若 $\mu_{i}=(\underbrace{0,\cdots, 0}_{n-1},-1)$，则 $\l_{i}=(\underbrace{0,\cdots,0}_{2n-1}, -1)$；若 $\mu_{i}=(1,\underbrace{0,\cdots,0}_{n-1})$，则 $\l_{i}=(1, \underbrace{0,\cdots,0}_{2n-1})$。类似地，用 $\Sht^{\un\l}_{2n,X}$ 记 $X$ 上秩 $2n$ 的迭代 shtuka、界由 $\un\l$ 给出的模叠 $\Sht^{(I_{1},\cdots, I_{r}),\le\un\l}_{\GL_{2n},X}$。

有映射
\begin{equation*}
\th^{\un\mu}_{\HD(n)}: \Sht^{\un\mu}_{n, X'}\to \Sht^{\un\l}_{2n,X}
\end{equation*}
它把 $X'$ 上秩 $n$ 的 shtuka $((x'_{i}), (\cE_{i}), (\a_{i}), \io)$ 送到 $X$ 上秩 $2n$ 的 shtuka $((\nu(x'_{i})), (\nu_{*}\cE_{i}),(\nu_{*}\a_{i}), \nu_{*}\io)$。把基本圈在 $\th^{\un\mu}_{\HD(n)}$ 下的像称为\emph{Heegner–Drinfeld 圈}。由命题 \ref{p:local model}(2)，$\dim \Sht^{\un\mu}_{n,X'}=rn$，而 $\dim \Sht^{\un\mu}_{2n,X}=2rn$，故有
\begin{equation*}
\dim  \Sht^{\un\mu}_{n,X'}=\frac{1}{2}\dim \Sht^{\un\l}_{2n,X}.
\end{equation*}
因此 $\Sht^{\un\mu}_{n,X'}$ 的基本类在 $\th^{\un\mu}_{\HD(n)}$ 下的像定义了 $\Sht^{\un\l}_{2n,X}$ 上的一个中间维数的圈。

为了看出这个圈与 Heegner 点的关系，我们考虑 $n=1$ 的情形以及上述构造带一条腿的变体。设 $\infty\in |X|$ 是在 $X'$ 中惰性的闭点，于是 $\infty'\in |X'|$ 是 $\infty$ 上方唯一的闭点。考虑 \S\ref{sss:one leg} 中的 Drinfeld 模（相对）曲线 $\Sht^{\Dr}_{2}$。设 $\un\mu=(\mu_{1},\mu_{2})$，其中作为 $\Gm$ 的余权 $\mu_{1}=-1$，$\mu_{2}=1$。令 $\cP_{X'}=\Sht^{(\{1\},\{2\}), \un\mu}_{\GL_{1}, X'}|_{U'\times \{\infty'\}}$；这是 $U'=X'-\infty'$ 上的 $\Pic_{X'}(k)$-主丛。我们构造映射
\begin{equation*}
\th^{1}_{\HD(1)}: \cP_{X'}\to \Sht^{\Dr}_{2}.
\end{equation*}
对 $\cP_{X'}$ 的一个 $S$-点 $\cL_{0}\leftarrow\cL_{0}(-x')\rightarrow\cL_{0}(-x'+y')\isom {}^{\t}\cL_{0}$（其中 $y':S\to \infty'$），考虑 $X'$ 上秩二的丛 $\cE=\nu_{*}\cL_{0}$，它在 $S\times\infty$ 上的纤维有分裂 $\cE|_{S\times\infty}=\cE_{y'}\op \cE_{\s(y')}$（其中 $\s\in \Gal(X'/X)$ 是非平凡元素）。给 $\cE|_{S\times \infty}$ 配上完全旗 $\cE_{y'}\subset \cE_{S\times\infty}$，并定义 $\cE(-\frac{1}{2})\subset \cE$ 为 $\cE_{y'}$ 在限制映射 $\cE\to \cE|_{S\times \infty}$ 下的原像。于是 $\cE(-\frac{1}{2})=\nu_{*}\cL_{0}(-\s(y'))$。因此 ${}^{\t}(\cE(-\frac{1}{2}))=\nu_{*}({}^{\t}\cL_{0})(-y')\cong \nu_{*}\cL_{0}(-x')$（这里用到 $y'\c\Fr_{S}=\s(y')$，因为 $\infty$ 是惰性的）。设 $\phi: {}^{\t}(\cE(-\frac{1}{2}))\to \cE$ 由嵌入 $\cL_{0}(-x')\incl \cL_{0}$ 诱导。这就定义了 $S$-点 $(\xi=\nu(x'), \cE, \{\cE(\frac{i}{2})\}, \phi)\in \Sht^{\Dr}_{2}(S)$。

利用 \S\ref{sss:dict} 中椭圆模与 shtuka 之间的词典，$\cP_{X'}$ 分类 $X$ 上秩二的椭圆模，其自同态由 $A'=\G(U', \cO_{U'})$ 给出；这是带复乘的椭圆曲线的函数域类比。因此 $\th(\cP_{X'})$ 是模曲线上 Heegner 点的函数域类比。带 Iwahori 水平结构的推广在 \cite{YZ2} 中完成。
\end{exam}

\begin{exam}[Gan–Gross–Prasad 圈]\label{ex:GGP}
设 $\nu: X'\to X$ 是非分歧二重覆叠。设 $\Ug_{n}$ 是 $X$ 上关于 $X'$ 的酉群概形。我们不细说非分裂群 $\Ug_{n}$ 本身，只把 $\Bun_{\Ug_{n}}$ 与 $\Sht^{r}_{\Ug_{n}}$ 的模含义写清楚。

设 $S$ 是 $k$ 上的测试概形。$X \times S$ 上关于 $\nu:X'\to X$ 的秩 $n$ \emph{Hermite}（Hermitian，亦称\emph{酉}）丛，是指 $X' \times S$ 上秩 $n$ 的向量丛 $\cF$，并带有同构 $h : \cF \xrightarrow{\sim} \sigma^* \cF^{\vee}$，满足 $\sigma^* h^{\vee} = h$。这里 $\cF^{\vee}=\un\Hom(\cF, \om_{X'})$ 是 $\cF$ 的 Serre 对偶。于是 $\Bun_{\Ug_{n}}$ 是这样的叠：其 $S$-点是 $X \times S$ 上关于 $\nu$ 的秩 $n$ Hermite 丛构成的群胚。

$X'\times S$ 上秩 $n$ Hermite 丛 $(\cF,h)$ 沿 $x': S\to X'$ 的一个初等下修正，是指另一个秩 $n$ Hermite 丛 $(\cF',h')$（同在 $X'\times S$ 上），连同同构 $\a:  \cF|_{X' \times S - \Gamma(x') - \Gamma(\sigma(x'))} \xrightarrow{\sim} \cF'|_{X' \times S - \Gamma(x')-\Gamma(\sigma(x'))}$，它与 Hermite 结构相容，并满足下述性质：存在秩 $n$ 的向量丛 $\cF^{\flat}$（必然唯一），使 $f$ 可以分解为
\begin{equation*}
\begin{tikzcd}
\cF & \cF^{\flat} \arrow[l, hook, "\oll{\a}"] \arrow[r, hook, "\orr{\a}"] & \cF'
\end{tikzcd}
\end{equation*}
使得 $\coker(\oll{\a})$ 在 $\Gamma(x')$ 上是秩 $1$ 局部自由，$\coker(\orr{\a})$ 在 $\Gamma(\sigma(x'))$ 上是秩 $1$ 局部自由。

设 $r$ 为偶数。定义叠 $\Sht_{\Ug_{n}}^{r}$，其 $S$-点由下述数据构成的群胚给出
$$((x'_{i})_{1\le i\le r}, (\cF_{i}, h_{i})_{0\le i\le r}, (\a_{i})_{1\le i\le r}, \io)$$
其中 $x'_{i}: S\to X'$，$(\cF_{i},h_{i})$ 是 $X'\times S$ 上秩 $n$ 的 Hermite 丛，$\a_{i}: \cF_{i-1}\dr\cF_{i}$ 是 $\cF_{i-1}$ 沿 $x'_{i}$ 的初等下修正，$\io$ 是 Hermite 丛的同构 $\cF_{r}\cong {}^{\t}\cF_{0}$。

现在固定 $X'$ 上关于 $\nu$ 的秩一 Hermite 丛 $\cL$；这等价于给一个线丛以及同构 $\Nm_{X'/X}\cL\cong \om_{X}$。有映射
\begin{equation*}
\ka_{\cL}:\Sht^{r}_{\Ug_{n}}\to \Sht^{r}_{\Ug_{n+1}}
\end{equation*}
它把 $((x'_{i}), (\cF_{i}), \cdots)$ 送到 $((x'_{i}), (\cF_{i}\op\cL),\cdots)$。考虑对角映射
\begin{equation*}
\th^{r}_{\GGP(n)}=(\id,\ka_{\cL}): \Sht^{r}_{\Ug_{n}}\to \Sht^{r}_{\Ug_{n}}\times_{X'^{r}}\Sht^{r}_{\Ug_{n+1}}.
\end{equation*}
由 \cite[Lemma 6.9(2)]{FYZ1}，映射 $\Sht^{r}_{\Ug_{n}}\to X'^{r}$ 光滑，相对维数为 $r(n-1)$。因此
$$\dim\Sht^{r}_{\Ug_{n}}=\frac{1}{2}(\dim\Sht^{r}_{\Ug_{n}}\times_{X'^{r}}\Sht^{r}_{\Ug_{n+1}}).$$
$\Sht^{r}_{\Ug_{n}}$ 的基本圈在 $\th^{r}_{\GGP(n)}$ 下的像定义了目标上的一个中间维数的圈。记其圈类为 $Z^{r}_{\GGP(n)}\in \cohog{2rn}{(\Sht^{r}_{\Ug_{n}}\times_{X'^{r}}\Sht^{r}_{\Ug_{n+1}})_{\ov k}}(rn)$。

受算术 Gan–Gross–Prasad 猜想 \cite{GGP} 的启发，我们预期：对 $\Ug_{n}(\AA_{F})\times \Ug_{n+1}(\AA_{F})$ 的尖点自守表示 $\Pi$，$Z^{r}_{\GGP(n)}$ 的 $\Pi$-等型分支的自交数可以良定义（注意它先验地并不是具有 proper 支集的圈），并且与 $\Pi$ 到 $\GL_{n}(\AA_{F'})\times\GL_{n+1}(\AA_{F'})$ 的基变换的 Rankin–Selberg $L$-函数的第 $r$ 个中心导数有关，其中 $F'=k(X')$。

\end{exam}

\begin{exam}[Kudla–Rapoport 圈]\label{ex:KR}
在 \cite{FYZ1} 中我们引入了一类代数圈，它们是酉 Shimura 簇上 Kudla–Rapoport 特殊圈 \cite{KR2} 的函数域类比。

沿用例 \ref{ex:GGP} 的设定。设 $\cE$ 是 $X'$ 上秩 $m$ 的向量丛。定义叠 $\cZ_{\KR(n)}^{r}(\cE)$，其 $S$-点是下述数据构成的群胚：
 \begin{itemize}
 \item 一个 Hermite shtuka $((x'_i)_{1\le i\le r}, (\cF_{i},h_{i})_{0\le i\le r}, (\a_{i})_{1\le i\le r}, \io) \in  \Sht_{\Ug_{n}}^{r}(S)$。
 \item $X'\times S$ 上凝聚层之间的映射 $t_i : \cE \boxtimes \cO_{S} \to \cF_i$（$i=0,1,\cdots, r$），使得同构 $\io : \cF_r \cong {}^{\t}\cF_0$ 把 $t_r$ 与 ${}^{\t}t_0$ 相缠，且对每个 $i = 1, \cdots, r$，映射 $t_{i-1}, t_{i}$ 被修正 $\a_i :  \cF_{i-1} \dr \cF_{i}$ 相缠，即下图交换。
\begin{equation}\label{EFsq}
\begin{tikzcd}
\cE\boxtimes\cO_{S} \arrow[r,equal] \arrow[d,"t_{0}"'] & \cE\boxtimes\cO_{S} \arrow[r,equal] \arrow[d,"t_1"'] & \cdots \arrow[r,equal] \arrow[d] & \cE\boxtimes\cO_{S} \arrow[r,"\sim"] \arrow[d,"t_{r}"'] & {}^{\t}(\cE \boxtimes\cO_{S}) \arrow[d,"{}^{\t}t_{0}"] \\
\cF_0 \arrow[r,dashed,"\a_0"'] & \cF_1 \arrow[r,dashed,"\a_1"'] & \cdots \arrow[r,dashed,"\a_r"'] & \cF_{r} \arrow[r,"\io","\sim"'] & {}^{\t}\cF_0
\end{tikzcd}
\end{equation}
 \end{itemize}

对 $((x'_i), (\cF_i,h_{i}), (\a_i), \io, (t_i)) \in \cZ_{\KR(n)}^r(\cE)(S)$，考虑复合
\begin{equation*}
\cE \boxtimes \cO_S \xr{t_i} \cF_i \xr{h_i} \s^* \cF_i^{\vee} \xr{\s^* t^{\vee}_{i}} \s^* \cE^{\vee} \boxtimes \cO_S.
\end{equation*}
由 \eqref{EFsq} 中的交换方块，上述映射与 $i$ 无关，并且下降为映射 $a: \cE\to \s^{*}\cE^{\vee}$，满足 $\s^{*}a^{\vee}=a$。换言之，$\cZ^{r}_{\KR(n)}(\cE)$ 的每个点都在 $\cE$ 上诱导一个可能退化的 Hermite 形式；据此有开闭子叠的分解
\begin{equation*}
\cZ_{\KR(n)}^{r}(\cE)=\coprod_{a\in \Herm(\cE)} \cZ^{r}_{\KR(n)}(\cE,a).
\end{equation*}
其中 $\Herm(\cE)$ 是 $\cE$ 上（可能退化的）Hermite 形式之集。有自然映射
\begin{equation*}
\th_{\KR(n)}^{r}(\cE,a): \cZ^{r}_{\KR(n)}(\cE,a)\to \Sht^{r}_{\Ug_{n}}
\end{equation*}
它是有限且非分歧的。

$\cZ^{r}_{\KR(n)}(\cE,a)$ 的预期维数是 $r(n-m)$。然而它的实际维数往往更大。在 \cite{FYZ2} 中，我们用导出代数几何给 $\cZ^{r}_{\KR(n)}(\cE,a)$ 配上一个预期维数 $r(n-m)$ 的虚拟基本类 $[\cZ^{r}_{\KR(n)}(\cE,a)]$；它在 $\th^{r}_{\KR(n)}(\cE,a)$ 下的像是酉 Shimura 簇上 Kudla–Rapoport 特殊圈的函数域类比。
\end{exam}

\subsection{相交理论与 $L$-函数的高阶导数}
经典上，自守形式的周期积分常常可以表示成自守 $L$-函数的特殊值。在几何情形中，Shimura 簇或 shtuka 模叠上的特殊圈可以看作周期积分的类比。在数域上，Shimura 簇上特殊圈的高度有时与自守 $L$-函数的导数有关，例如在 Gross–Zagier 公式或算术 Gan–Gross–Prasad 猜想中。在函数域上，不仅有类似的故事，还有进一步的推广：shtuka 可以有多条腿，这种灵活性使人能把特殊圈的相交数与 $L$-函数的高阶导数联系起来。我们简要综述这一方向的若干结果。关于动机与证明中所用技术的更详细讨论，参见 \cite{Y-ICM}。

\ssszh{高阶 Gross–Zagier 公式}

著名的 Gross–Zagier 公式用模形式的自守 $L$-函数的一阶导数，给出模曲线上（投影到 Hecke 本征分量的）Heegner 点高度之间的等式。该公式在数论中有大量应用，尤其在椭圆曲线的算术中。

在 \cite{YZ1} 与 \cite{YZ2} 中，我们给出 Gross–Zagier 公式的函数域类比的高阶导数／多条腿版本。先陈述无水平结构的版本，它在 \cite{YZ1} 中处理。

设 $G=\PGL_{2}$，$I=\{1,\cdots, r\}$ 分解为单点集 $I_{i}=\{i\}$，并设 $\l_{i}$ 是 $G$ 唯一的支配极小余权。把所得的 shtuka 模叠 $\Sht^{(I_{1},\cdots, I_{r}),\le\un\l}_{G}$ 简记为 $\Sht^{r}_{G}$。$\un\l$ 的容许性条件迫使 $r$ 为偶数。

设 $\nu:X'\to X$ 是非分歧二重覆叠。设 $T=(\Res_{X'/X}(\Gm\times X'))/\Gm$ 是 $X$ 上在 $X'$ 处分裂的一维非分裂环面。对 $X$ 上这样的群概形 $T$ 与序列 $\un\mu=(\mu_{1},\cdots, \mu_{r})\in\{\pm1\}^{r}$，有定义 $\Sht_{T}^{\un\mu}:=\Sht_{T}^{(I_{1},\cdots, I_{r}), \un\mu}$。具体地说，它可以等同于 $\Sht^{(I_{1},\cdots, I_{r}), \un\mu}_{\GL_{1}, X'}$（基曲线取 $X'$，见例 \ref{ex:Gm}）在 $\Pic_{X}(k)$ 作用下的商；该作用为拉回到 $X'$ 再与秩一 shtuka 作张量积。如例 \ref{ex:HD} 那样，有自然映射
\begin{equation*}
\th: \Sht^{\un\mu}_{T}\to \Sht^{r}_{G}.
\end{equation*}
考虑 $\th$ 的提升
\begin{equation*}
\th': \Sht^{\un\mu}_{T} \to \Sht'^{r}_{G}:=\Sht^{r}_{G}\times_{X^{r}}X'^{r}.
\end{equation*}
容易看出 $\Sht^{\un\mu}_{T}$ 是 proper 的、维数为 $r$，而由推论 \ref{c:Sht sm}，$\Sht'^{r}_{G}$ 是维数 $2r$ 的光滑概形。因此 Heegner–Drinfeld 圈 $\cZ^{\un\mu}_{T}:=\th'_{*}[\Sht^{\un\mu}_{T}]$ 是光滑 Deligne–Mumford 叠 $\Sht'^{r}_{G}$ 中的中间维数 proper 圈。设 $Z^{\un\mu}_{T}\in \cohoc{2r}{\Sht'^{r}_{G,\kbar},\Qlbar}(r)$ 是 $\cZ^{\un\mu}_{T}$ 的圈类。

现在设 $\pi$ 是 $G(\AA)$ 的处处非分歧尖点自守表示，系数在 $\Qlbar$ 中。在 Hecke 代数作用下，$Z^{\un\mu}_{T}$ 到 $\pi$-等型直和项上的投影是有意义的。把所得的上同调类记作 $Z^{\un\mu}_{T,\pi}\in \cohoc{2r}{\Sht'^{r}_{G,\kbar},\Qlbar}(r)$。

\begin{thm}[\cite{YZ1}]\label{th:YZ1} 圈类 $Z^{\un\mu}_{T,\pi}$ 的自交数为
\begin{equation*}
\j{Z^{\un\mu}_{T,\pi}, Z^{\un\mu}_{T,\pi}}_{\Sht'^{r}_{G}}=\frac{q^{2-2g}}{2(\log q)^{r}}\frac{\sL^{(r)}(\pi_{F'}, 1/2)}{L(\pi,\Ad,1)}
\end{equation*}
其中
\begin{itemize}
\item $\pi_{F'}$ 是 $\pi$ 到 $F'=k(X')$ 的基变换。
\item $\sL(\pi_{F'},s)=q^{4(g-1)(s-1/2)}L(\pi_{F'},s)$ 是 $\pi_{F'}$ 的规范化 $L$-函数，满足 $\sL(\pi_{F'},s)=\sL(\pi_{F'},1-s)$。
\end{itemize}
\end{thm}

在 \cite{YZ2} 中，我们把上述定理推广到允许自守表示 $\pi$ 带 Iwahori 水平结构（即局部表示 $\pi_{v}$ 或是非分歧的，或是 Steinberg 表示的非分歧扭），也允许二重覆叠 $\nu:X'\to X$ 有分歧。这一般性使将来能把结果应用于函数域上的 Birch–Swinnerton-Dyer 猜想，预览见 \cite[\S1.3]{YZ2}。

\ssszh{高阶 Siegel–Weil 公式与高阶 theta 级数}
经典的 Siegel–Weil 公式把辛群（相应地，酉群）上 Siegel–Eisenstein 级数的特殊值与 theta 函数联系起来；后者是数域上二次型（相应地，Hermite 型）表示数的生成级数。在 \cite{Kud} 中，Kudla 开始研究 Siegel–Weil 公式的算术版本：他定义了「算术 theta 函数」——Shimura 曲线整模型上算术圈的生成级数——并发现它与 $\Sp_{4}$ 上 Siegel–Eisenstein 级数的第一个中心导数的关系。在一系列论文 \cite{KR1} 与 \cite{KR2} 中，Kudla 与 Rapoport 通过定义正交群与酉群的 Shimura 簇整模型上特殊圈（Kudla–Rapoport 圈）的生成级数，发展了这一范式。他们猜想这些圈的算术相交数与 Siegel–Eisenstein 级数中心导数的非奇异 Fourier 系数之间的关系。这一猜想最近由 Chao Li 与 Wei Zhang 在 \cite{LZ} 中证明。

在 \cite{FYZ1} 中，我们证明了函数域上酉群的 Siegel–Weil 公式的高阶导数版本。考虑例 \ref{ex:KR} 中构造的特殊圈 $\cZ^{r}_{\KR(n)}(\cE,a)$。我们限定在 $\rk (\cE)=n$ 且 $a: \cE\to\s^{*}\cE^{\vee}$ 在一般点上是同构的情形（这样的 $a$ 称为非奇异的）。此时 $\cZ^{r}_{\KR(n)}(\cE,a)$ 是 $k$ 上的 proper 概形，但它的维数可能不等于预期维数（此时预期维数为 $0$）。我们先定义虚拟基本圈 $[\cZ^{r}_{\KR(n)}(\cE,a)]$，它是 proper 概形 $\cZ^{r}_{\KR(n)}(\cE,a)$ 上的 $0$-圈，从而其次数有定义。我们证明下述高阶 Siegel–Weil 公式

\begin{thm} 设 $r\in \NN$ 为偶数，$\cE$ 是 $X'$ 上秩 $n$ 的向量丛，$a: \cE\to \s^{*}\cE^{\vee}$ 是 $\cE$ 上的非奇异 Hermite 形式。则
\begin{equation*}
\deg [\cZ_{\KR(n)}^r(\cE, a)]=\frac{1}{(\log q)^r}  \left(\dfrac{d}{ds} \right)^r\Big|_{s=0} \left( q^{ds} \wt{E}_{a}(\cE,s,\Phi) \right).
\end{equation*}
这里 $d=-\deg(\cE)+n\deg\omega_X =-\chi(X',\cE)$，而 $\wt{E}_{a}(\cE,s,\Phi)$ 是酉群 $\Ug_{2n}$ 上 Siegel–Eisenstein 级数 $E(g,s,\Phi)$ 的第 $a$ 个 Fourier 系数的某种规范化，在 $\cE$ 处展开（把 $\cE$ 提升为 Siegel–Levi $\GL_{n}(\AA_{F})$ 中的元素）。
\end{thm}

证明与数域情形完全不同（后者基于对 $n$ 的归纳）。证明化归为一个来自 Springer 理论的反常层等式。

在 \cite{FYZ2} 中，我们构造特殊圈的完整生成级数，即对任意秩 $m\le n$ 的 $\cE$ 与 $\cE$ 上任意的 Hermite 形式 $a$，定义虚拟基本类 $[\cZ_{\KR(n)}^r(\cE, a)]$。特别地，$a$ 可以是奇异的甚至为零，此时虚拟基本类会涉及 $\Sht^{r}_{\Ug_{n}}$ 上万有丛的 Chern 类。令 $Z_{\KR(n)}^{r}(\cE,a)$ 为 $[\cZ_{\KR(n)}^r(\cE, a)]$ 在映射 $\th^{r}_{\KR(n)}(\cE,a)$ 下的像，看作 Chow 群 $\Ch_{r(n-m)}(\Sht^{r}_{\Ug_{n}})$ 中的元素。然后把这些代数圈装配成高阶 theta 级数。

更确切地说，固定 $m\le n$，考虑 $X$ 上关于二重覆叠 $X'/X$ 的拟分裂酉群 $\Ug_{2m}$ 及其标准的 Siegel 抛物子群 $P_m$。定义函数
\begin{equation*}
\wt Z^{r}_{m}: \Bun_{P_{m}}(k)\to \Ch_{r(n-m)}(\Sht^{r}_{\Ug_{n}})
\end{equation*}
如下。$\Bun_{P_{m}}(k)$ 中的一个点等同于一对 $(\cG,\cE)$，其中 $\cG$ 是 $X'$ 上秩 $2m$ 的 Hermite 向量丛（这里 Hermite 形式指的是同构 $\cG\isom \s^{*}\cG^{*}$，其中 $\cG^{*}$ 是 $\cG$ 的\emph{线性对偶}，而非 Serre 对偶），$\cE$ 是 $\cG$ 的 Lagrangian 子丛。然后用「Fourier 展开」定义 $\wt Z^{r}_{m}(\cG,\cE)$：
\begin{equation*}
\wt Z^{r}_{m}(\cG,\cE) =\chi(\det\cE)q^{n(\deg \cE-\deg\om_{X}) /2}\sum_{a\in \Herm(\cE)}\psi_{0}(\j{e_{\cG,\cE}, a})Z^{ r}_{\KR(n)}(\cE, a).
\end{equation*}
这里 $\psi_{0}:k\to \Qlbar^{\times}$ 是非平凡特征，$e_{\cG,\cE}\in \Ext^{1}(\s^{*}\cE^{*}, \cE)$ 是正合列 $\cE\incl \cG\surj \s^{*}\cE$ 的扩张类；在 Serre 对偶下它可以与 $a\in \Hom(\cE,\s^{*}\cE^{\vee})$ 配对。

\begin{conj}[模性猜想，{\cite[Conjecture 4.12]{FYZ2}}] 函数 $\wt Z^{r}_{m}$ 下降为函数
\begin{equation*}
Z^{r}_{m}:\Bun_{\Ug_{2m}}(k)\to \Ch_{r(n-m)}(\Sht^{r}_{\Ug_{n}}).
\end{equation*}
\end{conj}
换言之，类 $\wt Z^{r}_{m}(\cG,\cE)\in \Ch_{r(n-m)}(\Sht^{r}_{\Ug_{n}})$ 应当只依赖 Hermite 丛 $\cG$，而不依赖它的 Lagrangian 子丛 $\cE$。当 $r=0$ 时，$\Ch_{0}(\Sht^{0}_{\Ug_{n}})$ 就是 $\Bun_{\Ug_{n}}(k)$ 上的函数空间，经典的 Weil 表示理论表明此时 $\wt Z^{r}_{m}$ 确实下降到 $\Bun_{\Ug_{2m}}(k)$。所得 $\Bun_{\Ug_{2m}}(k)\times \Bun_{\Ug_{n}}(k)$ 上的二元函数就是经典的 theta 函数。

在 \cite{FYZ2} 中我们给出了模性猜想的若干证据；在 \cite{FYZ3} 中我们证明：若限制到 $\Sht^{r}_{\Ug_{n}}\to X'^{r}$ 的通用纤维上，并过渡到圈的上同调类，则猜想成立。

% ch05  参考文献（原文 2029–2276 行）
% 由脚本机械转换：条目标题保留英文，序号与原文一致。

\clearpage
\phantomsection
\addcontentsline{toc}{section}{参考文献}

\begin{thebibliography}{99}

\bibitem{ARH}
E. Arasteh Rad, U. Hartl,
Uniformizing the moduli stacks of global G-shtukas.
Int. Math. Res. Not. IMRN 2021, no. 21, 16121--16192.

\bibitem{And}
G.Anderson, 
$t$-motives.
Duke Math. J. 53 (1986), no. 2, 457--502.

\bibitem{Br}
P. Breutmann, 
Functoriality of Moduli Spaces of Global $G$-Shtukas. 
arXiv:1902.10602.

\bibitem{DL}
P. Deligne, G. Lusztig,
Representations of reductive groups over finite fields. 
Ann. of Math. (2) 103 (1976), no. 1, 103--161.

\bibitem{Dr-Ell}
V. Drinfeld,
Elliptic modules. Mat. Sb. (N.S.) 94(136) (1974), 594--627, 656. Elliptic modules II. Mat. Sb. (N.S.) 102(144) (1977), no. 2, 182--194, 325.

\bibitem{Dr-FBun}
V. Drinfeld,
Moduli varieties of $F$-sheaves. 
Funktsional. Anal. i Prilozhen. 21 (1987), no. 2, 23--41. 
 
\bibitem{Dr-Pet}
V. Drinfeld, 
Proof of the Petersson conjecture for $\GL(2)$ over a global field of characteristic $p$. 
Funktsional. Anal. i Prilozhen. 22 (1988), no. 1, 34--54, 96; translation in Funct. Anal. Appl. 22 (1988), no. 1, 28--43.

\bibitem{Dr-comp}
V.Drinfeld, 
Cohomology of compactified moduli varieties of $F$-sheaves of rank 2.  
Zap. Nauchn. Sem. Leningrad. Otdel. Mat. Inst. Steklov. (LOMI) 162 (1987), Avtomorfn. Funkts. i Teor. Chisel. III, 107--158, 189; translation in J. Soviet Math. 46 (1989), no. 2, 1789--1821.

\bibitem{FYZ1}
T. Feng, Z. Yun, W. Zhang,
Higher Siegel--Weil formula for unitary groups: the non-singular terms.
Invent. Math. 235 (2024), no. 2, 569--668.

\bibitem{FYZ2}
T. Feng, Z. Yun, W. Zhang,
Higher theta series for unitary groups over function fields.
arXiv:2110.07001. To appear in Ann. Sci. E.N.S.

\bibitem{FYZ3}
T. Feng, Z. Yun, W. Zhang,
Modularity of higher theta series I: cohomology of the generic fiber.
arXiv: 2308.10979.

\bibitem{GGP}
W-T. Gan, B.Gross, D.Prasad,
Symplectic local root numbers, central critical $L$ values, and restriction problems in the representation theory of classical groups. 
Sur les conjectures de Gross et Prasad. I. Ast\'erisque No. 346 (2012), 1--109.

\bibitem{GL}
A. Genestier, V. Lafforgue,
Chtoucas restreints pour les groupes r\'eductifs et param\'etrisation de Langlands locale.
arXiv:1709.00978.

\bibitem{Goss}
D. Goss, 
Basic structures of function field arithmetic.
Ergebnisse der Mathematik und ihrer Grenzgebiete (3), 35. Springer-Verlag, Berlin, 1996. xiv+422 pp.

\bibitem{Jan}
U. Jannsen,
Motives, numerical equivalence, and semi-simplicity, 
Invent. Math. 107, 1992, 447--452.

\bibitem{Kot}
R. Kottwitz,
$B(G)$ for all local and global fields.
arXiv:1401.5728.

\bibitem{Kud}
S.Kudla,
Central derivatives of Eisenstein series and height pairings, 
Ann. of Math. (2) 146 (1997), no. 3, 545--646.

\bibitem{KR1}
S. Kudla, M.Rapoport,
Special cycles on unitary Shimura varieties I: Unramified local theory, 
Invent. Math. 184 (2011), no. 3, 629--682.

\bibitem{KR2}
S. Kudla, M. Rapoport, 
Special cycles on unitary Shimura varieties II: Global theory, 
J. Reine Angew. Math. 697 (2014), 91--157.

\bibitem{Laff-comp}
L. Lafforgue,
Une compactification des champs classifiant les chtoucas de Drinfeld. 
J. Amer. Math. Soc. 11 (1998), no. 4, 1001--1036.

\bibitem{LLaff}
L. Lafforgue,
Chtoucas de Drinfeld et correspondance de Langlands. 
Invent. Math. 147 (2002), no. 1, 1--241.

\bibitem{VLaff}
V. Lafforgue,
Chtoucas pour les groupes r\'eductifs et param\'etrisation de Langlands globale.
J. Amer. Math. Soc. 31 (2018), no. 3, 719--891.

 
\bibitem{LR}
R. Langlands, M.Rapoport,
Shimuravariet\"aten und Gerben. 
J. Reine Angew. Math. 378 (1987), 113--220.

\bibitem{LRS}
G. Laumon, M. Rapoport, U. Stuhler,
$\cD$-elliptic sheaves and the Langlands correspondence.
Invent. Math. 113 (1993), no. 2, 217--338.

\bibitem{Laumon}
G. Laumon,
Drinfeld shtukas. 
In {\em Vector bundles on curves-- new directions} (Cetraro, 1995), 50--109, Lecture Notes in Math., 1649, Fond. CIME/CIME Found. Subser., Springer, Berlin, 1997.

\bibitem{LZ}
C. Li, W. Zhang,
Kudla-Rapoport cycles and derivatives of local densities. 
J. Amer. Math. Soc. 35 (2022), no. 3, 705--797. 

\bibitem{Milne}
J. Milne, 
Motives over finite fields.
In {\em Motives} (Seattle, WA, 1991), 401--459, 
Proc. Sympos. Pure Math., 55, Part 1, Amer. Math. Soc., Providence, RI, 1994.

\bibitem{MV}
I. Mirkovic, K.Vilonen,
Geometric Langlands duality and representations of algebraic groups over commutative rings. 
Ann. of Math. (2) 166 (2007), no. 1, 95--143.

\bibitem{Mum}
D. Mumford,
An algebro-geometric construction of commuting operators and of solutions to the Toda lattice equation, Korteweg deVries equation and related nonlinear equation. 
In {\em Proceedings of the International Symposium on Algebraic Geometry} (Kyoto Univ., Kyoto, 1977), pp. 115--153, Kinokuniya Book Store, Tokyo, 1978.

\bibitem{NDT}
T. Ng\^o Dac,
Compactification des champs de chtoucas et th\'eorie g\'eom\'etrique des invariants.
Ast\'erisque No. 313 (2007), 124 pp.

\bibitem{NDT-Sht}
T. Ng\^o Dac,
Introduction to the stacks of shtukas. 
Algebraic cycles, sheaves, shtukas, and moduli, 217--236, 
Trends Math., Birkh\"auser, Basel, 2008.
 
\bibitem{PR}
G. Pappas, M. Rapoport, 
Twisted loop groups and their affine flag varieties. With an appendix by T. Haines and Rapoport. Adv. Math. 219 (2008), no. 1, 118--198. 

\bibitem{Ra}
M. Rapoport,
A guide to the reduction modulo $p$ of Shimura varieties, 
in Automorphic Forms. I, 298, 2005, pp. 271--318.

\bibitem{S-Berkeley}
P. Scholze, J. Weinstein, 
Berkeley lectures on $p$-adic geometry. 
Annals of Mathematics Studies, 207. 
Princeton University Press, Princeton, NJ, 2020. x+250 pp.

\bibitem{SchICM}
P. Scholze, 
$p$-adic geometry. 
In {\em Proceedings of the International Congress of Mathematicians--Rio de Janeiro 2018}. Vol. I. Plenary lectures, 899--933, World Sci. Publ., Hackensack, NJ, 2018.

\bibitem{Var}
Y.Varshavsky, 
Moduli spaces of principal $F$-bundles. 
Selecta Math. (N.S.) 10 (2004), no. 1, 131--166. 

\bibitem{X}
C.Xue,
Finiteness of cohomology groups of stacks of shtukas as modules over Hecke algebras, and applications.
\'Epijournal G\'eom. Alg\'ebrique 4 (2020), Art. 6, 42 pp.

\bibitem{Xue}
C. Xue,
Lectures in the IHES summer school.

\bibitem{YZ1}
Z. Yun, W. Zhang, 
Shtukas and the Taylor expansion of $L$-functions. 
Ann. of Math. (2) 186 (2017), no. 3, 767--911.

\bibitem{YZ2}
Z. Yun, W. Zhang, 
Shtukas and the Taylor expansion of $L$-functions (II). 
Ann. of Math. (2) 189 (2019), no. 2, 393--526. 

\bibitem{Y-ICM}
Z. Yun, 
Hitchin type moduli stacks in automorphic representation theory. 
Proceedings of the International Congress of Mathematicians -- Rio de Janeiro 2018. Vol. II. Invited lectures, 1447--1476, World Sci. Publ.,
Hackensack, NJ, 2018.

\bibitem{Y-cycle}
Z. Yun,
Special cycles for Shtukas are closed.
Pure Appl. Math. Q. 18 (2022), no. 5, 2203--2220.

\bibitem{Zhu-Satake}
X. Zhu, 
An introduction to affine Grassmannians and the geometric Satake equivalence.
In {\em Geometry of Moduli Spaces and Representation Theory}, IAS/Park City Mathematical Series vol 24.

\bibitem{Zhu}
X. Zhu,
Coherent sheaves on the stack of Langlands parameters.
arXiv:2008.02998.

\end{thebibliography}


\end{document}
